% Les engrenages — PCT 3ème — Cours signé OnBuch+
% Programme officiel MINESEC. Compiler avec : tectonic N10-les-engrenages.tex
\newcommand{\DOCMATIERE}{PCT}
\newcommand{\DOCNIVEAU}{3ème}
\newcommand{\DOCMODULE}{Module 4 : Technologie et mécanique}
\newcommand{\DOCLECON}{N10}
\newcommand{\DOCTITRE}{Les engrenages}
% OnBuch+ — préambule « cours signés » ENRICHI (compilé avec tectonic / XeLaTeX)
% Système visuel « Pop » d'OnBuch+ : fond crème (jamais blanc), bordures encre
% épaisses, ombres « dures » décalées, titres Archivo Black en capitales.
% Version MATHÉMATIQUES : graphes (pgfplots/tikz), boîtes pédagogiques
% (définition, propriété, méthode, attention, le savais-tu, exercice résolu,
% exercice, TP, bilan), page de garde et signature OnBuch+ ancrée en bas.
%
% Macros attendues dans main.tex AVANT \input{preamble} :
%   \DOCMATIERE  \DOCNIVEAU  \DOCMODULE  \DOCLECON (numéro)  \DOCTITRE
\documentclass[11pt]{article}

\usepackage{fontspec}
\usepackage[french]{babel}
\usepackage[a4paper,margin=2cm,top=3.4cm,bottom=2.8cm,headheight=1.4cm,footskip=1.3cm]{geometry}
\usepackage{amsmath,amssymb,amsfonts}
\usepackage{mathtools} % \xmapsto, \coloneqq... souvent utilisés par les modèles
\usepackage{cancel} % simplifications barrées (bilans de forces)
% \abs{z} (module, valeur absolue) et \norm{u} (norme), utilisés par les modèles
\providecommand{\abs}{}\let\abs\relax\DeclarePairedDelimiter{\abs}{\lvert}{\rvert}
\providecommand{\norm}{}\let\norm\relax\DeclarePairedDelimiter{\norm}{\lVert}{\rVert}
\usepackage[table]{xcolor} % [table] : fournit aussi \rowcolors (alternance de lignes)
\usepackage{pagecolor}
\usepackage{tikz}
\usetikzlibrary{arrows.meta,positioning,calc,shapes.geometric,shapes.misc,shapes.arrows,decorations.pathmorphing,decorations.pathreplacing,decorations.markings,patterns,shadows,mindmap,trees,fit,backgrounds,matrix,intersections,angles,quotes}
\usepackage{pgfplots}
\usepackage[version=4]{mhchem} % équations nucléaires \ce{^{235}_{92}U}
\usepackage[european,siunitx]{circuitikz} % schémas électriques
\pgfplotsset{compat=1.18}
\usepgfplotslibrary{fillbetween,groupplots} % aires entre courbes (calcul intégral)
\usepackage{enumitem}
\usepackage{fancyhdr}
% [most] charge toutes les bibliothèques tcolorbox (listings, minted, raster,
% documentation...) et gonfle inutilement la mémoire TeX sur un document long
% (~300 boîtes) : on ne charge que ce qui est réellement utilisé.
\usepackage[skins,breakable]{tcolorbox}
\usepackage{graphicx}
\usepackage{colortbl}
\usepackage{booktabs}
\usepackage{tabularx}
\usepackage{multirow}
\usepackage{diagbox} % case d'en-tête diagonale des tableaux à double entrée
% Colonnes extensibles centrée / gauche / droite (C, L, R), souvent utilisées par les modèles
\newcolumntype{C}{>{\centering\arraybackslash}X}
\newcolumntype{L}{>{\raggedright\arraybackslash}X}
\newcolumntype{R}{>{\raggedleft\arraybackslash}X}
\usepackage{array}
\usepackage{multicol}
\usepackage{siunitx}
% parse-numbers=false : accepte aussi \SI{2,5 \times 10^{-3}}{...}
\sisetup{locale=FR,output-decimal-marker={,},group-minimum-digits=4,parse-numbers=false}
\DeclareSIUnit\ohm{\ensuremath{\symup{\Omega}}} % U+2126 absent de la police
\DeclareSIUnit\division{div} % réglages d'oscilloscope (ms/div, V/div)
\DeclareSIUnit\tr{tr} % tours (vitesse de rotation, ex. \SI{1500}{\tr\per\minute})
\DeclareSIUnit\molar{M} % concentration molaire (ex. \SI{3}{\molar}, \SI{1}{\micro\molar})
\DeclareSIUnit\an{an} % année (le modèle confond parfois avec \year, un compteur TeX) — ex. \SI{5}{\kilo\meter\per\an}
\DeclareSIUnit\FCFA{FCFA} % franc CFA (ex. \SI{500}{\FCFA\per\kilo\gram})
\DeclareSIUnit\degreeBrix{\textdegree Brix} % teneur en sucre d'un sirop/jus (réfractométrie)
\DeclareSIUnit\month{mois} % le modèle utilise parfois \month, un compteur TeX, comme unité de durée
\DeclareSIUnit\microliter{\micro\liter} % le modèle écrit parfois \microliter en un seul token
\DeclareSIUnit\million{millions} % dénombrement (ex. \SI{15}{\million\per\mL} spermatozoïdes)
\usepackage{qrcode}
% \degree : commande fréquemment utilisée par les modèles pour un angle ou une
% température en dehors de \SI{}{} (non fournie par les paquets chargés).
\providecommand{\degree}{^{\circ}}
% \qed (fin de démonstration), utilisé par les modèles sans amsthm
\providecommand{\qed}{\hfill$\square$}
% \barre : double barre des valeurs interdites dans un tableau de variations (tkz-tab)
\providecommand{\barre}{\ensuremath{\|}}
% environnement proof (amsthm non chargé) utilisé par les modèles
\ifdefined\proof\else\newenvironment{proof}[1][Démonstration]{\par\noindent\textit{#1.}\ }{\qed\par}\fi
\usepackage{lastpage}
\usepackage{hyperref}
\hypersetup{hidelinks}

% ============================================================
% Palette « Pop » OnBuch+ — mêmes hex que la classe P (plus_ui.dart)
% ============================================================
\definecolor{poporange}{HTML}{F59321}
\definecolor{popdark}{HTML}{E07A0C}
\definecolor{popcream}{HTML}{FFF6E8}
\definecolor{popink}{HTML}{1C1714}
\definecolor{poppurple}{HTML}{7A5AE0}
\definecolor{popgreen}{HTML}{2E9E6A}
\definecolor{poppink}{HTML}{E0457B}
\definecolor{popblue}{HTML}{2F7DE1}
\definecolor{popgold}{HTML}{C98A12}
\definecolor{popmuted}{HTML}{8A7F76}
\definecolor{popline}{HTML}{EADFD2}
\definecolor{popgray}{HTML}{8A7F76}
\colorlet{popgrey}{popgray}
% Alias tolérants (le modèle nomme parfois les couleurs différemment)
\colorlet{popwhite}{white}
\colorlet{popblack}{popink}
\colorlet{popred}{poppink}
\colorlet{popyellow}{popgold}
% Teintes claires (fonds de tableaux/encadrés) — le modèle nomme parfois
% les couleurs avec un suffixe L (ex. popblueL) sans les avoir définies.
\colorlet{poporangeL}{poporange!20}
\colorlet{popdarkL}{popdark!20}
\colorlet{poppurpleL}{poppurple!20}
\colorlet{popgreenL}{popgreen!20}
\colorlet{poppinkL}{poppink!20}
\colorlet{popblueL}{popblue!20}
\colorlet{popgoldL}{popgold!20}
\colorlet{popredL}{popred!20}
\colorlet{popgrayL}{popgray!20}
% Teintes claires pour fonds de tableaux / zones
\colorlet{popblueL}{popblue!10!white}
\colorlet{poporangeL}{poporange!14!white}
\colorlet{popgreenL}{popgreen!12!white}
\colorlet{poppurpleL}{poppurple!12!white}
\colorlet{poppinkL}{poppink!12!white}
\colorlet{popgoldL}{popgold!14!white}

\pagecolor{popcream}

% ============================================================
% Typographie
% ============================================================
\defaultfontfeatures{Ligatures=TeX}
\setmainfont{PlusJakartaSans-Medium.ttf}[Path=../../../fonts/,BoldFont=PlusJakartaSans-Bold.ttf]
% Maths en sans-serif (Fira Math) avec chiffres et lettres droites de Plus
% Jakarta, pour que les formules se fondent dans le texte.
\usepackage[math-style=ISO,bold-style=ISO]{unicode-math}
\setmathfont{FiraMath-Regular.otf}[Path=../../../fonts/]
\setmathfont{PlusJakartaSans-Medium.ttf}[Path=../../../fonts/,range={up/num,up/{Latin,latin},\mathup}]
\usepackage{icomma} % virgule décimale sans espace en mode math
% Caractères mathématiques Unicode absents des polices de texte (Plus Jakarta,
% Archivo Black) : rendus par leur équivalent mathématique, en texte comme en
% formule — sinon ils s'affichent en carré vide (ex. « Arithmétique dans ℤ »).
\usepackage{newunicodechar}
\newunicodechar{ℝ}{\ensuremath{\mathbb{R}}}
\newunicodechar{ℤ}{\ensuremath{\mathbb{Z}}}
\newunicodechar{ℂ}{\ensuremath{\mathbb{C}}}
\newunicodechar{ℕ}{\ensuremath{\mathbb{N}}}
\newunicodechar{ℚ}{\ensuremath{\mathbb{Q}}}
\newunicodechar{𝔻}{\ensuremath{\mathbb{D}}}
\newunicodechar{α}{\ensuremath{\alpha}}
\newunicodechar{β}{\ensuremath{\beta}}
\newunicodechar{γ}{\ensuremath{\gamma}}
\newunicodechar{δ}{\ensuremath{\delta}}
\newunicodechar{ε}{\ensuremath{\varepsilon}}
\newunicodechar{θ}{\ensuremath{\theta}}
\newunicodechar{λ}{\ensuremath{\lambda}}
\newunicodechar{μ}{\ensuremath{\mu}}
\newunicodechar{σ}{\ensuremath{\sigma}}
\newunicodechar{τ}{\ensuremath{\tau}}
\newunicodechar{φ}{\ensuremath{\varphi}}
\newunicodechar{ψ}{\ensuremath{\psi}}
\newunicodechar{ω}{\ensuremath{\omega}}
\newunicodechar{Σ}{\ensuremath{\Sigma}}
% majuscules produites par \MakeUppercase dans les titres (α-aminés -> Α)
\newunicodechar{Α}{\ensuremath{\alpha}}
\newunicodechar{Β}{\ensuremath{\beta}}
\newunicodechar{Γ}{\ensuremath{\Gamma}}
\newunicodechar{Δ}{\ensuremath{\Delta}}
\newunicodechar{Ε}{\ensuremath{\varepsilon}}
\newunicodechar{Θ}{\ensuremath{\Theta}}
\newunicodechar{Λ}{\ensuremath{\Lambda}}
\newunicodechar{Μ}{\ensuremath{\mu}}
\newunicodechar{Τ}{\ensuremath{\tau}}
\newunicodechar{Φ}{\ensuremath{\Phi}}
\newunicodechar{Ψ}{\ensuremath{\Psi}}
\newunicodechar{Ω}{\ensuremath{\Omega}}
\newunicodechar{↦}{\ensuremath{\mapsto}}
\newunicodechar{∈}{\ensuremath{\in}}
\newunicodechar{∉}{\ensuremath{\notin}}
\newunicodechar{∧}{\ensuremath{\wedge}}
\newunicodechar{⇔}{\ensuremath{\Leftrightarrow}}
\newunicodechar{⟺}{\ensuremath{\Longleftrightarrow}}
\newunicodechar{⇒}{\ensuremath{\Rightarrow}}
\newunicodechar{⟹}{\ensuremath{\Longrightarrow}}
\newunicodechar{∘}{\ensuremath{\circ}}
\newunicodechar{∪}{\ensuremath{\cup}}
\newunicodechar{∩}{\ensuremath{\cap}}
\newunicodechar{≡}{\ensuremath{\equiv}}
\newunicodechar{⊂}{\ensuremath{\subset}}
\newunicodechar{⊥}{\ensuremath{\perp}}
\newunicodechar{∃}{\ensuremath{\exists}}
\newunicodechar{∀}{\ensuremath{\forall}}
\newunicodechar{∛}{\ensuremath{\sqrt[3]{\,}}}
\newunicodechar{∜}{\ensuremath{\sqrt[4]{\,}}}
\newunicodechar{⃗}{\ensuremath{{}^{\to}}}
\newunicodechar{ⁿ}{\ifmmode^{n}\else\textsuperscript{\ensuremath{n}}\fi}
\newunicodechar{ˣ}{\ifmmode^{x}\else\textsuperscript{\ensuremath{x}}\fi}
\newunicodechar{ᵇ}{\ifmmode^{b}\else\textsuperscript{\ensuremath{b}}\fi}
\newunicodechar{ʳ}{\ifmmode^{r}\else\textsuperscript{\ensuremath{r}}\fi}
\newunicodechar{ᵘ}{\ifmmode^{u}\else\textsuperscript{\ensuremath{u}}\fi}
\newunicodechar{ʸ}{\ifmmode^{y}\else\textsuperscript{\ensuremath{y}}\fi}
\newunicodechar{ᵃ}{\ifmmode^{a}\else\textsuperscript{\ensuremath{a}}\fi}
\newunicodechar{ᵉ}{\ifmmode^{e}\else\textsuperscript{\ensuremath{e}}\fi}
\newunicodechar{ᵏ}{\ifmmode^{k}\else\textsuperscript{\ensuremath{k}}\fi}
\newunicodechar{ᵖ}{\ifmmode^{p}\else\textsuperscript{\ensuremath{p}}\fi}
\newunicodechar{ᴺ}{\ifmmode^{N}\else\textsuperscript{\ensuremath{N}}\fi}
\newunicodechar{ᶜ}{\ifmmode^{c}\else\textsuperscript{\ensuremath{c}}\fi}
\newunicodechar{ʰ}{\ifmmode^{h}\else\textsuperscript{\ensuremath{h}}\fi}
\newunicodechar{ᵠ}{\ifmmode^{q}\else\textsuperscript{\ensuremath{q}}\fi}
\newunicodechar{ₙ}{\ifmmode_{n}\else\textsubscript{\ensuremath{n}}\fi}
\newunicodechar{ₐ}{\ifmmode_{a}\else\textsubscript{\ensuremath{a}}\fi}
\newunicodechar{ᵢ}{\ifmmode_{i}\else\textsubscript{\ensuremath{i}}\fi}
\newunicodechar{ᵣ}{\ifmmode_{r}\else\textsubscript{\ensuremath{r}}\fi}
\newunicodechar{ₖ}{\ifmmode_{k}\else\textsubscript{\ensuremath{k}}\fi}
\newunicodechar{ᵦ}{\ifmmode_{\beta}\else\textsubscript{\ensuremath{\beta}}\fi}
\newunicodechar{ₓ}{\ifmmode_{x}\else\textsubscript{\ensuremath{x}}\fi}
\newfontfamily\popdisplay{ArchivoBlack-Regular.ttf}[Path=../../../fonts/]
\newfontfamily\poplogo{SpaceGrotesk-Bold.ttf}[Path=../../../fonts/]
\newfontfamily\popsemi{PlusJakartaSans-SemiBold.ttf}[Path=../../../fonts/]
\newfontfamily\popxbold{PlusJakartaSans-ExtraBold.ttf}[Path=../../../fonts/]
\newfontfamily\popxboldls{PlusJakartaSans-ExtraBold.ttf}[Path=../../../fonts/,LetterSpace=3.0]

\setlist[enumerate]{leftmargin=1.7em,itemsep=5pt,topsep=4pt}
\setlist[itemize]{leftmargin=1.7em,itemsep=4pt,topsep=4pt,label={\color{poporange}\textbullet}}
\setlength{\parindent}{0pt}
\setlength{\parskip}{5pt}
\renewcommand{\arraystretch}{1.25}

% pgfplots : style Pop par défaut
\pgfplotsset{
  every axis/.append style={axis line style={popink,line width=1pt},tick style={popink},
    grid style={popline,line width=0.5pt},label style={font=\small},tick label style={font=\footnotesize}},
  popcurve/.style={poporange,line width=1.8pt,no markers,smooth},
}
% Même style disponible dans un \draw TikZ simple (hors pgfplots)
\tikzset{popcurve/.style={poporange,line width=1.8pt}}
\tikzset{popgrid/.style={popline,very thin}}
\colorlet{popcurve}{poporange} % le nom du style est parfois utilisé comme couleur

% ============================================================
% Pill / squiggle
% ============================================================
\newcommand{\poppill}[3]{%
  \tikz[baseline=-0.55ex]\node[draw=popink,line width=1.2pt,rounded corners=2.6pt,%
    fill=#2,inner xsep=6.5pt,inner ysep=3pt]{%
    \popxboldls\fontsize{7}{7}\selectfont\color{#3}\MakeUppercase{#1}};%
}
\newcommand{\popsquiggle}[1]{%
  \tikz[baseline=-0.4ex]{
    \draw[#1,line width=2.6pt,line cap=round]
      plot [smooth] coordinates {(0,0.09) (0.34,0.01) (0.68,0.21) (1.02,0.01) (1.36,0.10)};
  }%
}
% Mot-clé surligné (marqueur orange sous le texte)
\newcommand{\cle}[1]{{\popxbold\color{popdark}#1}}

% ============================================================
% En-tête / pied
% ============================================================
\pagestyle{fancy}
\fancyhf{}
\renewcommand{\headrulewidth}{0pt}
\renewcommand{\footrulewidth}{0pt}
\lhead{\raisebox{-0.15ex}{\poplogo\fontsize{13}{13}\selectfont\color{popink}OnBuch\color{poporange}+}}
\rhead{\poppill{\DOCMATIERE\ \textbullet\ \DOCNIVEAU\ \textbullet\ Leçon \DOCLECON}{white}{popink}}
\lfoot{\popxbold\fontsize{7.6}{7.6}\selectfont\color{popmuted}\MakeUppercase{Cours signé OnBuch+}\ \textbullet\ {\popsemi\DOCTITRE}}
\rfoot{\tikz[baseline=-0.6ex]\node[rounded corners=8.5pt,fill=popink,minimum height=17pt,minimum width=17pt,inner xsep=5pt,inner sep=0]{\popxbold\fontsize{8}{8}\selectfont\color{popcream}\thepage\,/\,\pageref*{LastPage}};}

% ============================================================
% Titres
% ============================================================
\newcommand{\chaptitre}[2]{%
  \begin{tcolorbox}[enhanced,colback=poporange,colframe=popink,boxrule=2.6pt,arc=11pt,
      drop shadow=popink,left=15pt,right=15pt,top=12pt,bottom=11pt,before skip=3pt,after skip=18pt]
    \poppill{#2}{popink}{popcream}\par\vspace{9pt}
    {\raggedright\hyphenpenalty=10000\exhyphenpenalty=10000\popdisplay\fontsize{22}{24}\selectfont\color{popcream}\MakeUppercase{#1}\par}
  \end{tcolorbox}
}
% Section numérotée : pastille orange avec le numéro + titre Archivo
\newcounter{popsec}
\newcommand{\coursec}[1]{%
  \stepcounter{popsec}\setcounter{popsub}{0}%
  \par\vspace{16pt}\noindent
  \begin{minipage}{\linewidth}
  \tikz[baseline=-0.7ex]\node[draw=popink,line width=1.6pt,fill=poporange,rounded corners=4pt,minimum size=22pt,inner sep=0]{\popdisplay\fontsize{12}{12}\selectfont\color{popcream}\thepopsec};%
  \hspace{8pt}{\popdisplay\fontsize{15}{17}\selectfont\color{popink}\MakeUppercase{#1}}\par
  \vspace{4pt}\noindent{\color{poporange}\rule{36pt}{3.6pt}}
  \end{minipage}\par\nopagebreak\vspace{8pt}
}
\newcounter{popsub}
\newcommand{\courssub}[1]{%
  \stepcounter{popsub}%
  \par\vspace{9pt}\noindent{\popxbold\fontsize{11.5}{13}\selectfont\color{popink}{\color{poporange}\thepopsec.\thepopsub}\hspace{6pt}#1}\par\nopagebreak\vspace{4pt}
}

% ============================================================
% Boîtes pédagogiques — même recette : fond blanc, bordure épaisse colorée,
% ombre dure encre, badge plein. Titre optionnel [#1].
% ============================================================
\tcbset{popbase/.style={breakable,enhanced,colback=white,boxrule=2.3pt,arc=9pt,
  shadow={3pt}{-3pt}{0pt}{popink},left=14pt,right=14pt,top=11pt,bottom=11pt,before skip=11pt,after skip=15pt,
  fontupper=\popsemi\fontsize{10.6}{14.5}\selectfont\color{popink},
  fontlower=\popsemi\fontsize{10.6}{14.5}\selectfont\color{popink}}}
% Coche dessinée (le glyphe ✓ n'existe pas dans les polices du document)
\newcommand{\popcheck}[1]{\tikz[baseline=-0.5ex]\draw[#1,line width=1.6pt,line cap=round,line join=round](-0.13,0.0)--(-0.04,-0.09)--(0.13,0.1);}
\providecommand{\coche}{\popcheck{popgreen}}
\providecommand{\resultat}[1]{\par\smallskip\noindent\fcolorbox{popgreen}{popgreenL}{\parbox{\dimexpr\linewidth-2\fboxsep-2\fboxrule\relax}{\textbf{Résultat.} #1}}\par\smallskip}
\newcommand{\popboxhead}[3]{\poppill{#1}{#2}{white}\ifx\relax#3\relax\else\hspace{6pt}{\popxbold\fontsize{10.6}{12}\selectfont\color{popink}#3}\fi\par\vspace{7pt}}

\newtcolorbox{aretenir}[1][]{popbase,colframe=popgreen,before upper={\popboxhead{À retenir}{popgreen}{#1}}}
\newtcolorbox{exemplebox}[1][]{popbase,colframe=poporange,before upper={\popboxhead{Exemple}{poporange}{#1}}}
\newtcolorbox{definition}[1][]{popbase,colframe=popblue,before upper={\popboxhead{Définition}{popblue}{#1}}}
\newtcolorbox{propriete}[1][]{popbase,colframe=poppurple,before upper={\popboxhead{Propriété}{poppurple}{#1}}}
\newtcolorbox{methode}[1][]{popbase,colframe=popgold,before upper={\popboxhead{Méthode}{popgold}{#1}}}
\newtcolorbox{attention}[1][]{popbase,colframe=poppink,before upper={\popboxhead{Attention}{poppink}{#1}}}
\newtcolorbox{experience}[1][]{popbase,colframe=popblue,colback=popblueL,before upper={\popboxhead{Expérience}{popblue}{#1}}}
% « Le savais-tu ? » : carte violette pleine (contexte, culture, Cameroun)
\newtcolorbox{savaistu}[1][]{popbase,colback=poppurple,colframe=popink,
  fontupper=\popsemi\fontsize{10.4}{14.2}\selectfont\color{white},
  before upper={\poppill{Le savais-tu\,?}{white}{poppurple}\ifx\relax#1\relax\else\hspace{6pt}{\popxbold\color{white}#1}\fi\par\vspace{7pt}}}
% Exercice résolu : énoncé en haut, \tcblower puis la solution sur fond crème
\newcounter{popexo}\newcounter{popexr}
\newtcolorbox{exoresolu}[1][]{popbase,colframe=poporange,colbacklower=popcream,
  segmentation style={popink,line width=1.2pt,dashed},
  before upper={\stepcounter{popexr}\popboxhead{Exercice résolu \thepopexr}{poporange}{#1}},
  before lower={\poppill{Solution}{popink}{popcream}\par\vspace{6pt}}}
% Exercice (non résolu en place) — la correction est dans la partie Corrigés
\newtcolorbox{exercice}[1][]{popbase,colframe=popink,
  before upper={\stepcounter{popexo}\popboxhead{Exercice \thepopexo}{popink}{#1}}}
% Corrigé
\newtcolorbox{corrige}[1][]{popbase,colframe=popgreen,colback=popgreenL,
  before upper={\popboxhead{Corrigé}{popgreen}{#1}}}
% Objectifs / prérequis / bilan
\newtcolorbox{objectifs}{popbase,colframe=popink,colback=poporangeL,
  before upper={\popboxhead{Objectifs de la leçon}{popink}{}}}
\newtcolorbox{prerequis}{popbase,colframe=popink,colback=popblueL,
  before upper={\popboxhead{Prérequis}{popblue}{}}}
\newtcolorbox{bilan}{popbase,colframe=popink,colback=white,boxrule=2.8pt,
  before upper={\popboxhead{Fiche bilan}{poporange}{L'essentiel à connaître pour l'examen}}}
% Figure encadrée avec légende
\newtcolorbox{popfigure}[1][]{enhanced,breakable=false,colback=white,colframe=popline,boxrule=1.4pt,arc=8pt,
  left=10pt,right=10pt,top=10pt,bottom=8pt,before skip=10pt,after skip=12pt,halign=center,
  lower separated=false,halign lower=center,
  fontlower=\popsemi\fontsize{9}{11.5}\selectfont\color{popmuted},
  #1}
\newcommand{\legende}[1]{\par\vspace{6pt}{\popsemi\fontsize{9}{11.5}\selectfont\color{popmuted}#1}}

% ============================================================
% Signature OnBuch+
% ============================================================
\newcommand{\poplogoinline}[1]{{\poplogo\fontsize{#1}{#1}\selectfont\color{popink}OnBuch\color{poporange}+}}
\newcommand{\signatureonbuch}{%
  \begin{tcolorbox}[enhanced,colback=popink,colframe=popink,boxrule=2.2pt,arc=10pt,
      drop shadow=poporange,left=14pt,right=14pt,top=10pt,bottom=10pt,before skip=0pt,after skip=0pt]
    \tikz[baseline=-0.7ex]\node[circle,fill=popgreen,draw=popcream,line width=1.3pt,minimum size=22pt,inner sep=0]{\popcheck{white}};%
    \hspace{9pt}\begin{minipage}[c]{0.62\linewidth}
      {\popxbold\fontsize{12}{13}\selectfont\color{popcream}Signé\ \textbullet\ {\poplogo OnBuch\color{poporange}+}}\par\vspace{2pt}
      {\raggedright\popsemi\fontsize{8}{10.5}\selectfont\color{popline}\MakeUppercase{Équipe pédagogique OnBuch+}\\\MakeUppercase{Conforme au programme officiel MINESEC}\par}
    \end{minipage}\hfill
    {\poplogo\fontsize{22}{22}\selectfont\color{popcream}OnBuch\color{poporange}+}
  \end{tcolorbox}%
}

% ============================================================
% Page de garde
% #1 titre · #2 matière · #3 niveau · #4 module · #5 n° leçon · #6 durée ·
% #7 description / objectifs (texte ou itemize)
% ============================================================
\newcommand{\pagegarde}[7]{%
  \thispagestyle{empty}
  \newgeometry{a4paper,margin=1.7cm,top=1.6cm,bottom=1.4cm}%
  \begin{tikzpicture}[remember picture,overlay]
    \fill[poporange!18] ([xshift=-3.2cm,yshift=-3.6cm]current page.north east) circle (4.6cm);
    \fill[poppurple!14] ([xshift=2.4cm,yshift=7.6cm]current page.south west) circle (3.8cm);
  \end{tikzpicture}%
  {\poplogo\fontsize{30}{30}\selectfont\color{popink}OnBuch\color{poporange}+}\hfill
  \raisebox{0.6ex}{\poppill{Cours signé \textbullet\ Premium}{popink}{popcream}}\par
  \vspace{1pt}\popsquiggle{poppurple}\par
  \vspace{8pt}
  \begin{tcolorbox}[enhanced,colback=poporange,colframe=popink,boxrule=2.8pt,arc=13pt,
      drop shadow=popink,left=18pt,right=18pt,top=12pt,bottom=13pt,after skip=10pt]
    \poppill{#2 \textbullet\ #3}{popink}{popcream}\hspace{5pt}\poppill{#4}{white}{popink}\par\vspace{8pt}
    {\popdisplay\fontsize{14}{14}\selectfont\color{popink}LEÇON #5}\par\vspace{4pt}
    {\raggedright\hyphenpenalty=10000\exhyphenpenalty=10000\popdisplay\fontsize{25}{27}\selectfont\color{popcream}\MakeUppercase{#1}\par}
    \vspace{8pt}
    {\popxbold\fontsize{9.5}{9.5}\selectfont\color{popink}\MakeUppercase{Durée conseillée : #6}}
  \end{tcolorbox}
  \begin{tcolorbox}[enhanced,colback=white,colframe=popink,boxrule=2.2pt,arc=10pt,
      drop shadow=popink,left=16pt,right=16pt,top=10pt,bottom=10pt,after skip=8pt]
    \poppill{Dans cette leçon}{poppurple}{white}\par\vspace{5pt}
    \popsemi\fontsize{9.3}{12.6}\selectfont\color{popink}#7
  \end{tcolorbox}
  \noindent\begin{minipage}[t]{0.487\linewidth}
    \begin{tcolorbox}[enhanced,colback=popgreen,colframe=popink,boxrule=2.2pt,arc=10pt,
        drop shadow=popink,left=11pt,right=11pt,top=10pt,bottom=10pt,height=3.1cm,valign=center]
      \tcbox[colback=white,colframe=popink,boxrule=1.3pt,arc=5pt,boxsep=0pt,nobeforeafter,
        left=3pt,right=3pt,top=3pt,bottom=3pt,valign=center]{\qrcode[height=1.9cm]{https://play.google.com/store/apps/details?id=com.luvvix.onbuch&hl=fr}}%
      \hspace{8pt}\begin{minipage}[c]{0.5\linewidth}
        {\popdisplay\fontsize{11}{12.5}\selectfont\color{white}\MakeUppercase{Télécharge OnBuch}}\par\vspace{4pt}
        {\raggedright\popsemi\fontsize{8.4}{11}\selectfont\color{white}Gratuit \textbullet\ Google Play\\Tuteur IA, annales, concours, résultats\par}
      \end{minipage}
    \end{tcolorbox}
  \end{minipage}\hfill
  \begin{minipage}[t]{0.487\linewidth}
    \begin{tcolorbox}[enhanced,colback=poppurple,colframe=popink,boxrule=2.2pt,arc=10pt,
        drop shadow=popink,left=11pt,right=11pt,top=10pt,bottom=10pt,height=3.1cm,valign=center]
      \tikz[baseline=-0.7ex]\node[circle,fill=white,draw=popink,line width=1.3pt,minimum size=1.6cm,inner sep=0]{\popdisplay\fontsize{24}{24}\selectfont\color{poppurple}+};%
      \hspace{8pt}\begin{minipage}[c]{0.6\linewidth}
        {\popdisplay\fontsize{11}{12.5}\selectfont\color{white}\MakeUppercase{Passe à OnBuch+}}\par\vspace{4pt}
        {\raggedright\popsemi\fontsize{8.4}{11}\selectfont\color{white}Tous les cours signés, Tuteur IA illimité, fiches PDF — onglet OnBuch+ de l'app\par}
      \end{minipage}
    \end{tcolorbox}
  \end{minipage}
  \vspace{8pt}
  \popcontenu
  \vfill
  \signatureonbuch
  \restoregeometry
  \newpage
}

% Rangée « Ce cours contient » (page de garde)
\newcommand{\poptile}[3]{% #1 couleur #2 gros texte #3 légende
  \begin{tcolorbox}[enhanced,colback=#1,colframe=popink,boxrule=2pt,arc=9pt,shadow={2.5pt}{-2.5pt}{0pt}{popink},
      left=6pt,right=6pt,top=9pt,bottom=9pt,halign=center,valign=center,height=1.75cm,width=\linewidth,nobeforeafter]
    {\popdisplay\fontsize{12.5}{13}\selectfont\color{white}\MakeUppercase{#2}}\par\vspace{3pt}
    {\centering\popsemi\fontsize{7.8}{9.5}\selectfont\color{white}#3\par}
  \end{tcolorbox}}
\newcommand{\popcontenu}{%
  \par\noindent{\popxboldls\fontsize{7.5}{7.5}\selectfont\color{popmuted}CE COURS SIGNÉ CONTIENT}\par\vspace{7pt}
  \noindent\begin{minipage}[t]{0.235\linewidth}\poptile{popblue}{Cours}{complet, illustré, pas à pas}\end{minipage}\hfill
  \begin{minipage}[t]{0.235\linewidth}\poptile{poporange}{Méthodes}{et exercices résolus}\end{minipage}\hfill
  \begin{minipage}[t]{0.235\linewidth}\poptile{poppink}{Exercices}{type examen + corrigés}\end{minipage}\hfill
  \begin{minipage}[t]{0.235\linewidth}\poptile{popgreen}{Fiche bilan}{l'essentiel à retenir}\end{minipage}\par}

% Sommaire
\newtcolorbox{sommaire}{enhanced,colback=white,colframe=popink,boxrule=2.2pt,arc=10pt,
  drop shadow=popink,left=18pt,right=18pt,top=13pt,bottom=13pt,before skip=6pt,after skip=20pt,
  before upper={\poppill{Sommaire}{poppurple}{white}\par\vspace{9pt}\popsemi\fontsize{10.6}{14.5}\selectfont\color{popink}}}

% Signature de fin de cours — poussée tout en bas de la dernière page
\newcommand{\findecours}{%
  \par\vfill\nopagebreak
  \begin{center}\popsquiggle{poporange}\end{center}\vspace{4pt}
  \signatureonbuch
}
\AtBeginDocument{\def\llbracket{\mathopen{[\mkern-3.5mu[}}\def\rrbracket{\mathclose{]\mkern-3.5mu]}}} % intervalles d'entiers (glyphe ⟦ absent de la police)
% Glyphes absents de Fira Math : redéfinis à partir de glyphes disponibles
\AtBeginDocument{%
  \def\setminus{\mathbin{\backslash}}\let\smallsetminus\setminus
  \def\vdots{\mathord{\vbox{\baselineskip3.2pt\lineskiplimit0pt\kern4pt\hbox{.}\hbox{.}\hbox{.}}}}%
  \def\ddots{\mathinner{\mkern1mu\raise6.5pt\hbox{.}\mkern2mu\raise3.6pt\hbox{.}\mkern2mu\raise0.7pt\hbox{.}\mkern1mu}}%
  \def\triangle{\mathord{\scalebox{0.9}{$\symup{\Delta}$}}}%
  \def\nabla{\mathord{\raisebox{\depth}{\scalebox{1}[-1]{$\symup{\Delta}$}}}}%
  \def\checkmark{\popcheck{popdark}}%
  \def\star{\mathbin{\ast}}\def\bigstar{\mathord{\ast}}%
  \def\diamond{\mathbin{\scalebox{0.6}{$\Diamond$}}}%
  \def\bigcup{\mathop{\mathchoice{\raisebox{-0.3em}{\scalebox{1.7}{$\cup$}}}{\scalebox{1.25}{$\cup$}}{\cup}{\cup}}\displaylimits}%
  \def\bigcap{\mathop{\mathchoice{\raisebox{-0.3em}{\scalebox{1.7}{$\cap$}}}{\scalebox{1.25}{$\cap$}}{\cap}{\cap}}\displaylimits}%
  \def\lesseqqgtr{\mathrel{\lessgtr}}\def\bigoplus{\mathop{\oplus}\displaylimits}%
}
\providecommand{\ssi}{si et seulement si} % abréviation fréquente
\AtBeginDocument{\providecommand{\wideparen}{\overparen}} % arc de cercle

%% FIGFIX-BEGIN v5 — correctifs globaux des figures (docs/onbuch-generation/tools/figfix)
\usepackage{adjustbox}\usepackage{etoolbox}\tcbuselibrary{raster}
% toute figure TikZ/pgfplots plus large que la ligne est réduite à la largeur disponible
\BeforeBeginEnvironment{tikzpicture}{\begin{adjustbox}{max width=\linewidth}}
\AfterEndEnvironment{tikzpicture}{\end{adjustbox}}
% légendes pgfplots lisibles par-dessus les courbes
\pgfplotsset{every axis/.append style={legend style={fill=white,fill opacity=0.92,text opacity=1,draw=black!15,inner sep=3pt}}}
% étiquettes d'arêtes (syntaxe edge["..."]) avec fond blanc
\tikzset{every edge quotes/.append style={fill=white,inner sep=1.2pt}}
%% FIGFIX-END
\begin{document}

\pagegarde{Les engrenages}{PCT}{3ème}{Module 4 : Technologie et mécanique}{N10}{1 séance (1 h chacune)}{Cette leçon présente les engrenages, organes de transmission de mouvement par roues dentées. L'élève apprend à définir un engrenage, à en citer les avantages et les inconvénients, et à calculer un rapport de transmission dans des situations concrètes de la vie courante.
\vspace{4pt}
\begin{multicols}{2}\small
\begin{itemize}
\item À la fin de cette leçon, je sais définir un engrenage et une roue dentée.
\item À la fin de cette leçon, je sais identifier les éléments d'un engrenage (roue menante, roue menée, dents, module) sur un schéma ou un objet réel.
\item À la fin de cette leçon, je sais décrire le sens de rotation des roues d'un engrenage extérieur.
\item À la fin de cette leçon, je sais citer les avantages et les inconvénients de la transmission par engrenages.
\item À la fin de cette leçon, je sais définir et calculer le rapport de transmission d'un engrenage.
\item À la fin de cette leçon, je sais déterminer si un engrenage est réducteur ou multiplicateur de vitesse.
\end{itemize}\end{multicols}}

\begin{sommaire}
\begin{multicols}{2}
\begin{enumerate}
\item Découverte expérimentale des engrenages
\item Définitions et vocabulaire
\item Types d'engrenages et sens de rotation
\item Avantages et inconvénients des engrenages
\item Rapport de transmission
\item Sécurité et entretien des machines à engrenages
\item Activité d'intégration
\item Méthodes et exercices résolus
\item Exercices
\item Corrigés des exercices
\item Fiche bilan
\end{enumerate}
\end{multicols}
\end{sommaire}

\begin{objectifs}
\begin{itemize}
\item À la fin de cette leçon, je sais définir un engrenage et une roue dentée.
\item À la fin de cette leçon, je sais identifier les éléments d'un engrenage (roue menante, roue menée, dents, module) sur un schéma ou un objet réel.
\item À la fin de cette leçon, je sais décrire le sens de rotation des roues d'un engrenage extérieur.
\item À la fin de cette leçon, je sais citer les avantages et les inconvénients de la transmission par engrenages.
\item À la fin de cette leçon, je sais définir et calculer le rapport de transmission d'un engrenage.
\item À la fin de cette leçon, je sais déterminer si un engrenage est réducteur ou multiplicateur de vitesse.
\item À la fin de cette leçon, je sais interpréter un schéma de transmission et résoudre un calcul simple de vitesse de rotation.
\item À la fin de cette leçon, je sais appliquer les consignes de sécurité liées à l'utilisation des machines à engrenages.
\end{itemize}
\end{objectifs}

\begin{prerequis}
\begin{itemize}
\item Notion de mouvement de rotation et de vitesse de rotation (tours par minute)
\item Notion de transmission de mouvement (leçon sur les poulies et courroies)
\item Calculs de proportionnalité et de rapports (mathématiques, 4ème-3ème)
\item Lecture et interprétation de schémas techniques simples
\end{itemize}
\end{prerequis}

\newpage

\coursec{Découverte expérimentale des engrenages}

\courssub{Situation d'accroche : les roues dentées sont partout}

Au quartier, le réparateur de vélos explique à un client pourquoi son vélo monte plus facilement les côtes quand il passe sur le grand plateau arrière. À l'atelier de menuiserie de Nkongsamba, la perceuse à colonne change de vitesse grâce à des poulies et des roues dentées. Dans le moulin à manioc du village, un moteur entraîne la meule par un système de roues dentées. Le soir, en démontant une vieille montre qui ne fonctionne plus, on découvre une dizaine de petites roues en laiton, toutes garnies de dents, qui s'emboîtent les unes dans les autres.

Partout, la même idée revient : des \cle{roues dentées} qui tournent en s'entraînant mutuellement. Mais comment ces roues transmettent-elles le mouvement ? Pourquoi la vitesse change-t-elle quand on change de roue ? Et pourquoi le réparateur de vélos parle-t-il de « petit » ou « grand » plateau ?

\textbf{Problématique.} Comment les roues dentées transmettent-elles un mouvement de rotation d'un arbre à un autre, et que devient la vitesse de rotation au passage ?

Pour répondre à ces questions, nous n'allons pas croire le professeur sur parole : nous allons \emph{observer}, \emph{manipuler} et \emph{compter}. C'est la démarche d'investigation : observation, hypothèse, expérience, interprétation, conclusion.

\courssub{Observation d'objets de la vie courante}

Avant toute expérience, ouvrons les yeux autour de nous. De nombreux objets familiers contiennent des roues dentées.

\begin{itemize}
\item \textbf{La montre mécanique} : un ressort fournit l'énergie ; une cascade de roues dentées de tailles différentes fait tourner l'aiguille des secondes 60 fois plus vite que celle des minutes, et celle des minutes 12 fois plus vite que celle des heures.
\item \textbf{La perceuse} (à main ou à colonne) : le moteur tourne très vite ; des roues dentées réduisent la vitesse pour que le foret tourne plus lentement mais avec plus de force, afin de percer le bois ou le métal.
\item \textbf{La boîte de vitesses} de la moto-taxi ou de la voiture : le conducteur choisit un couple de roues dentées ; en première, le moteur tourne vite et les roues lentement (force pour démarrer ou monter une côte) ; en quatrième, c'est l'inverse (vitesse sur route plate).
\item \textbf{Le moulin à manioc ou à maïs} : le moteur entraîne la meule par des roues dentées ou des poulies ; la meule tourne à la bonne vitesse pour écraser le produit sans s'emballer.
\item \textbf{Le dérailleur du vélo} : la chaîne passe sur des plateaux et des pignons de diamètres différents ; changer de vitesse, c'est choisir une nouvelle combinaison de roues dentées.
\end{itemize}

\begin{popfigure}
\begin{center}
\begin{tikzpicture}[scale=0.9, every node/.style={font=\small}]
% Cadre du vélo simplifié
\draw[thick, popdark] (0,0) circle (0.9);
\draw[thick, popdark] (6.2,0) circle (0.9);
\draw[thick, popdark] (0,0) -- (2.6,1.9) -- (6.2,0);
\draw[thick, popdark] (0,0) -- (3.4,0.55) -- (6.2,0);
\draw[thick, popdark] (2.6,1.9) -- (3.4,0.55);
\draw[thick, popdark] (2.6,1.9) -- (2.35,2.5);
\draw[thick, popdark] (6.2,0) -- (6.6,1.7) -- (7.3,1.75);
% Plateau (pédalier)
\fill[popblue] (3.4,0.55) circle (0.55);
\foreach \a in {0,30,...,330} \fill[popblue] (3.4,0.55) ++(\a:0.55) circle (0.06);
\fill[popcream] (3.4,0.55) circle (0.35);
% Pignon arrière
\fill[poporange] (6.2,0) circle (0.28);
\foreach \a in {0,45,...,315} \fill[poporange] (6.2,0) ++(\a:0.28) circle (0.05);
\fill[popcream] (6.2,0) circle (0.15);
% Chaîne
\draw[thick, popmuted] (3.4,1.1) -- (6.2,0.28);
\draw[thick, popmuted] (3.4,0.0) -- (6.2,-0.28);
% Pédales
\draw[thick, popdark] (3.4,0.55) -- (3.0,0.1);
\fill[popdark] (2.95,0.05) rectangle (3.15,0.15);
% Légendes fléchées
\node[popblue, font=\small\bfseries] at (1.6,-1.3) {1. Plateau (roue dentée menante)};
\draw[->, popblue, thick] (2.4,-1.1) -- (3.2,0.15);
\node[poporange, font=\small\bfseries] at (7.6,-1.3) {2. Pignon (roue dentée menée)};
\draw[->, poporange, thick] (7.0,-1.1) -- (6.35,-0.25);
\node[popmuted, font=\small\bfseries] at (4.8,1.6) {3. Chaîne};
\draw[->, popmuted, thick] (4.8,1.4) -- (4.8,0.75);
\end{tikzpicture}
\end{center}
\legende{Schéma simplifié de la transmission d'un vélo : le pédalier (1) entraîne, par la chaîne (3), le pignon (2) fixé à la roue arrière. Plateau et pignon sont des roues dentées.}
\end{popfigure}

\begin{attention}[Avant de manipuler : la règle d'or de l'investigation]
Ne jamais affirmer les résultats avant d'avoir observé. Dire « je sais déjà que la petite roue tourne plus vite » sans l'avoir vérifié, c'est faire de la mémoire, pas de la science. Dans cette leçon, chaque conclusion devra être justifiée par un \cle{constat expérimental} : ce que nous avons réellement vu et compté pendant la manipulation.
\end{attention}

\courssub{Expérience de classe : deux roues dentées engrenées}

\begin{experience}[Deux roues dentées en carton]
\textbf{Matériel} : deux roues dentées découpées dans du carton rigide ou en plastique (par exemple une grande roue de 24 dents et une petite de 12 dents), deux pointes ou clous pour servir d'axes, une planche support, un feutre.

\textbf{Protocole} :
\begin{enumerate}
\item Fixer les deux roues sur la planche par leurs axes, de façon que les dents de l'une pénètrent dans les creux de l'autre : on dit que les roues \cle{engrenent}.
\item Marquer d'un trait de feutre une dent de repère sur chaque roue.
\item Faire tourner lentement la grande roue, appelée roue \cle{menante} (celle qui commande), et observer la petite roue, appelée roue \cle{menée} (celle qui est entraînée).
\item Compter les tours : combien de tours fait la petite roue pendant que la grande fait exactement un tour ?
\end{enumerate}

\textbf{Observations} :
\begin{itemize}
\item Les deux roues tournent en \textbf{sens opposés} : si la menante tourne dans le sens des aiguilles d'une montre, la menée tourne dans le sens inverse.
\item La petite roue tourne \textbf{plus vite} : pendant que la grande roue de 24 dents fait 1 tour, la petite de 12 dents fait 2 tours.
\item Les dents ne glissent pas les unes sur les autres : une dent de la menante pousse exactement une dent de la menée.
\end{itemize}
\end{experience}

\begin{methode}[Bien mener l'expérience d'observation]
\begin{enumerate}
\item \textbf{Compter les dents} de chaque roue avant de commencer, et noter les nombres : $Z_1$ pour la roue menante, $Z_2$ pour la roue menée.
\item \textbf{Marquer un repère} (trait de feutre) sur une dent de chaque roue pour repérer facilement un tour complet.
\item \textbf{Faire tourner lentement} la roue menante, en comptant ses tours à voix haute.
\item \textbf{Compter les tours} de la roue menée pendant un nombre entier de tours de la menante.
\item \textbf{Noter les résultats} dans un tableau, puis seulement après, chercher une conclusion.
\end{enumerate}
\end{methode}

\begin{popfigure}
\begin{center}
\begin{tikzpicture}[scale=1.0]
% Grande roue (menante)
\foreach \a in {0,15,...,345}
  \fill[popblue] (0,0) ++(\a:1.18) circle (0.09);
\fill[popblue!70] (0,0) circle (1.05);
\fill[popcream] (0,0) circle (0.85);
\fill[popblue] (0,0) circle (0.12);
% Repère menante
\fill[popink] (0,0) ++(90:1.18) circle (0.14);
% Petite roue (menée) : centre à distance 1.18+0.62 = 1.80
\foreach \a in {0,30,...,330}
  \fill[poporange] (1.80,0) ++(\a:0.60) circle (0.09);
\fill[poporange!75] (1.80,0) circle (0.48);
\fill[popcream] (1.80,0) circle (0.34);
\fill[poporange] (1.80,0) circle (0.10);
% Repère menée
\fill[popink] (1.80,0) ++(180:0.60) circle (0.13);
% Flèches de rotation
\draw[->, very thick, popdark] (-0.55,1.55) arc (150:30:0.75);
\node[popdark, font=\small\bfseries] at (-0.9,1.7) {sens 1};
\draw[->, very thick, popdark] (2.5,0.95) arc (60:150:0.85);
\node[popdark, font=\small\bfseries] at (3.1,1.15) {sens 2 (inverse)};
% Légendes
\node[popblue, font=\small\bfseries] at (0,-1.75) {Roue menante : 24 dents};
\node[poporange, font=\small\bfseries] at (2.6,-1.75) {Roue menée : 12 dents};
% Axes
\node[font=\footnotesize] at (0,0.28) {axe};
\node[font=\footnotesize] at (1.80,0.24) {axe};
\end{tikzpicture}
\end{center}
\legende{Deux roues dentées engrenées. Les dents de la roue menante (bleue) poussent celles de la roue menée (orange) : les deux roues tournent en sens opposés. Les points noirs sont les repères de feutre pour compter les tours.}
\end{popfigure}

\courssub{Les trois constats expérimentaux}

Reprenons calmement chaque observation et transformons-la en constat bien formulé. Un constat est une phrase qui décrit \emph{uniquement} ce qu'on a vu, sans interprétation.

\textbf{Constat 1 : les sens de rotation sont opposés.} Quand la roue menante tourne dans le sens des aiguilles d'une montre, la roue menée tourne dans le sens inverse. C'est logique : au point de contact, les dents de la menante poussent celles de la menée vers le bas si elles-mêmes descendent... mais du côté opposé de l'axe de la menée, ce qui la fait tourner à l'envers. Pour un engrenage \cle{extérieur} (dents à l'extérieur des roues, le cas le plus courant), deux roues qui engrenent tournent toujours en sens contraires.

\textbf{Constat 2 : la roue qui a moins de dents tourne plus vite.} Avec une menante de 24 dents et une menée de 12 dents, la menée fait 2 tours pendant que la menante en fait 1. Pourquoi ? Parce que les dents passent une par une au point de contact : quand 24 dents de la menante sont passées (1 tour complet), 24 dents de la menée ont aussi été poussées... mais la menée n'en possède que 12, donc elle a fait 2 tours ! Le nombre de dents qui passent au contact est le même pour les deux roues : c'est la clé de tout le mécanisme.

\textbf{Constat 3 : la transmission se fait sans glissement.} Contrairement à une courroie tendue entre deux poulies, qui peut \emph{patiner} quand l'effort est trop grand (on entend parfois un couinement au démarrage d'un moteur), les dents s'emboîtent : elles ne peuvent pas glisser sans se casser. Le mouvement est donc transmis de façon \cle{rigoureuse} : si la menante avance d'une dent, la menée avance exactement d'une dent. C'est pour cela que les moteurs qui doivent rester synchronisés (dans une horloge, par exemple) utilisent des engrenages et non de simples frottements.

\begin{exoresolu}[Compter les tours]
Sur le dispositif de la classe, la roue menante possède 24 dents et la roue menée 12 dents. On fait faire 3 tours complets à la roue menante. Combien de tours la roue menée effectue-t-elle ? Dans quel sens tourne-t-elle si la menante tourne dans le sens des aiguilles d'une montre ?
\tcblower
\textbf{Solution détaillée.}
\begin{enumerate}
\item En un tour, la menante fait passer 24 dents au point de contact. En 3 tours, elle fait passer $3 \times 24 = 72$ dents.
\item La menée reçoit ces 72 dents une par une. Comme elle possède 12 dents, elle effectue :
\[ \frac{72}{12} = 6 \text{ tours.} \]
\item L'engrenage est extérieur : la menée tourne donc en \textbf{sens inverse} des aiguilles d'une montre.
\end{enumerate}
\textbf{Conclusion} : la roue menée fait 6 tours, en sens inverse de la roue menante. On remarque qu'elle tourne 2 fois plus vite, car elle a 2 fois moins de dents.
\end{exoresolu}

\begin{attention}[Pièges fréquents]
\begin{itemize}
\item Ne confonds pas \textbf{menante} et \textbf{menée} : la menante est celle qu'on fait tourner (le moteur, la pédale), la menée est celle qui reçoit le mouvement.
\item Ne dis pas « la petite roue tourne plus vite » sans préciser \emph{laquelle est menante} : si c'est la petite qui commande, c'est la grande qui tourne plus lentement !
\item Le sens inverse n'est vrai que pour un engrenage \textbf{extérieur} avec un nombre impair de roues... retiens simplement : deux roues extérieures qui engrenent tournent en sens contraires.
\end{itemize}
\end{attention}

\courssub{Conclusion de la découverte}

\begin{aretenir}[Ce que l'expérience nous a appris]
Des roues dentées qui engrenent \cle{transmettent un mouvement de rotation} d'un arbre vers un autre arbre proche. Au passage :
\begin{itemize}
\item le \textbf{sens} de rotation est inversé (engrenage extérieur) ;
\item la \textbf{vitesse} de rotation change : la roue qui a le moins de dents tourne le plus vite ;
\item la transmission se fait \textbf{sans glissement}, dent par dent.
\end{itemize}
C'est exactement ce qui se passe dans le vélo du quartier, la perceuse de Nkongsamba et le moulin à manioc : on choisit la taille des roues dentées pour obtenir la vitesse (et la force) désirées.
\end{aretenir}

\begin{savaistu}[Des engrenages vieux de plus de 2000 ans]
Les engrenages ne sont pas une invention moderne ! Il y a plus de 2000 ans, les moulins à eau utilisaient déjà de grandes roues dentées en bois pour transformer la rotation verticale de la roue à aubes en rotation de la meule qui écrasait le grain. Les horloges mécaniques, apparues au Moyen Âge, doivent leur précision à des trains d'engrenages en laiton. Plus extraordinaire encore : le « mécanisme d'Anticythère », retrouvé dans une épave grecque et datant d'environ 100 ans avant J.-C., contenait plus de 30 engrenages en bronze qui calculaient les positions du Soleil et de la Lune. Les roues dentées de ton vélo sont donc les héritières d'une très longue histoire !
\end{savaistu}

\begin{exercice}[À toi de jouer]
Dans le moulin à manioc du village, le moteur fait tourner une roue dentée de 40 dents qui engrene avec une roue de 10 dents fixée sur l'arbre de la meule.
\begin{enumerate}
\item Quelle roue est la menante ? Laquelle est la menée ?
\item La meule tourne-t-elle dans le même sens que le moteur ? Justifie.
\item Pendant que le moteur fait 1 tour, combien de tours fait la meule ?
\item Le moteur tourne à 300 tours par minute. Calcule la vitesse de rotation de la meule en tours par minute.
\end{enumerate}
\end{exercice}

\begin{corrige}[Exercice : le moulin à manioc]
\begin{enumerate}
\item La roue de 40 dents, entraînée par le moteur, est la \textbf{menante} ; la roue de 10 dents, fixée à la meule, est la \textbf{menée}.
\item Non : il s'agit d'un engrenage extérieur, donc la meule tourne en \textbf{sens inverse} du moteur.
\item En 1 tour, la menante fait passer 40 dents ; la menée, qui n'a que 10 dents, fait donc $\frac{40}{10} = 4$ tours.
\item La meule tourne 4 fois plus vite que le moteur : $300 \times 4 = 1200$ tours par minute.
\end{enumerate}
\end{corrige}

\coursec{Définitions et vocabulaire}

Dans la section précédente, tu as observé des engrenages en action : deux roues crantées qui tournent ensemble, comme dans une vieille horloge, un moulin à manioc ou la boîte de vitesses d'une moto-taxi. Tu as constaté que lorsqu'une roue tourne, elle entraîne l'autre. Pour décrire précisément ces mécanismes — et surtout pour réussir les exercices du BEPC — il faut maintenant apprendre le \cle{vocabulaire} exact des techniciens et des ingénieurs. C'est l'objet de cette section : donner un nom rigoureux à chaque élément et définir chaque notion.

\courssub{La roue dentée}

\textbf{Observation préalable.}\par

Prends le temps de regarder attentivement une roue d'engrenage, par exemple celle d'un vieux réveil mécanique ou d'un jouet démonté. Que remarques-tu ? Ce n'est pas une simple roue lisse comme une roue de bicyclette : tout autour de son bord, on voit des saillies identiques, régulièrement espacées, qu'on appelle des \cle{dents}. C'est précisément cette particularité qui lui permet d'entraîner une autre roue sans glisser.

\begin{definition}[Roue dentée]
Une \cle{roue dentée} est une roue qui comporte, sur sa périphérie, des \cle{dents} régulièrement espacées, destinées à s'engrêner avec les dents d'une autre roue pour transmettre un mouvement de rotation.
\end{definition}

Deux caractéristiques essentielles décrivent une roue dentée :

\begin{itemize}
\item son \cle{nombre de dents}, noté \cle{$Z$} : c'est tout simplement le nombre total de dents réparties sur toute la périphérie de la roue. Par exemple, une roue qui porte 20 dents a $Z = 20$ ;
\item son \cle{diamètre primitif}, noté $d$ : c'est le diamètre du cercle de référence (appelé \emph{cercle primitif}) sur lequel les dents des deux roues s'engrènent. On le représente en pointillés sur les schémas. Ce cercle imaginaire passe « au milieu » des dents, ni tout à fait à leur base, ni tout à fait à leur sommet.
\end{itemize}

\begin{attention}[Le diamètre primitif est un savoir admis]
Au programme de la classe de 3ème, la notion de diamètre primitif est simplement \emph{présentée} : tu dois savoir le repérer sur un schéma (cercle en pointillés) et comprendre qu'il sert de référence pour l'engrènement. On ne te demandera pas de le calculer au BEPC.
\end{attention}

\begin{popfigure}
\begin{center}
\begin{tikzpicture}[scale=1.05]
% Cercle primitif en pointillés
\draw[dashed,popblue,thick] (0,0) circle (1.5);
% Corps de la roue
\draw[fill=popblueL,thick,popdark] (0,0) circle (1.25);
% Dents (12 dents)
\foreach \a in {0,30,...,330}{
  \begin{scope}[rotate=\a]
    \draw[fill=poporangeL,thick,popdark] (1.15,-0.16) rectangle (1.75,0.16);
  \end{scope}
}
% Axe
\draw[fill=white,thick] (0,0) circle (0.18);
\fill[popdark] (0,0) circle (0.05);
% Cote diamètre primitif
\draw[<->,popblue,thick] (-1.5,0) -- (1.5,0);
\node[popblue,below] at (0,-0.08) {\small diamètre primitif $d$};
% Légendes fléchées
\draw[->,popdark] (2.6,1.6) -- (1.75,1.05);
\node[popdark,right] at (2.6,1.6) {\small une dent};
\draw[->,popdark] (-2.8,-1.4) -- (-1.15,-0.7);
\node[popdark,left] at (-2.8,-1.4) {\small corps de la roue};
\node[popdark] at (0,-2.3) {\small Nombre de dents : $Z = 12$};
\end{tikzpicture}
\end{center}
\legende{Schéma d'une roue dentée : les dents sont régulièrement espacées sur la périphérie ; le cercle primitif (pointillés bleus) est le cercle de référence de l'engrènement ; le nombre de dents se note $Z$.}
\end{popfigure}

\courssub{L'engrenage}

\textbf{De l'observation à la définition.}\par

Reprenons la démarche d'investigation de la séance précédente. Tu as posé deux roues dentées côte à côte, dents contre dents, puis tu as fait tourner l'une d'elles à la main. \emph{Observation} : la deuxième roue se met à tourner, sans qu'on la touche. \emph{Interprétation} : les dents de la première roue poussent les dents de la seconde, comme les maillons d'une chaîne. Le mouvement de rotation est donc \emph{transmis} d'une roue à l'autre. \emph{Conclusion} : deux roues dentées dont les dents s'engrènent forment un mécanisme de transmission de mouvement. Ce mécanisme porte un nom : l'engrenage.

\begin{definition}[Engrenage]
Un \cle{engrenage} est un ensemble de deux roues dentées (ou plus) dont les dents \cle{s'engrènent}, c'est-à-dire sont en prise les unes avec les autres, de manière à transmettre un \cle{mouvement de rotation} d'une roue à l'autre.
\end{definition}

\begin{aretenir}[L'idée essentielle]
Dans un engrenage, il n'y a \textbf{pas de glissement} entre les roues : chaque dent de la première roue pousse exactement une dent de la seconde. C'est ce qui distingue l'engrenage de deux simples roues lisses posées l'une contre l'autre, qui pourraient patiner.
\end{aretenir}

\courssub{Roue menante et roue menée}

Dans un engrenage, les deux roues ne jouent pas le même rôle. Pour bien les distinguer, pose-toi toujours la question : \emph{d'où vient le mouvement ?}

\begin{definition}[Roue menante et roue menée]
\begin{itemize}
\item La \cle{roue menante} (ou roue \emph{motrice}) est celle qui \textbf{reçoit le mouvement} du moteur ou de l'opérateur (la main, une manivelle, une pédale). C'est elle qui \emph{donne} le mouvement.
\item La \cle{roue menée} (ou roue \emph{réceptrice}) est celle qui \textbf{reçoit le mouvement de la roue menante}. C'est elle qui \emph{reçoit} le mouvement et le transmet à la suite du mécanisme.
\end{itemize}
\end{definition}

\begin{exemplebox}[La perceuse]
Dans une perceuse électrique, le moteur fait tourner un petit pignon solidaire de son axe. Ce pignon engrène avec une roue dentée plus grande, solidaire de l'axe qui porte le foret. Ici :
\begin{itemize}
\item le petit pignon du moteur est la \cle{roue menante} : c'est le moteur qui lui donne le mouvement ;
\item la grande roue est la \cle{roue menée} : elle reçoit le mouvement du pignon et entraîne le foret.
\end{itemize}
\end{exemplebox}

\begin{attention}[Erreur fréquente : confondre roue menante et roue menée]
Sur un schéma, beaucoup d'élèves désignent la roue menante « au hasard », souvent en choisissant la plus grande. C'est faux ! La taille ne décide pas du rôle. \textbf{Méthode infaillible} : repère d'où vient le mouvement (moteur, manivelle, pédale, flèche d'entrée sur le schéma). La roue reliée à cette source de mouvement est la menante ; l'autre est la menée. Au BEPC, lis toujours l'énoncé jusqu'au bout avant de nommer les roues.
\end{attention}

\courssub{Le pignon et le nombre de dents}

Lorsque les deux roues d'un engrenage n'ont pas la même taille — ce qui est le cas le plus fréquent — on donne un nom particulier à la plus petite.

\begin{definition}[Pignon]
Dans un engrenage de deux roues, le \cle{pignon} est la \textbf{plus petite} des deux roues dentées. La plus grande est simplement appelée la \emph{roue} (ou la couronne).
\end{definition}

Pour distinguer les deux roues dans les calculs et les schémas, on numérote leurs nombres de dents :

\begin{itemize}
\item \cle{$Z_1$} : nombre de dents de la roue \textbf{menante} ;
\item \cle{$Z_2$} : nombre de dents de la roue \textbf{menée}.
\end{itemize}

\begin{exemplebox}[Exemple de notation]
Un engrenage est constitué d'un pignon menant de $15$ dents et d'une roue menée de $45$ dents. On écrit : $Z_1 = 15$ (menante) et $Z_2 = 45$ (menée). Le pignon est ici la roue menante, car c'est la plus petite : $Z_1 < Z_2$.
\end{exemplebox}

\begin{popfigure}
\begin{center}
\begin{tikzpicture}[scale=0.95]
% ===== Roue menante (petite, à gauche) =====
\begin{scope}[shift={(-1.55,0)}]
  \draw[dashed,popblue,thick] (0,0) circle (0.9);
  \draw[fill=poporangeL,thick,popdark] (0,0) circle (0.72);
  \foreach \a in {0,45,...,315}{
    \begin{scope}[rotate=\a]
      \draw[fill=poporange,thick,popdark] (0.62,-0.11) rectangle (1.06,0.11);
    \end{scope}
  }
  \draw[fill=white,thick] (0,0) circle (0.13);
  \draw[->,popdark,thick] (0.35,0.55) arc (60:150:0.65);
\end{scope}
% ===== Roue menée (grande, à droite) =====
\begin{scope}[shift={(1.75,0)}]
  \draw[dashed,popblue,thick] (0,0) circle (1.35);
  \draw[fill=popgreenL,thick,popdark] (0,0) circle (1.12);
  \foreach \a in {22.5,67.5,...,337.5}{
    \begin{scope}[rotate=\a]
      \draw[fill=popgreen,thick,popdark] (1.0,-0.11) rectangle (1.5,0.11);
    \end{scope}
  }
  \draw[fill=white,thick] (0,0) circle (0.13);
  \draw[->,popdark,thick] (-0.5,0.85) arc (120:30:0.95);
\end{scope}
% Étiquettes
\node[popdark,align=center] at (-1.55,-1.85) {\small \textbf{Roue menante} (pignon)\\ \small $Z_1 = 8$ dents};
\node[popdark,align=center] at (1.75,-1.85) {\small \textbf{Roue menée}\\ \small $Z_2 = 12$ dents};
% Flèche mouvement moteur
\draw[->,very thick,popink] (-3.4,0) -- (-2.75,0);
\node[popink,left,align=center] at (-3.4,0) {\small mouvement\\ \small du moteur};
% Point de contact
\fill[poppink] (-0.55,0) circle (0.07);
\node[poppink,above] at (-0.55,0.12) {\small engrènement};
\end{tikzpicture}
\end{center}
\legende{Engrenage à deux roues : le mouvement du moteur arrive sur la roue menante (ici le pignon, $Z_1 = 8$) ; les dents s'engrènent au point de contact et entraînent la roue menée ($Z_2 = 12$), qui tourne en sens inverse.}
\end{popfigure}

\courssub{Condition de fonctionnement : le même module}

\textbf{Une expérience qui ne marche pas.}\par

Essaie de faire engrèner deux roues dentées dont les dents n'ont pas la même taille : par exemple une roue à grandes dents espacées et une roue à petites dents serrées. \emph{Observation} : les dents se bloquent, se chevauchent ou sautent ; la rotation est impossible ou irrégulière. \emph{Conclusion} : deux roues ne peuvent engrèner correctement que si leurs dents ont la \textbf{même taille}.

\begin{propriete}[Condition d'engrènement]
Pour que deux roues dentées engrènent correctement, elles doivent avoir le \cle{même module}, c'est-à-dire des dents de \textbf{même taille} (même hauteur et même espacement). Le module, noté $m$, est le nombre qui caractérise la taille des dents ; il est le même pour toutes les roues d'un même engrenage.
\end{propriete}

\begin{attention}[Ne pas confondre « même module » et « même nombre de dents »]
Deux roues d'un engrenage ont presque toujours des nombres de dents \emph{différents} ($Z_1 \neq Z_2$), et c'est justement ce qui permet de modifier la vitesse de rotation. Mais leurs dents doivent avoir la \emph{même taille} (même module). Retiens : \textbf{même module, obligatoire ; même nombre de dents, pas du tout nécessaire}.
\end{attention}

\begin{exoresolu}[Identifier les éléments d'un engrenage]
\textbf{Énoncé.} Sur la boîte de vitesses d'une moto-taxi, un mécanicien observe un engrenage : l'arbre qui vient du moteur porte une roue de $20$ dents qui engrène avec une roue de $60$ dents fixée sur l'arbre de sortie.
\begin{enumerate}
\item Quelle est la roue menante ? La roue menée ?
\item Quelle roue est le pignon ?
\item Donner les valeurs de $Z_1$ et $Z_2$.
\item Les deux roues ont des dents de tailles différentes. L'engrenage peut-il fonctionner correctement ?
\end{enumerate}
\tcblower
\textbf{Solution détaillée.}
\begin{enumerate}
\item Le mouvement vient du moteur. La roue de $20$ dents, fixée sur l'arbre moteur, est donc la \textbf{roue menante}. La roue de $60$ dents, qui reçoit le mouvement de la première, est la \textbf{roue menée}.
\item Le \textbf{pignon} est la plus petite des deux roues : c'est la roue de $20$ dents (ici, c'est aussi la roue menante).
\item Par convention, l'indice $1$ désigne la menante et l'indice $2$ la menée : $Z_1 = 20$ et $Z_2 = 60$.
\item \textbf{Non.} Si les dents n'ont pas la même taille, les deux roues n'ont pas le même module : les dents se bloqueront ou sauteront. Pour fonctionner, un engrenage exige le \textbf{même module} pour les deux roues.
\end{enumerate}
\end{exoresolu}

\begin{savaistu}[Le module, une grandeur normalisée]
Les fabricants du monde entier utilisent des modules normalisés (en millimètres) : $m = 1$ ; $1,5$ ; $2$ ; $3$\ldots{} Ainsi, un pignon de rechange acheté à Douala peut remplacer un pignon usé d'une machine importée, pourvu qu'on choisisse le bon module et le bon nombre de dents. C'est la \emph{normalisation} qui rend possible la réparation des machines, des moulins à maïs du village aux groupes électrogènes.
\end{savaistu}

\begin{aretenir}[Bilan du vocabulaire à connaître par cœur]
\begin{itemize}
\item \textbf{Roue dentée} : roue à dents régulièrement espacées sur sa périphérie.
\item \textbf{Engrenage} : ensemble de roues dentées en prise transmettant un mouvement de rotation.
\item \textbf{Roue menante} : reçoit le mouvement du moteur ou de l'opérateur ($Z_1$ dents).
\item \textbf{Roue menée} : reçoit le mouvement de la roue menante ($Z_2$ dents).
\item \textbf{Pignon} : la plus petite roue de l'engrenage.
\item \textbf{Diamètre primitif} : diamètre du cercle de référence de l'engrènement (cercle en pointillés).
\item \textbf{Condition de fonctionnement} : les deux roues doivent avoir le \textbf{même module}.
\end{itemize}
Ces définitions sont la base indispensable pour aborder la suite de la leçon : les types d'engrenages, le sens de rotation et le rapport de transmission.
\end{aretenir}

\coursec{Types d'engrenages et sens de rotation}

\courssub{Transition depuis le vocabulaire}
Dans la section précédente, nous avons appris à nommer les éléments d'une roue dentée : le \cle{primitif}, le \cle{pied}, la \cle{tête}, le \cle{module} et le \cle{pas}. Nous savons maintenant \emph{ce qu'est} une dent. Voyons maintenant \emph{comment} on assemble ces roues pour transmettre le mouvement. Le type d'assemblage détermine le sens de rotation des arbres et la direction des forces. C'est ce que nous allons découvrir par l'observation et le schéma.

\courssub{Engrenage extérieur : sens de rotation inverses}
C'est le montage le plus courant. Deux roues dentées sont placées côte à côte, leurs dents s'engrènent \emph{extérieurement}. Leurs axes de rotation sont \textbf{parallèles}.

\begin{experience}[Observation du sens de rotation d'un engrenage extérieur]
\textbf{Matériel :} Deux roues dentées de modules identiques (ex. $m=2$) fixées sur un support (planchette), un feutre pour marquer une dent de référence sur chaque roue.
\textbf{Protocole :}
\begin{enumerate}
\item Marquer une dent sur la roue motrice (R1) et une dent sur la roue menée (R2).
\item Faire tourner lentement la roue motrice (R1) dans le sens des aiguilles d'une montre (sens horaire).
\item Observer le sens de rotation de la roue menée (R2).
\end{enumerate}
\textbf{Observations :} Lorsque la roue motrice tourne dans le sens horaire, la roue menée tourne dans le sens \emph{anti-horaire} (trigonométrique). Les flèches tangentielles au point de contact pointent dans des directions opposées.
\textbf{Interprétation :} Au point de contact, la dent de la roue motrice \emph{pousse} la dent de la roue menée. Pour que le contact roule sans glisser, les vitesses linéaires au primitif sont égales et de sens contraires. Les axes étant parallèles, les rotations sont inverses.
\textbf{Conclusion :} Dans un engrenage extérieur, les deux roues tournent en \textbf{sens inverses}.
\end{experience}

\begin{popfigure}
\begin{tikzpicture}[scale=0.9, >=Stealth]
% Roue 1 (gauche)
\def\R{2.5} % rayon primitif
\def\r{2.0} % rayon pied
\def\Rt{2.9} % rayon tete
\def\n{12} % nombre de dents
\foreach \i in {1,...,\n} {
  \pgfmathsetmacro{\angle}{360/\n*\i}
  \draw[popblue, thick] (\angle:\r) -- (\angle:\Rt);
  \draw[popblue, fill=popblue!20] (\angle:\Rt) arc (\angle:\angle+360/\n*0.5:\Rt) -- (\angle+360/\n*0.5:\r) arc (\angle+360/\n*0.5:\angle:\r) -- cycle;
}
\draw[popblue, thick] (0,0) circle (\R);
\draw[popdark, thick] (0,0) circle (0.3);
\draw[->, poporange, thick] (0,0) -- ++(90:1.5) node[above] {$N_1$};
\draw[->, poporange, thick] (1.5,0) arc (0:90:1.5) node[right] {Sens horaire};
\node[popdark] at (0,-3.5) {Roue motrice (R1)};

% Roue 2 (droite) - centre à distance R1+R2 = 5
\begin{scope}[xshift=5cm]
\foreach \i in {1,...,\n} {
  \pgfmathsetmacro{\angle}{360/\n*\i + 180/\n} % décalage demi-dent pour engrener
  \draw[popgreen!70!black, thick] (\angle:\r) -- (\angle:\Rt);
  \draw[popgreen!70!black, fill=popgreen!20] (\angle:\Rt) arc (\angle:\angle+360/\n*0.5:\Rt) -- (\angle+360/\n*0.5:\r) arc (\angle+360/\n*0.5:\angle:\r) -- cycle;
}
\draw[popgreen!70!black, thick] (0,0) circle (\R);
\draw[popdark, thick] (0,0) circle (0.3);
\draw[->, poporange, thick] (0,0) -- ++(90:1.5) node[above] {$N_2$};
\draw[->, poporange, thick] (-1.5,0) arc (180:90:1.5) node[left] {Sens anti-horaire};
\node[popdark] at (0,-3.5) {Roue menée (R2)};
\end{scope}

% Flèches vitesse au contact
\draw[->, popred, very thick] (2.5,0) -- ++(0,1) node[above, pos=0.5, xshift=0.3cm] {$\vec{v}_1$};
\draw[->, popred, very thick] (2.5,0) -- ++(0,-1) node[below, pos=0.5, xshift=0.3cm] {$\vec{v}_2$};
\draw[dashed, popmuted] (0,0) -- (5,0);
\node[popmuted] at (2.5, -0.5) {Point de contact C};
\end{tikzpicture}
\legende{Engrenage extérieur : les deux axes sont parallèles, les roues tournent en sens inverses. Les vitesses linéaires au point de contact C sont opposées.}
\end{popfigure}

\begin{propriete}[Sens de rotation en engrenage extérieur]
Dans un engrenage à denture extérieure (axes parallèles), la roue motrice et la roue menée tournent en \textbf{sens opposés}.
Si la roue motrice tourne dans le sens horaire, la roue menée tourne dans le sens anti-horaire, et inversement.
\end{propriete}

\courssub{Engrenage intérieur (couronne) : même sens de rotation}
Ici, la roue menée est une \cle{couronne dentée intérieurement}. Le pignon (petite roue) s'engrène \emph{à l'intérieur} de la couronne. Les axes restent parallèles.

\begin{popfigure}
\begin{tikzpicture}[scale=0.9, >=Stealth]
% Couronne (grande roue creuse) - Centre (0,0)
\def\Rc{3.5} % rayon primitif couronne
\def\rc{3.0} % rayon pied couronne (intérieur)
\def\Rct{4.0} % rayon tête couronne (extérieur)
\def\nc{24}
% Anneau extérieur
\draw[poppurple, thick] (0,0) circle (\Rct);
% Cercle primitif
\draw[poppurple, thick] (0,0) circle (\Rc);
% Cercle pied (intérieur)
\draw[poppurple, thick] (0,0) circle (\rc);
% Dents intérieures (vers l'intérieur)
\foreach \i in {1,...,\nc} {
  \pgfmathsetmacro{\angle}{360/\nc*\i}
  \draw[poppurple, thick] (\angle:\rc) -- (\angle:\Rc-0.2);
  \draw[poppurple, fill=poppurple!20] (\angle:\rc) -- (\angle:\Rc-0.2) -- (\angle+360/\nc*0.4:\Rc-0.2) -- (\angle+360/\nc*0.4:\rc) -- cycle;
}
\draw[popdark, thick] (0,0) circle (0.3);
\draw[->, poporange, thick] (0,0) -- ++(90:1.8) node[above] {$N_2$};
\draw[->, poporange, thick] (-2,0) arc (180:90:2) node[left] {Sens horaire};
\node[popdark] at (0,-4.5) {Couronne menée (R2)};

% Pignon (petite roue) à l'intérieur - Centre (0, -2.0) car Rc - Rp = 3.5 - 1.5 = 2.0
\def\Rp{1.5}
\def\rp{1.0}
\def\Rpt{1.9}
\def\np{10}
\begin{scope}[yshift=-2.0cm]
\foreach \i in {1,...,\np} {
  \pgfmathsetmacro{\angle}{360/\np*\i}
  \draw[popblue, thick] (\angle:\rp) -- (\angle:\Rpt);
  \draw[popblue, fill=popblue!20] (\angle:\Rpt) arc (\angle:\angle+360/\np*0.5:\Rpt) -- (\angle+360/\np*0.5:\rp) arc (\angle+360/\np*0.5:\angle:\rp) -- cycle;
}
\draw[popblue, thick] (0,0) circle (\Rp);
\draw[popdark, thick] (0,0) circle (0.3);
\draw[->, poporange, thick] (0,0) -- ++(90:1.2) node[above] {$N_1$};
\draw[->, poporange, thick] (1.2,0) arc (0:90:1.2) node[right] {Sens horaire};
\node[popdark] at (0,-2.5) {Pignon moteur (R1)};
\end{scope}

% Flèches vitesse au contact (bas, point C à (0, -3.5))
\draw[->, popred, very thick] (0,-3.5) -- ++(0,-1) node[below, pos=0.5, xshift=0.3cm] {$\vec{v}_1$};
\draw[->, popred, very thick] (0,-3.5) -- ++(0,-1) node[below, pos=0.5, xshift=-0.3cm] {$\vec{v}_2$};
\node[popmuted] at (0, -3.5) {$\bullet$} node[above=2pt, popmuted] {C};
\draw[dashed, popmuted] (0,0) -- (0,-3.5) -- (0,-2);
\end{tikzpicture}
\legende{Engrenage intérieur (pignon dans couronne) : axes parallèles, les deux roues tournent dans le \textbf{même sens}. Les vitesses linéaires au contact C ont même direction.}
\end{popfigure}

\begin{propriete}[Sens de rotation en engrenage intérieur]
Dans un engrenage intérieur (pignon dans une couronne), le pignon et la couronne tournent dans le \textbf{même sens}.
\end{propriete}

\begin{attention}[Piège fréquent : Confondre intérieur et extérieur]
\begin{itemize}
\item \textbf{Extérieur :} Dents face à face $\rightarrow$ Sens \textbf{inverses}. (Cas général : moto, vélo, boîte de vitesse).
\item \textbf{Intérieur :} Dents dos à dos (pignon dans couronne) $\rightarrow$ Sens \textbf{identiques}. (Cas : différentiel arrière de voiture, certains réducteurs compacts).
\end{itemize}
Au BEPC, on te demande souvent : « Préciser le sens de rotation de la roue menée. » Regarde toujours si les dents s'engrènent par l'extérieur ou par l'intérieur !
\end{attention}

\courssub{Engrenages à axes concourants : les roues coniques}
Jusqu'ici, les axes étaient parallèles. Pour transmettre le mouvement entre deux axes qui se coupent (généralement \textbf{perpendiculaires}), on utilise des \cle{roues coniques} (ou engrenages coniques). Les primitives sont des cônes qui roulent l'un sur l'autre sans glisser.

\begin{popfigure}
\begin{tikzpicture}[scale=1.1, >=Stealth]
% Représentation 2D projetée de deux cônes engrenés (axes perpendiculaires)
% Apex commun à l'origine (0,0)
% Cônes : génératrices et base
% Roue 1 : Axe horizontal (Ox), cône vers la gauche. Angle demi-sommet 45 deg (exemple).
% Base à x = -3. Rayon base = 2.5.
\def\L{3.0} % distance apex-base
\def\R{2.5} % rayon base
% Cônes : lignes de génératrices
\draw[popblue, thick] (0,0) -- (-\L, \R);  % génératrice sup
\draw[popblue, thick] (0,0) -- (-\L, -\R); % génératrice inf
\draw[popblue, thick] (-\L, -\R) -- (-\L, \R); % base
% Axe
\draw[popdark, thick] (0,0) -- (1,0); % prolongement axe
\draw[->, poporange, thick] (1,0) -- ++(1,0) node[right] {Axe $N_1$};
\draw[->, poporange, thick] (-\L, \R+0.5) arc (90:-90:0.5) node[below] {Rotation};
\node[popblue] at (-\L, -3.5) {Roue conique 1 (Horizontale)};

% Roue 2 : Axe vertical (Oy), cône vers le bas. Angle demi-sommet 45 deg.
% Base à y = -3. Rayon base = 2.5.
\draw[popgreen!70!black, thick] (0,0) -- (-\R, -\L); % génératrice gauche
\draw[popgreen!70!black, thick] (0,0) -- (\R, -\L);  % génératrice droite
\draw[popgreen!70!black, thick] (-\R, -\L) -- (\R, -\L); % base
% Axe
\draw[popdark, thick] (0,0) -- (0,1); % prolongement axe
\draw[->, poporange, thick] (0,1) -- ++(0,1) node[above] {Axe $N_2$};
\draw[->, poporange, thick] (\R+0.5, -\L) arc (0:90:0.5) node[right] {Rotation};
\node[popgreen!70!black] at (0, 3.5) {Roue conique 2 (Verticale)};

% Denture symbolique (lignes radiales sur les cônes)
\foreach \y in {-2.0, -1.5, -1.0, -0.5, 0, 0.5, 1.0, 1.5, 2.0} {
  \pgfmathsetmacro{\x}{-\L * (1 - abs(\y)/\R)}; % approx sur génératrice
  \ifdim\y pt=0pt \else \draw[popmuted, thin] (\x, \y) -- (\x*0.5, \y*0.5); \fi
}
\foreach \x in {-2.5, -2.0, -1.5, -1.0, -0.5} {
  \pgfmathsetmacro{\y}{\R * (1 + \x/\L)}; % approx
  \draw[popmuted, thin] (\x, \y) -- (\x*0.5, \y*0.5);
  \draw[popmuted, thin] (\x, -\y) -- (\x*0.5, -\y*0.5);
}
\end{tikzpicture}
\legende{Engrenage conique : transmission entre deux axes concourants (ici perpendiculaires). Les primitives sont des cônes roulant sans glisser.}
\end{popfigure}

\begin{exemplebox}[La perceuse à main (manivelle)]
La perceuse à main (ou « vilebrequin ») est un exemple classique d'engrenage conique.
\begin{itemize}
\item La manivelle actionne un grand pignon conique (axe horizontal).
\item Celui-ci entraîne une petite roue conique (pignon) fixée sur l'axe du foret (axe vertical).
\item Le mouvement de rotation passe ainsi de l'horizontal à la verticale.
\item \textbf{Avantage :} Changement de direction du mouvement sans perte de puissance significative, encombrement réduit.
\end{itemize}
\end{exemplebox}

\begin{savaistu}[Au Cameroun : les motoculteurs et les engrenages coniques]
Dans les zones agricoles (Dschang, Foumbot, Ngaoundéré), les \textbf{motoculteurs} (tracteurs marche-pied) sont très utilisés pour le labour. Leur boîte de vitesse et leur prise de force utilisent des \textbf{engrenages coniques} pour transmettre la puissance du moteur (arbre horizontal) vers les roues (axes horizontaux) mais surtout vers l'outil porté à l'arrière (arbre vertical ou perpendiculaire). La robustesse des dents coniques (souvent en acier trempé) est essentielle pour supporter les chocs du sol caillouteux.
\end{savaistu}

\courssub{Pignon et crémaillère : rotation $\leftrightarrow$ translation}
C'est un cas particulier d'engrenage où l'une des « roues » a un rayon infini : c'est une \cle{crémaillère} (barre droite dentée). Le pignon (roue dentée standard) roule sur la crémaillère.

\begin{popfigure}
\begin{tikzpicture}[scale=1.0, >=Stealth]
% Crémaillère (en bas)
\def\mod{0.5} % module visuel
\def\pas{3.14159*\mod} % pi * m
\def\h{0.8} % hauteur dent
% Base
\draw[popdark, thick] (-4, -1) -- (4, -1);
% Dents crémaillère (profil trapézoïdal simplifié, largeur dent = largeur vide = pas/2)
\foreach \i in {-6,-5,...,6} {
  \pgfmathsetmacro{\x}{\i*\pas};
  \draw[poppurple, thick] (\x, -1) -- (\x, -1+\h) -- (\x+0.5*\pas, -1+\h) -- (\x+0.5*\pas, -1);
  \draw[poppurple, fill=poppurple!20] (\x, -1) -- (\x, -1+\h) -- (\x+0.5*\pas, -1+\h) -- (\x+0.5*\pas, -1) -- cycle;
}
\node[poppurple] at (0, -2) {Crémaillère (translation $v$)};

% Pignon (au-dessus)
\def\Rp{2.0}
\def\rp{1.5}
\def\Rpt{2.4}
\def\np{12}
\begin{scope}[yshift=2.5cm] % Distance centre pignon - pied crémaillère ~ Rp + jeu
\foreach \i in {1,...,\np} {
  \pgfmathsetmacro{\angle}{360/\np*\i - 90} % Dent en bas alignée
  \draw[popblue, thick] (\angle:\rp) -- (\angle:\Rpt);
  \draw[popblue, fill=popblue!20] (\angle:\Rpt) arc (\angle:\angle+360/\np*0.5:\Rpt) -- (\angle+360/\np*0.5:\rp) arc (\angle+360/\np*0.5:\angle:\rp) -- cycle;
}
\draw[popblue, thick] (0,0) circle (\Rp);
\draw[popdark, thick] (0,0) circle (0.2);
\draw[->, poporange, thick] (0,0) -- ++(90:1.5) node[above] {Axe $N$};
\draw[->, poporange, thick] (1.5,0) arc (0:90:1.5) node[right] {$\omega$ (horaire)};
\node[popblue] at (0, 3.5) {Pignon (rotation)};
\end{scope}

% Flèches transformation
\draw[->, popred, very thick] (0, 0.5) -- ++(0, -1.5) node[left, midway] {$\vec{v}_C$};
\draw[->, popred, very thick] (0, -1) -- ++(1.5, 0) node[below, midway] {Translation $v$};
\node[popmuted] at (0, 0.5) {$\bullet$} node[above=2pt, popmuted] {Contact C};
\end{tikzpicture}
\legende{Système pignon-crémaillère : transformation du mouvement de rotation (pignon) en mouvement de translation rectiligne (crémaillère), et inversement.}
\end{popfigure}

\begin{propriete}[Pignon - Crémaillère]
\begin{itemize}
\item Si le \textbf{pignon tourne} (vitesse angulaire $\omega$), la \textbf{crémaillère se déplace} en translation (vitesse $v = \omega \times R_p$).
\item Si la \textbf{crémaillère se déplace}, le \textbf{pignon tourne}.
\item Le sens de translation dépend du sens de rotation du pignon (horaire $\rightarrow$ translation vers la droite si pignon au-dessus).
\end{itemize}
\end{propriete}

\begin{exemplebox}[Direction de voiture à crémaillère]
C'est l'application la plus connue au quotidien.
\begin{enumerate}
\item Le conducteur tourne le \textbf{volant} (rotation).
\item L'arbre de colonne entraîne un \textbf{pignon} en bout d'arbre.
\item Le pignon engrène sur une \textbf{crémaillère} transversale.
\item La crémaillère se déplace latéralement (translation).
\item Des biellettes relient la crémaillère aux \textbf{roues directrices} : elles pivottent.
\end{enumerate}
C'est un système précis, réversible et peu encombrant.
\end{exemplebox}

\begin{exemplebox}[Perceuse à colonne (descente de la broche)]
Sur une perceuse à colonne d'atelier, la manivelle actionne un pignon qui engrène sur une crémaillère fixée sur le corps de la perceuse (la « colonne »). Tourner la manivelle fait monter ou descendre la broche verticalement. C'est une transformation rotation $\rightarrow$ translation verticale.
\end{exemplebox}

\courssub{Train d'engrenages : plusieurs roues en série}
Dans de nombreuses machines (boîte de vitesse de moto, horloge, mélangeuse), plusieurs engrenages sont montés en série : la roue menée du premier étage devient la roue motrice du suivant. On appelle cela un \cle{train d'engrenages}.

\begin{popfigure}
\begin{tikzpicture}[scale=0.9, >=Stealth]
% 3 roues alignées horizontalement
\def\R{1.5}
\def\r{1.0}
\def\Rt{1.9}
\def\n{10}
% Roue 1 (Moteur)
\begin{scope}[xshift=0cm]
\foreach \i in {1,...,\n} {
  \pgfmathsetmacro{\angle}{360/\n*\i}
  \draw[popblue, thick] (\angle:\r) -- (\angle:\Rt);
  \draw[popblue, fill=popblue!20] (\angle:\Rt) arc (\angle:\angle+360/\n*0.5:\Rt) -- (\angle+360/\n*0.5:\r) arc (\angle+360/\n*0.5:\angle:\r) -- cycle;
}
\draw[popblue, thick] (0,0) circle (\R);
\draw[popdark, thick] (0,0) circle (0.2);
\draw[->, poporange, thick] (0,0) -- ++(90:1.2) node[above] {$N_1$};
\draw[->, poporange, thick] (1.2,0) arc (0:90:1.2) node[right] {Horaire};
\node[popdark] at (0,-2.8) {R1 (Motrice)};
\end{scope}

% Roue 2 (Intermédiaire / Menée 1 / Motrice 2)
\begin{scope}[xshift=3.2cm] % distance 2R = 3.0 + jeu
\foreach \i in {1,...,\n} {
  \pgfmathsetmacro{\angle}{360/\n*\i + 180/\n}
  \draw[popgreen!70!black, thick] (\angle:\r) -- (\angle:\Rt);
  \draw[popgreen!70!black, fill=popgreen!20] (\angle:\Rt) arc (\angle:\angle+360/\n*0.5:\Rt) -- (\angle+360/\n*0.5:\r) arc (\angle+360/\n*0.5:\angle:\r) -- cycle;
}
\draw[popgreen!70!black, thick] (0,0) circle (\R);
\draw[popdark, thick] (0,0) circle (0.2);
\draw[->, poporange, thick] (0,0) -- ++(90:1.2) node[above] {$N_2$};
\draw[->, poporange, thick] (-1.2,0) arc (180:90:1.2) node[left] {Anti-horaire};
\node[popdark] at (0,-2.8) {R2 (Intermédiaire)};
\end{scope}

% Roue 3 (Menée finale)
\begin{scope}[xshift=6.4cm]
\foreach \i in {1,...,\n} {
  \pgfmathsetmacro{\angle}{360/\n*\i}
  \draw[poppurple, thick] (\angle:\r) -- (\angle:\Rt);
  \draw[poppurple, fill=poppurple!20] (\angle:\Rt) arc (\angle:\angle+360/\n*0.5:\Rt) -- (\angle+360/\n*0.5:\r) arc (\angle+360/\n*0.5:\angle:\r) -- cycle;
}
\draw[poppurple, thick] (0,0) circle (\R);
\draw[popdark, thick] (0,0) circle (0.2);
\draw[->, poporange, thick] (0,0) -- ++(90:1.2) node[above] {$N_3$};
\draw[->, poporange, thick] (1.2,0) arc (0:90:1.2) node[right] {Horaire};
\node[popdark] at (0,-2.8) {R3 (Menée finale)};
\end{scope}

% Flèches contact
\draw[->, popred, thick] (1.6,0) -- ++(0,0.8) node[above] {$\vec{v}_{12}$};
\draw[->, popred, thick] (1.6,0) -- ++(0,-0.8) node[below] {$\vec{v}_{21}$};
\draw[->, popred, thick] (4.8,0) -- ++(0,0.8) node[above] {$\vec{v}_{23}$};
\draw[->, popred, thick] (4.8,0) -- ++(0,-0.8) node[below] {$\vec{v}_{32}$};
\end{tikzpicture}
\legende{Train de 3 engrenages extérieurs. R1 horaire $\rightarrow$ R2 anti-horaire $\rightarrow$ R3 horaire. Nombre impair de roues (3) : R1 et R3 tournent dans le \textbf{même sens}.}
\end{popfigure}

\begin{methode}[Déterminer le sens de rotation final d'un train d'engrenages]
\textbf{Objectif :} Trouver le sens de la dernière roue connaissant le sens de la première.
\textbf{Principe :} Chaque engrenage extérieur inverse le sens. Chaque engrenage intérieur conserve le sens.
\textbf{Étapes :}
\begin{enumerate}
\item \textbf{Dessiner le schéma} simplifié (cercles pour les primitives, flèches pour les axes).
\item \textbf{Fixer le sens de la roue motrice} (donnée de l'énoncé : ex. horaire).
\item \textbf{Propager roue par roue} :
  \begin{itemize}
  \item Contact \textbf{extérieur} : flèche opposée (inversion).
  \item Contact \textbf{intérieur} (couronne) : flèche même sens (conservation).
  \item Pignon-crémaillère : pas de rotation pour la crémaillère (translation).
  \end{itemize}
\item \textbf{Compter les inversions} (raccourci pour train tout extérieur) :
  \begin{itemize}
  \item Si le train ne contient que des engrenages \textbf{extérieurs} :
    \begin{itemize}
    \item Nombre \textbf{pair} de roues $\rightarrow$ Sens \textbf{inverses} (1ère et dernière).
    \item Nombre \textbf{impair} de roues $\rightarrow$ Sens \textbf{identiques}.
    \end{itemize}
  \end{itemize}
\item \textbf{Conclure} en écrivant clairement : « La roue finale tourne dans le sens [horaire/anti-horaire]. »
\end{enumerate}
\end{methode}

\begin{exoresolu}[Sens de rotation dans une boîte de vitesse simplifiée]
\textbf{Énoncé :} Un train d'engrenages comporte 4 roues extérieures ($R_1$, $R_2$, $R_3$, $R_4$) montées en série sur des axes parallèles. $R_1$ est la roue motrice (entrée de boîte), $R_4$ est la roue menée (sortie de boîte). $R_1$ tourne dans le sens horaire. Déterminer le sens de rotation de $R_4$.
\textbf{Solution détaillée :}
\begin{enumerate}
\item \textbf{Analyse :} 4 roues $\rightarrow$ 3 contacts extérieurs ($R_1/R_2$, $R_2/R_3$, $R_3/R_4$). Tous sont des engrenages extérieurs.
\item \textbf{Propagation (Méthode pas à pas) :}
  \begin{itemize}
  \item $R_1$ : Horaire (donné).
  \item Contact $R_1/R_2$ (extérieur) $\rightarrow$ $R_2$ : \textbf{Anti-horaire}.
  \item Contact $R_2/R_3$ (extérieur) $\rightarrow$ $R_3$ : \textbf{Horaire}.
  \item Contact $R_3/R_4$ (extérieur) $\rightarrow$ $R_4$ : \textbf{Anti-horaire}.
  \end{itemize}
\item \textbf{Raccourci (Parité) :} Nombre de roues $N = 4$ (Pair). Pour un train tout extérieur : Pair $\rightarrow$ Sens inverses entre entrée et sortie.
\item \textbf{Conclusion :} La roue de sortie $R_4$ tourne dans le sens \textbf{anti-horaire}.
\end{enumerate}
\end{exoresolu}

\begin{attention}[Erreur classique : Compter les contacts au lieu des roues]
La règle de la parité s'applique sur le \textbf{nombre de roues} ($N$), pas le nombre de contacts ($N-1$).
\begin{itemize}
\item $N=2$ (1 contact) : Pair $\rightarrow$ Inverses. (Correct : 1 inversion).
\item $N=3$ (2 contacts) : Impair $\rightarrow$ Identiques. (Correct : 2 inversions = retour au sens initial).
\item $N=4$ (3 contacts) : Pair $\rightarrow$ Inverses. (Correct : 3 inversions = sens inverse).
\end{itemize}
\textbf{Astuce :} Dessine toujours les 3 premières flèches sur ton brouillon. C'est plus sûr que la règle de parité si le train mélange extérieurs et intérieurs !
\end{attention}

\courssub{Synthèse des types d'engrenages}

\begin{center}
\begin{tabularx}{\linewidth}{|l|X|X|X|}\hline
\rowcolor{popblueL}
\textbf{Type} & \textbf{Axes} & \textbf{Sens rotation (Entrée/Sortie)} & \textbf{Application courante (Cameroun / Vie courante)} \\ \hline
\textbf{Extérieur (cylindrique droit)} & Parallèles & \textbf{Inverses} & Boîte de vitesse moto/voiture, vélo, mélangeuse, groupe électrogène \\ \hline
\textbf{Intérieur (Pignon + Couronne)} & Parallèles & \textbf{Identiques} & Différentiel arrière voiture, réducteur de moteur de portillon automatique \\ \hline
\textbf{Conique} & Concourants (souvent $\perp$) & Inverses (selon montage) & Perceuse à main, transmission motoculteur, direction camion (vis sans fin + conique) \\ \hline
\textbf{Pignon - Crémaillère} & Perpendiculaires (Rotation $\perp$ Translation) & Rotation $\leftrightarrow$ Translation & Direction de voiture (crémaillère), perceuse à colonne, portique de levage \\ \hline
\textbf{Train (série extérieurs)} & Parallèles & Pair : Inverses \newline Impair : Identiques & Horlogerie, boîte de vitesse multi-rapports, transmission finale moto (pignon/couronne = 2 roues) \\ \hline
\end{tabularx}
\end{center}

\begin{exercice}[Application : Le système de levage d'un forage]
\textbf{Contexte :} Un forage villageois (pompe à main type India Mark II) utilise un système de levage. La manivelle (axe horizontal) entraîne un pignon conique qui engrène sur une grande roue conique (axe vertical). Cette roue conique est solidaire d'un tambour sur lequel s'enroule un câble (translation verticale de la tige de pompe).
\begin{enumerate}
\item Nommer les deux types d'engrenages (ou organes de transmission) présents dans cette chaîne de transmission.
\item Si on tourne la manivelle dans le sens horaire (vu de l'opérateur), et que le pignon conique est placé \emph{au-dessus} de la grande roue conique (attaque par le haut), dans quel sens tourne la grande roue conique (vue du dessus) ?
\item La grande roue conique étant solidaire du tambour, dans quel sens le câble s'enroule-t-il (montée ou descente de la tige) ?
\end{enumerate}
\end{exercice}

\begin{corrige}[Exercice : Système de levage d'un forage]
\begin{enumerate}
\item \textbf{Types :} \textbf{Engrenage conique} (manivelle $\rightarrow$ axe vertical) + \textbf{Tambour / Treuil} (organe de transformation rotation $\rightarrow$ translation). Note : le tambour n'est pas un engrenage à dents.
\item \textbf{Sens grande roue (Méthode des flèches au contact) :}
  \begin{itemize}
  \item Le pignon conique (axe horizontal) tourne \textbf{horaire} (vu de l'opérateur, donc vecteur rotation vers la gauche si manivelle à gauche).
  \item Au contact (sommet des cônes, pignon au-dessus), la vitesse tangentielle du pignon est dirigée \emph{vers l'arrière} (loin de l'opérateur).
  \item Pour rouler sans glisser, la grande roue conique (axe vertical) doit avoir au même point une vitesse tangentielle \emph{vers l'arrière}.
  \item Vu du dessus (axe vertical), une vitesse tangentielle vers l'arrière correspond à une rotation \textbf{horaire}.
  \end{itemize}
  \textbf{Réponse :} La grande roue conique tourne dans le sens \textbf{horaire} (vue du dessus).
\item \textbf{Sens câble :} Grande roue horaire (vue du dessus) $\rightarrow$ Tambour horaire. Si le câble sort par le haut du tambour vers la poulie de renvoi, un tambour horaire \textbf{enroule} le câble $\rightarrow$ \textbf{Montée} de la tige de pompe (et remontée de l'eau).
\end{enumerate}
\end{corrige}

\begin{aretenir}[Résumé : Types et sens de rotation]
\begin{itemize}
\item \textbf{Engrenage extérieur (cylindrique) :} Axes \textbf{parallèles}, sens \textbf{inverses}.
\item \textbf{Engrenage intérieur (couronne) :} Axes \textbf{parallèles}, sens \textbf{identiques}.
\item \textbf{Engrenage conique :} Axes \textbf{concourants} (souvent $\perp$), changement de direction du mouvement.
\item \textbf{Pignon - Crémaillère :} Transformation \textbf{Rotation $\leftrightarrow$ Translation}.
\item \textbf{Train d'engrenages (extérieurs) :} Sens final = Sens initial si \textbf{nombre impair} de roues ; Sens final = Sens inverse si \textbf{nombre pair} de roues.
\item \textbf{Méthode infaillible :} Dessiner les flèches roue par roue au contact.
\end{itemize}
\end{aretenir}

\coursec{Avantages et inconvénients des engrenages}

\courssub{Transition depuis les types d'engrenages}

Nous venons d'étudier les différents types d'engrenages (droit, hélicoïdal, conique, vis sans fin) et comment ils imposent le sens de rotation des arbres. Maintenant, posons-nous la question qu'un technicien se pose avant de choisir une solution technique : \emph{pourquoi utiliser des engrenages plutôt que des poulies-courroies, des chaînes ou des frottements ?} La réponse se trouve dans le bilan de leurs avantages et de leurs inconvénients, que nous allons établir à partir d'observations concrètes.

\courssub{Les quatre grands avantages des engrenages}

\begin{experience}[Observation : pas de glissement, effort important]
\textbf{Matériel :} un kit de construction type LEGO Technic ou Meccano (deux engrenages de même taille sur axes parallèles), une courroie élastique et deux poulies de même diamètre, un dynamomètre, des masses.
\textbf{Protocole :}
\begin{enumerate}
    \item Monter les deux engrenages en prise. Tourner l'axe moteur d'un tour complet. Observer l'axe récepteur.
    \item Remplacer par les poulies et la courroie. Tourner l'axe moteur d'un tour complet. Observer l'axe récepteur (surtout si on freine légèrement l'axe récepteur à la main).
    \item Avec le dynamomètre, mesurer le couple maximal transmissible avant patinage (courroie) ou blocage (engrenage).
\end{enumerate}
\textbf{Observations :}
\begin{itemize}
    \item Engrenages : l'axe récepteur fait \textbf{exactement} un tour. Pas de décalage. On peut transmettre un gros effort sans que les dents ne « sautent » (si bien dimensionnées).
    \item Courroie : l'axe récepteur fait \emph{presque} un tour, mais un léger glissement (patinage) apparaît dès qu'on freine. La courroie peut même sauter des crans si elle est dentée, ou glisser si elle est lisse.
    \item Couple : l'engrenage supporte un couple bien plus élevé avant rupture que la courroie avant patinage.
\end{itemize}
\textbf{Interprétation :} Le contact dent/dent assure un verrouillage géométrique (forme fermée), alors que le contact courroie/poulie repose sur le frottement (adhérence), donc une forme ouverte sujette au glissement.
\end{experience}

\begin{popfigure}
\begin{center}
\begin{tikzpicture}[scale=1.0, font=\small, >=Stealth]
% Engrenages
\begin{scope}[xshift=-4cm]
\draw[popblue, thick, fill=popblueL] (0,0) circle (1.5);
\draw[popblue, thick, fill=popblueL] (3,0) circle (1.5);
\foreach \a in {0,30,...,330} {
  \draw[popdark, thick] ({1.5*cos(\a)},{1.5*sin(\a)}) -- ({1.7*cos(\a)},{1.7*sin(\a)});
  \draw[popdark, thick] ({3+1.5*cos(\a)},{1.5*sin(\a)}) -- ({3+1.7*cos(\a)},{1.7*sin(\a)});
}
\draw[popdark, thick, ->] (0,-2) arc (-90:-270:1);
\node[popblue, align=center] at (1.5,-2.5) {\textbf{Engrenages}\\\textit{Prise de forme, pas de glissement}};
\draw[popdark, <->] (0,0) -- (3,0) node[midway, above] {Entraxe fixe};
\end{scope}

% Courroie
\begin{scope}[xshift=4cm]
\draw[poporange, thick, fill=poporangeL] (0,0) circle (1.5);
\draw[poporange, thick, fill=poporangeL] (3,0) circle (1.5);
\draw[poporange, thick] (0,1.5) -- (3,1.5);
\draw[poporange, thick] (0,-1.5) -- (3,-1.5);
\draw[popdark, dashed, ->] (0,-2) arc (-90:-260:1);
\node[poporange, align=center] at (1.5,-2.5) {\textbf{Courroie}\\\textit{Adhérence, glissement possible}};
\draw[popdark, <->] (0,0) -- (3,0) node[midway, above] {Entraxe réglable};
\end{scope}
\end{tikzpicture}
\end{center}
\legende{Comparaison schématique : engrenages (prise de forme rigide, entraxe imposé) vs courroie (adhérence, entraxe tolérant, risque de patinage).}
\end{popfigure}

\begin{propriete}[Transmission sans glissement — Savoir admis au BEPC]
L'engrenage transmet le mouvement de rotation \textbf{sans glissement} (pas de patinage) entre l'arbre moteur et l'arbre récepteur. Le rapport de transmission est donc \textbf{constant et exact}, égal au rapport des nombres de dents (ou des diamètres primitifs), contrairement à la transmission par courroie où le glissement fait varier légèrement le rapport.
\end{propriete}

\textbf{Avantage 1 : Précision et absence de patinage.} Comme le montre l'expérience, les dents s'engrènent comme des « doigts » qui se serrent. Il n'y a pas de glissement relatif au point de contact (roulement pur). C'est indispensable pour les \cle{montre}s, les \cle{compteurs kilométriques} de moto-taxi, les \cle{robots industriels} ou la \cle{direction} d'une voiture : un tour d'entrée donne \emph{toujours} le même angle de sortie.

\textbf{Avantage 2 : Transmission de grands efforts (couples élevés).} Les dents répartissent l'effort sur une surface de contact en acier trempé. Une boîte de vitesses de voiture ou le réducteur d'un \cle{groupe électrogène} (marque Elemax ou Kipor, très courant au Cameroun) transmet des couples de plusieurs centaines de newton-mètres dans un volume réduit. Une courroie aurait besoin de largeurs énormes ou patinerait.

\textbf{Avantage 3 : Rapport de transmission constant et exact.} Puisque $Z_1 \cdot N_1 = Z_2 \cdot N_2$ (ou $D_{p1} \cdot N_1 = D_{p2} \cdot N_2$) sans terme de glissement, le rapport $k = \frac{N_1}{N_2} = \frac{Z_2}{Z_1}$ est \textbf{figé par la géométrie}. C'est la base des \cle{boîtes de vitesses} : on change de rapport en changeant les paires d'engrenages en prise, pas en modifiant un glissement.

\textbf{Avantage 4 : Bon rendement et grande durabilité (si lubrifié).} Le rendement d'une paire d'engrenages droits bien lubrifiés dépasse \textbf{98~\%} (souvent 99~\% par étage). Sur une moto-taxi (type Honda CG 125, Yamaha DT, ou les motos chinoises type Haojue), la boîte de vitesses baigne dans l'huile moteur : les engrenages peuvent faire \textbf{100~000~km} sans usure significative. C'est bien plus durable qu'une courroie de distribution (à changer tous les 60~000~km ou 5 ans) ou qu'une chaîne de distribution (qui s'allonge).

\begin{exemplebox}[La boîte de vitesses d'une moto-taxi à Douala]
Ouvre le carter moteur d'une moto-taxi : tu vois un \textbf{bain d'huile} où trempent les arbres primaire et secondaire avec leurs pignons. Pourquoi ?
\begin{itemize}
    \item L'huile assure la \textbf{lubrification} (film d'huile entre dents, pas de contact métal-métal direct).
    \item Elle assure le \textbf{refroidissement} (évacue la chaleur due aux frottements).
    \item Elle assure le \textbf{nettoyage} (emporte les microparticules d'usure vers le filtre ou le fond du carter).
\end{itemize}
C'est grâce à cet entretien simple (vidange tous les 3~000~km ou 6 mois) que les engrenages tiennent la distance.
\end{exemplebox}

\courssub{Les quatre inconvénients majeurs}

\begin{attention}[Piège fréquent : croire que l'engrenage est « increvable »]
Un engrenage \textbf{mal lubrifié} s'use très vite. Le frottement métal-métal sous haute pression (contact de Hertz) crée une usure par \cle{pitting} (picotage), puis par \cle{écaillage} des dents, jusqu'à la \cle{rupture} brutale. Un bruit de « grincement » ou de « ronronnement » anormal est un signal d'alerte : \emph{arrête la machine et vérifie le niveau d'huile !}
\end{attention}

\textbf{Inconvénient 1 : Bruit de fonctionnement.} Le choc successif des dents (surtout en denture droite) génère des vibrations et du bruit aérien. À grande vitesse ($> 3000$~tr/min), le « sifflement » ou « ronronnement » devient gênant. C'est pour cela que :
\begin{itemize}
    \item Les \textbf{boîtes de vitesses automobiles} utilisent des dentures \textbf{hélicoïdales} (contact progressif, plus silencieux).
    \item La \textbf{marche arrière} reste souvent à denture droite (coût moindre, usage court) : c'est pour cela qu'elle fait ce \textbf{bruit caractéristique « whiiiir »} quand tu recules ta voiture ou ta moto (si elle a une marche arrière, rare sur les motos).
\end{itemize}

\begin{savaistu}[Le bruit de la marche arrière]
Sur la plupart des boîtes de vitesses manuelles (voitures, camions, tracteurs), l'engrenage de marche arrière est à \textbf{denture droite} (coût de fabrication plus simple, axe intermédiaire). Les dentures droites sont bruyantes car toute la dent entre en contact d'un coup. Les rapports avant (1ère, 2ème, 3ème, 4ème, 5ème) sont en \textbf{denture hélicoïdale} : le contact commence par une extrémité de la dent et progresse, ce qui amortit le choc et réduit le bruit. Donc, le bruit strident en marche arrière est \textbf{normal} : c'est la signature acoustique de la denture droite !
\end{savaistu}

\textbf{Inconvénient 2 : Lubrification obligatoire et entretien régulier.} Contrairement à une courroie (sèche, pas d'entretien) ou une chaîne (graissage externe périodique), l'engrenage \textbf{doit} baigner dans l'huile (bain d'huile, projection, circulation forcée sous pression). Sans huile, la durée de vie chute de \textbf{100~000~h} à \textbf{quelques heures}. Au Cameroun, sur les \textbf{groupes électrogènes} des quartiers (quartier Makepe, Bonabéri, etc.), la première cause de casse moteur est l'oubli de la vidange : l'huile noircie, chargée en limaille, devient abrasive et use les pignons de distribution.

\textbf{Inconvénient 3 : Coût de fabrication et précision de montage.} Tailler des dents d'engrenage (fraisage, taillage par génératrice, rectification) demande des machines-outils coûteuses (fraiseuses à engrenages, machines à tailler) et un contrôle métrologique strict (pas, profil, pas de base). De plus, l'\textbf{entraxe} (distance entre axes) doit être respecté au \textbf{centième de millimètre près} : trop serré $\rightarrow$ grippage, trop large $\rightarrow$ jeu excessif, bruit, choc, usure prématurée. Une courroie tolère quelques millimètres d'erreur sur l'entraxe ; un engrenage, non.

\textbf{Inconvénient 4 : Usure des dents et risque de rupture.} Même bien lubrifié, l'engrenage s'use lentement (usure adhesive, fatigue de contact). Le profil de dent s'affine, le jeu de fonctionnement augmente. En cas de \textbf{surcharge} (démarrage brutal, blocage récepteur), la dent peut casser net (rupture fragile à la base de la dent). Une dent cassée tombe dans le carter, circule dans l'huile, et peut détruire tout le train d'engrenages (effet « grenade »).

\begin{popfigure}
\begin{center}
\begin{tikzpicture}[scale=1.2, font=\small, >=Stealth]
% Dent neuve
\begin{scope}[xshift=-4.5cm]
\draw[popdark, thick, fill=popgreenL] (0,0) -- (0.8,0) -- (0.8,1.5) -- (0.6,1.5) -- (0.6,0.3) 
    .. controls (0.5,0.5) and (0.3,0.5) .. (0.2,0.3) -- (0,0.3) -- cycle;
\draw[popdark, thick] (0.8,0) -- (1.6,0) -- (1.6,1.5) -- (1.4,1.5) -- (1.4,0.3)
    .. controls (1.3,0.5) and (1.1,0.5) .. (1.0,0.3) -- (0.8,0.3) -- cycle;
\draw[popdark, thick] (1.6,0) -- (2.4,0) -- (2.4,1.5) -- (2.2,1.5) -- (2.2,0.3)
    .. controls (2.1,0.5) and (1.9,0.5) .. (1.8,0.3) -- (1.6,0.3) -- cycle;
\node[popgreen, align=center] at (1.2,-0.7) {\textbf{Dent neuve}\\\textit{Profil théorique parfait}};
\draw[popgreen, <->] (0.8,1.6) -- (1.6,1.6) node[midway, above] {$p$ (pas)};
\end{scope}

% Dent usée
\begin{scope}[xshift=4.5cm]
\draw[popdark, thick, fill=poporangeL] (0,0) -- (0.8,0) -- (0.75,1.5) -- (0.55,1.5) -- (0.55,0.35) 
    .. controls (0.45,0.55) and (0.25,0.55) .. (0.15,0.35) -- (0,0.35) -- cycle;
\draw[popdark, thick, dashed] (0.8,0) -- (0.8,1.5); % profil théorique
\draw[popdark, thick, fill=poporangeL] (0.85,0) -- (1.6,0) -- (1.55,1.5) -- (1.35,1.5) -- (1.35,0.35)
    .. controls (1.25,0.55) and (1.05,0.55) .. (0.95,0.35) -- (0.85,0.35) -- cycle;
\draw[popdark, thick, dashed] (1.6,0) -- (1.6,1.5);
\draw[popdark, thick, fill=poporangeL] (1.65,0) -- (2.4,0) -- (2.35,1.5) -- (2.15,1.5) -- (2.15,0.35)
    .. controls (2.05,0.55) and (1.85,0.55) .. (1.75,0.35) -- (1.65,0.35) -- cycle;
\draw[popdark, thick, dashed] (2.4,0) -- (2.4,1.5);
\node[poporange, align=center] at (1.2,-0.7) {\textbf{Dent usée}\\\textit{Affaiblissement, jeu accru}};
\draw[<->, poporange] (0.55,1.6) -- (1.35,1.6) node[midway, above] {$p' > p$};
\draw[poporange, <->] (0.75,1.5) -- (0.8,1.5) node[midway, above, xshift=0.1cm] {usure};
\end{scope}
\end{tikzpicture}
\end{center}
\legende{Comparaison profil dent neuve (vert, plein) vs dent usée (orange, trait plein) et profil théorique (trait pointillé). L'usure augmente le pas effectif et réduit l'épaisseur de dent (risque de rupture).}
\end{popfigure}

\courssub{Choix technologique : Engrenages vs Poulies-Courroies}

Le tableau comparatif ci-dessous synthétise les critères de choix pour un concepteur. Rappelle-toi la leçon sur les poulies-courroies (Séance 13) : la courroie est souple, l'engrenage est rigide.

\begin{popfigure}
\begin{center}
\begin{tikzpicture}[font=\small, scale=0.9, >=Stealth]
% Styles
\tikzset{
    header/.style={fill=popblue, text=white, font=\bfseries, align=center, minimum height=1.2em},
    crit/.style={fill=popblueL, font=\bfseries, align=left, minimum height=1.5em},
    cell/.style={align=center, minimum height=1.5em},
    adv/.style={fill=popgreenL, text=popdark},
    disadv/.style={fill=poporangeL, text=popdark},
    neutral/.style={fill=popcream, text=popdark},
}
% Tableau
\matrix (tab) [matrix of nodes, nodes in empty cells, nodes={draw=popline, inner sep=3pt}, column sep=-\pgflinewidth, row sep=-\pgflinewidth,
    column 1/.style={nodes={crit, minimum width=3.2cm}},
    column 2/.style={nodes={cell, minimum width=3.5cm}},
    column 3/.style={nodes={cell, minimum width=3.5cm}},
]
{
|[header]| Critère & |[header]| \textbf{Engrenages} & |[header]| \textbf{Poulies-Courroies} \\
Précision / Rapport constant & |[adv]| Excellent (exact, sans glissement) & |[disadv]| Moyen (glissement 1 à 3\%, élongation) \\
Couple transmissible / Compacité & |[adv]| Très fort (acier, contact surfacique) & |[neutral]| Modéré (nécessite grandes largeurs) \\
Bruit & |[disadv]| Élevé (denture droite) à Modéré (hélicoïdale) & |[adv]| Faible (souplesse, amortissement) \\
Coût fabrication & |[disadv]| Élevé (usinage précision, trempe) & |[adv]| Faible (tôlerie, moulage, caoutchouc) \\
Entraxe & |[disadv]| Fixe, précis (centième mm) & |[adv]| Réglable, tolérant (tendeur) \\
Entretien / Lubrification & |[disadv]| Obligatoire (huile, vidange, filtre) & |[adv]| Quasi nul (contrôle tension, usure) \\
Rendement & |[adv]| Excellent (98-99\% / étage) & |[neutral]| Bon (95-98\% selon tension) \\
Durée de vie & |[adv]| Très longue (si entretenu) & |[neutral]| Limitée (vieillissement caoutchouc, usure) \\
Amortissement chocs/vibrations & |[disadv]| Nul (rigide) $\rightarrow$ couplage élastique requis & |[adv]| Naturel (souplesse courroie) \\
Sécurité (rupture) & |[disadv]| Risque projection éclats (carter obligatoire) & |[adv]| Rupture courroie = arrêt simple \\
};
% Légende couleurs
\begin{scope}[yshift=-14cm, xshift=-5cm]
\node[adv, minimum width=2cm] (l1) {Avantage};
\node[disadv, minimum width=2cm, right=0.2cm of l1] (l2) {Inconvénient};
\node[neutral, minimum width=2cm, right=0.2cm of l2] (l3) {Neutre / Moyen};
\end{scope}
\end{tikzpicture}
\end{center}
\legende{Tableau comparatif pour guider le choix technologique : Engrenages (rigides, précis, puissants, bruyants, chers) vs Poulies-Courroies (souples, tolérants, silencieux, économiques, glissement).}
\end{popfigure}

\begin{methode}[Comment choisir entre engrenage et courroie ?]
\begin{enumerate}
    \item \textbf{Précision exigée ?} (horlogerie, robotique, boîte de vitesses, compteur) $\rightarrow$ \textbf{Engrenage}.
    \item \textbf{Gros couple, espace réduit ?} (réducteur moto, voiture, groupe électrogène, broyeur) $\rightarrow$ \textbf{Engrenage}.
    \item \textbf{Arbres éloignés, entraxe variable, besoin de silence, budget serré ?} (ventilateur, compresseur frigo, alternateur voiture, convoyeur) $\rightarrow$ \textbf{Courroie} (ou chaîne si couple moyen + précision).
    \item \textbf{Amortir les à-coups ?} (démarrage moteur thermique sur machine-outil) $\rightarrow$ \textbf{Courroie} ou \textbf{accouplement élastique} + engrenage.
    \item \textbf{Environnement sale/poussière/extérieur ?} $\rightarrow$ \textbf{Engrenage en carter étanche huilé} (mieux que courroie qui glisse si huile/poussière).
\end{enumerate}
\end{methode}

\courssub{Exercice résolu : Choix d'une transmission pour une pompe à eau}

\begin{exoresolu}[Choix technologique pour une pompe de forage]
\textbf{Énoncé :} Un forage à Ngaoundéré doit pomper de l'eau à 40~m de profondeur. Le moteur thermique (monocylindre 4 temps, 7~ch, 3600~tr/min) entraîne une pompe volumétrique qui doit tourner à 600~tr/min. L'entraxe moteur-pompe est de 300~mm. L'installation est en plein air, poussiéreuse. Proposer la transmission la plus adaptée (engrenages ou courroie) et justifier.
\tcblower
\textbf{Solution détaillée :}
\begin{enumerate}
    \item \textbf{Rapport de transmission :} $k = \frac{N_{\text{moteur}}}{N_{\text{pompe}}} = \frac{3600}{600} = 6$. Rapport important.
    \item \textbf{Contraintes :}
    \begin{itemize}
        \item Moteur thermique monocylindre $\rightarrow$ \textbf{fortes vibrations, à-coups cycliques}.
        \item Extérieur, poussière $\rightarrow$ risque d'abrasion sur courroie, grippage sur engrenage nu.
        \item Entraxe 300~mm $\rightarrow$ compatible avec les deux.
        \item Couple : 7~ch $\approx$ 5,15~kW. Couple moteur $C_m = \frac{P}{\omega} = \frac{5150}{3600 \times 2\pi/60} \approx 13,7$~N$\cdot$m. Couple pompe $C_p = 6 \times 13,7 \approx 82$~N$\cdot$m. Couple modéré.
    \end{itemize}
    \item \textbf{Analyse :}
    \begin{itemize}
        \item \emph{Engrenages :} Précis, compacts pour ce couple. Mais \textbf{très sensibles aux à-coups} (risque cassure dents) et à la poussière (nécessite carter étanche coûteux, lubrification). Bruit élevé.
        \item \emph{Courroie (type V ou dentée) :} \textbf{Amortit naturellement les à-coups} du monocylindre (souplesse). Tolère la poussière (courroie lisse moins sensible que dentée). Silencieuse. Entraxe 300~mm facile avec tendeur. Coût faible. Glissement négligeable pour une pompe (pas de précision angulaire requise).
    \end{itemize}
    \item \textbf{Choix :} \textbf{Transmission par courroie trapézoïdale (type V) avec tendeur automatique.} C'est la solution standard sur les motopompes et groupes électrogènes portatifs au Cameroun (ex. pompes Robin, Honda, ou marques chinoises type Lifan).
    \item \textbf{Remarque :} Si on exigeait un rapport \emph{exact} (ex. alternateur synchronisé sur réseau), on choisirait une \textbf{courroie dentée} (synchrone) ou un \textbf{engrenage en carter huilé}.
\end{enumerate}
\end{exoresolu}

\courssub{Synthèse et fiche de révision}

\begin{aretenir}[Bilan : Engrenages — Avantages / Inconvénients]
\textbf{Avantages (Pourquoi on les choisit) :}
\begin{itemize}
    \item Transmission \textbf{rigide} : pas de glissement, rapport \textbf{exact} et constant ($k = Z_2/Z_1$).
    \item Transmettent de \textbf{grands couples} dans un volume réduit (acier trempé, contact surfacique).
    \item \textbf{Rendement excellent} ($> 98~\%$ par étage) et \textbf{durée de vie très longue} (100~000~km moto, 20~ans industrie) \textbf{SI lubrification parfaite}.
    \item Compacité axiale (arbres parallèles rapprochés).
\end{itemize}
\textbf{Inconvénients (Contraintes à gérer) :}
\begin{itemize}
    \item \textbf{Bruit} (choc dents), surtout denture droite $\rightarrow$ denture hélicoïdale pour réduire.
    \item \textbf{Lubrification impérative} (bain d'huile, projection, circulation) + vidanges régulières. Sans huile = casse rapide.
    \item \textbf{Coût élevé} (usinage précision, traitement thermique, métrologie).
    \item \textbf{Montage précis} : entraxe au 1/100~mm, alignement axes rigoureux.
    \item \textbf{Usure / Rupture} : pitting, écaillage, rupture dent si surcharge ou défaut huile. Carter de protection obligatoire (éclats).
    \item \textbf{Pas d'amortissement} : transmet les chocs $\rightarrow$ prévoir accouplement élastique en amont si moteur thermique.
\end{itemize}
\textbf{Règle de choix :} \emph{Précision + Couple fort + Compacité + Entretien possible} $\rightarrow$ \textbf{Engrenage}. \emph{Distance + Silence + À-coups + Budget + Tolérances} $\rightarrow$ \textbf{Courroie / Chaîne}.
\end{aretenir}

\begin{exercice}[Application : Diagnostic d'une boîte de vitesses]
\begin{enumerate}
    \item Un chauffeur de taxi (Toyota Corolla) entend un sifflement aigu en 3ème et 4ème vitesse, qui disparaît au point mort. L'huile de boîte a été vidangée il y a 3 ans (préconisation : 60~000~km ou 2 ans). Quel est le diagnostic probable ? Quelle est la cause racine ?
    \item Sur une moto-taxi (Honda CG 125), le kick démarreur ne « prend » plus par moments : la pédale descend dans le vide sans tourner le moteur. On entend un « clac » métallique. Quel organe d'engrenage est en cause ? Pourquoi ?
    \item Un agriculteur à Dschang veut motoriser une décortiqueuse de maïs avec un moteur diesel 10~ch (1500~tr/min). L'arbre de la machine doit tourner à 500~tr/min. Entraxe 400~mm. Environnement très poussiéreux (poussière de maïs). Proposer la transmission (engrenages ou courroie) et justifier par 3 arguments tirés du cours.
\end{enumerate}
\end{exercice}

\begin{corrige}[Exercice 1]
\begin{enumerate}
    \item \textbf{Diagnostic :} Usure prématurée des pignons de 3ème/4ème (souvent sur le même arbre secondaire) ou usure des roulements de cet arbre. \textbf{Cause racine :} Huile trop vieille, chargée en limaille, perte de pouvoir lubrifiant (viscosité, additifs épuisés) $\rightarrow$ contact métal-métal $\rightarrow$ pitting puis écaillage des dents $\rightarrow$ bruit de sifflement (denture hélicoïdale usée). \textbf{Action :} Vidange immédiate, contrôle limaille dans l'huile (aimant bouchon), si bruit persiste : dépose boîte, remplacement pignons/roulements usés.
    \item \textbf{Organe :} Le \textbf{roue libre du kick} (ou cliquet de kick). C'est un engrenage à cliquets (dents en rampe + ressort) qui doit s'enclencher dans le sens « kick » et tourner libre dans l'autre sens. \textbf{Cause :} Usure des rampes (dents arrondies) ou ressort cassé/faible $\rightarrow$ le cliquet ne « mord » plus sur la came. Graisse sèche ou absence de graisse dans le carter kick (souvent oublié à la vidange boîte). \textbf{Réparation :} Démontage carter kick, nettoyage, remplacement jeu cliquet/ressort, regraissage.
    \item \textbf{Choix :} \textbf{Courroie trapézoïdale (type V) ou courroie dentée large.} \textbf{Justifications :}
    \begin{enumerate}
        \item \textbf{Amortissement des à-coups :} Moteur diesel monocylindre/bicylindre = vibrations cycliques violentes. La courroie protège la décortiqueuse (arbre, pales) et le moteur (vilebrequin) des chocs de transmission (inconvénient engrenage : rigidité).
        \item \textbf{Environnement poussiéreux :} Poussière de maïs abrasive. Une courroie lisse (type V) tolère la poussière sur les poulies (nettoyage centrifuge). Un engrenage nu s'use en heures ; un engrenage en carter huilé demande un carter étanche complexe et cher (joints, respirateur, filtre huile).
        \item \textbf{Coût et maintenance :} Courroie + 2 poulies = coût faible, montage tolérant (entraxe 400~mm facile), changement courroie en 10~min sans outil spécial. Engrenages = usinage spécial, carter, huile, alignement précis, coût 5 à 10 fois supérieur.
    \end{enumerate}
\end{enumerate}
\end{corrige}

\coursec{Rapport de transmission}

\courssub{Transition : de la forme à la mesure}

Dans la section précédente, nous avons découvert que les engrenages permettent de transmettre un mouvement de rotation entre deux axes, de changer le sens de rotation et d'adapter la vitesse ou la force. Mais comment \emph{quantifier} cette adaptation ? Un mécanicien qui répare une moto-taxi à Douala, un cycliste qui gravit la côte de Nkolbisson ou un ingénieur qui règle le groupe électrogène de l'hôpital de district ont tous besoin d'un \textbf{nombre} précis pour choisir les bonnes roues dentées. Ce nombre, c'est le \textbf{rapport de transmission}. Nous allons le construire pas à pas, en partant d'un constat expérimental simple.

\courssub{Constat expérimental : l'égalité des dents au contact}

\begin{experience}[Le nombre de dents qui s'engrènent est le même]
\textbf{Matériel :} deux roues dentées de tailles différentes (ex. $Z_1 = 12$ dents et $Z_2 = 36$ dents) montées sur un support, un feutre.
\textbf{Protocole :}
\begin{enumerate}
    \item Marque d'un trait de feutre une dent de la petite roue (menante) et la dent correspondante de la grande roue (menée) au point de contact.
    \item Fais tourner lentement la roue menante d'un tour complet ($N_1 = 1$ tour).
    \item Compte combien de dents de la roue menée ont passé au point de contact.
    \item Recommence pour plusieurs tours de la menante.
\end{enumerate}
\textbf{Observations :}
\begin{itemize}
    \item Quand la petite roue fait \textbf{1 tour}, ses $12$ dents passent toutes au contact. La grande roue avance de $12$ dents aussi.
    \item Comme la grande roue a $36$ dents, elle n'a fait que $\frac{12}{36} = \frac{1}{3}$ de tour.
    \item Si la menante fait $3$ tours ($3 \times 12 = 36$ dents), la menée fait exactement $1$ tour ($36$ dents).
\end{itemize}
\textbf{Interprétation :} Au point de contact, \textbf{chaque dent de la roue menante pousse une dent de la roue menée}. Le nombre de dents qui défilent par unité de temps est \emph{identique} pour les deux roues. C'est la clé de tout le raisonnement.
\legende{Expérience : comptage des dents au point de contact. La petite roue (menante) entraîne la grande roue (menée).}
\end{experience}

\courssub{Établissement de la relation fondamentale}

Notons :
\begin{itemize}
    \item $Z_1$ : nombre de dents de la \textbf{roue menante} (celle qui reçoit le mouvement, ex. le pignon du moteur).
    \item $Z_2$ : nombre de dents de la \textbf{roue menée} (celle qui reçoit le mouvement transmis, ex. la roue arrière).
    \item $N_1$ : vitesse de rotation de la menante (en tours par minute, \si{tr/min} ou \si{min^{-1}}).
    \item $N_2$ : vitesse de rotation de la menée (en \si{tr/min}).
\end{itemize}

En $1$ minute, la menante fait $N_1$ tours. Le nombre total de dents qui passent au contact est $Z_1 \times N_1$.
De même, pour la menée : $Z_2 \times N_2$ dents passent au contact.
D'après notre constat expérimental, ces deux quantités sont égales :

\begin{propriete}[Relation fondamentale des engrenages]
Dans un engrenage, le produit \og nombre de dents $\times$ vitesse de rotation \fg{} se conserve :
\[
\boxed{Z_1 \times N_1 = Z_2 \times N_2}
\]
Cette égalité vaut pour n'importe quel couple de roues en prise directe (sans chaîne ni courroie).
\end{propriete}

\courssub{Définition du rapport de transmission}

À partir de la relation fondamentale, on isole le quotient des vitesses ou des nombres de dents.

\begin{definition}[Rapport de transmission $r$]
Le \textbf{rapport de transmission} $r$ est le quotient du nombre de dents de la roue menée par le nombre de dents de la roue menante. Il est aussi égal au quotient de la vitesse de la menante par la vitesse de la menée :
\[
\boxed{r = \frac{Z_2}{Z_1} = \frac{N_1}{N_2}}
\]
$r$ est un nombre \textbf{sans unité} (c'est un rapport de deux grandeurs de même nature).
\end{definition}

\begin{attention}[Piège classique : l'ordre des indices !]
\begin{itemize}
    \item Au numérateur : toujours la \textbf{menée} ($Z_2$, $N_2$ pour les dents ; $N_1$ pour la vitesse).
    \item Au dénominateur : toujours la \textbf{menante} ($Z_1$, $N_1$ pour les dents ; $N_2$ pour la vitesse).
    \item Astuce mnémotechnique : \og \textbf{M}enante commence par \textbf{M} comme \textbf{M}oins (dénominateur pour les dents, numérateur pour la vitesse) \fg. Ou plus simplement : \og \textbf{Menée} sur \textbf{Menante} \fg{} pour les dents ($Z_2/Z_1$).
    \item \textbf{Vérification systématique :} Si $r > 1$ (réducteur), alors $N_2 < N_1$. Si $r < 1$ (multiplicateur), alors $N_2 > N_1$. Si ton résultat contredit cela, tu as inversé $Z_1$ et $Z_2$.
\end{itemize}
\end{attention}

\courssub{Cas réducteur et cas multiplicateur}

La valeur de $r$ par rapport à $1$ dicte le comportement du mécanisme.

\begin{propriete}[Réducteur et multiplicateur]
\begin{itemize}
    \item \textbf{Cas réducteur ($r > 1$, c'est-à-dire $Z_2 > Z_1$) :} La roue menée est plus grande. Elle tourne \textbf{moins vite} ($N_2 < N_1$) mais elle fournit \textbf{plus de force} (couple plus grand). C'est le cas du démarrage d'une moto, de la montée en vélo (petit plateau / grand pignon) ou du treuil de levage.
    \item \textbf{Cas multiplicateur ($r < 1$, c'est-à-dire $Z_2 < Z_1$) :} La roue menée est plus petite. Elle tourne \textbf{plus vite} ($N_2 > N_1$) mais avec \textbf{moins de force}. C'est le cas de la grande vitesse sur route plate (grand plateau / petit pignon) ou du ventilateur entraîné par un moteur lent.
    \item \textbf{Cas direct ($r = 1$, $Z_2 = Z_1$) :} Vitesses égales, sens inverse. Utilisé pour transmettre le mouvement à distance sans changer le régime.
\end{itemize}
\end{propriete}

\begin{popfigure}
\begin{tikzpicture}[scale=0.9, every node/.style={font=\small}]
    % Styles
    \tikzset{
        gear/.style={draw=popdark, thick, fill=popblueL},
        shaft/.style={draw=popdark, fill=popmuted},
        labelbox/.style={draw=popblue, fill=popblueL, rounded corners=3pt, inner sep=3pt, font=\small}
    }
    % Menante (petit pignon) - Z1=20
    \begin{scope}[rotate=0]
        \foreach \i in {1,...,20} {
            \draw[gear] ({18*\i}:1.2) -- ({18*\i}:1.5) -- ({18*\i+9}:1.5) -- ({18*\i+9}:1.2) -- cycle;
        }
        \fill[shaft] (0,0) circle (0.4);
        \node[labelbox, align=center] at (0,-2.2){Roue menante (moteur)\\\textbf{Z$_1$ = 20 dents}\\\textbf{N$_1$ = 900 \si{tr/min}}};
        \draw[->, thick, poporange] (0,0) -- (-1.5,0) node[midway, above] {Entrée};
    \end{scope}
    
    % Contact point
    \draw[popline, dashed] (3.2,0) -- (3.2,-2.5);
    \node[popline, font=\tiny] at (3.2,-2.7) {Point de contact};
    
    % Menée (grande roue) - Z2=60
    \begin{scope}[xshift=7cm]
        \foreach \i in {1,...,60} {
            \draw[gear] ({6*\i}:2.8) -- ({6*\i}:3.2) -- ({6*\i+3}:3.2) -- ({6*\i+3}:2.8) -- cycle;
        }
        \fill[shaft] (0,0) circle (0.6);
        \node[labelbox, align=center] at (0,-4.0){Roue menée (sortie)\\\textbf{Z$_2$ = 60 dents}\\\textbf{N$_2$ = 300 \si{tr/min}}};
        \draw[->, thick, popgreen] (0,0) -- (1.5,0) node[midway, above] {Sortie};
    \end{scope}
    
    % Arrow between centers
    \draw[<->, thick, popdark] (1.8,0) -- (5.2,0) node[midway, above, font=\small] {Axe d'entraînement};
    
    % Ratio box
    \node[draw=popgold, fill=popgoldL, rounded corners=5pt, inner sep=5pt, align=center, font=\small\bfseries] at (3.5, 4.2) {
        Calcul du rapport : $r = \frac{Z_2}{Z_1} = \frac{60}{20} = 3$ \\
        $r > 1 \Rightarrow$ \textcolor{popgreen}{\textbf{Réducteur}} \\
        $N_2 = \frac{N_1}{r} = \frac{900}{3} = 300\ \si{tr/min}$
    };
\end{tikzpicture}
\legende{Engrenage réducteur : petit pignon menant ($Z_1=20$) entraînant une grande roue menée ($Z_2=60$). La vitesse est divisée par 3, le couple est multiplié par 3 (idéal pour démarrer une charge lourde).}
\end{popfigure}

\courssub{Exemple chiffré guidé : application directe}

\begin{exemplebox}[Calcul de la vitesse de sortie d'un réducteur]
Un moteur électrique de pompe à eau (type groupe électrogène ENEO de secours) tourne à \textbf{$N_1 = 1\,500\ \si{tr/min}$}. Son arbre porte un pignon de \textbf{$Z_1 = 25$ dents}. Ce pignon engrène directement sur une roue de pompe de \textbf{$Z_2 = 75$ dents}.
\textbf{Question :} Quelle est la vitesse de rotation $N_2$ de la roue de pompe ? Quel type d'engrenage est-ce ?
\textbf{Résolution :}
\begin{enumerate}
    \item \textbf{Identifier menante et menée :} Le moteur entraîne $\rightarrow$ Menante ($Z_1=25$, $N_1=1500$). La pompe est entraînée $\rightarrow$ Menée ($Z_2=75$, $N_2=?$).
    \item \textbf{Calculer le rapport $r$ :} $r = \frac{Z_2}{Z_1} = \frac{75}{25} = 3$.
    \item \textbf{Interpréter $r$ :} $r = 3 > 1$, c'est un \textbf{réducteur}. La pompe tournera moins vite que le moteur. Cohérent : une pompe a besoin de force (couple), pas de vitesse excessive.
    \item \textbf{Calculer $N_2$ :} On utilise $r = \frac{N_1}{N_2} \Rightarrow N_2 = \frac{N_1}{r} = \frac{1500}{3} = 500\ \si{tr/min}$.
    \item \textbf{Vérification (produit conservé) :} $Z_1 N_1 = 25 \times 1500 = 37\,500$. $Z_2 N_2 = 75 \times 500 = 37\,500$. \textbf{Égalité respectée.}
\end{enumerate}
\textbf{Réponse :} La roue de pompe tourne à \textbf{500 \si{tr/min}}. C'est un engrenage \textbf{réducteur} (rapport 3).
\end{exemplebox}

\courssub{Méthode générale de résolution}

\begin{methode}[Calculer une inconnue dans un engrenage (vitesse ou nombre de dents)]
\begin{enumerate}
    \item \textbf{Lire l'énoncé et identifier :} Quelle roue est \textbf{menante} (moteur, pédalier, entrée) ? Quelle roue est \textbf{menée} (roue arrière, pompe, sortie) ? Noter $Z_1, N_1, Z_2, N_2$ (connues et inconnue).
    \item \textbf{Écrire la relation fondamentale :} $Z_1 \times N_1 = Z_2 \times N_2$.
    \item \textbf{Isoler l'inconnue :} Faire le produit en croix ou diviser selon la place de l'inconnue.
    \item \textbf{Calculer le rapport $r = Z_2/Z_1$ (optionnel mais recommandé) :} Vérifier si $r>1$ (réducteur $\rightarrow N_2 < N_1$) ou $r<1$ (multiplicateur $\rightarrow N_2 > N_1$).
    \item \textbf{Conclure avec une phrase :} Donner la valeur numérique, l'unité (\si{tr/min} ou nombre de dents) et le type d'engrenage.
\end{enumerate}
\end{methode}

\courssub{Visualisation graphique : l'hyperbole de la transmission}

Pour bien sentir la relation inverse entre $Z_2$ et $N_2$ (quand $Z_1$ et $N_1$ sont fixés), traçons la courbe $N_2 = \frac{Z_1 N_1}{Z_2}$.

\begin{popfigure}
\begin{tikzpicture}
\begin{axis}[
    width=\linewidth, height=8cm,
    axis lines=middle,
    xlabel={$Z_2$ (dents de la roue menée)}, ylabel={$N_2$ (\si{tr/min})},
    xmin=0, xmax=100, ymin=0, ymax=2000,
    xtick={0,20,25,40,60,80,100}, ytick={0,500,1000,1500,2000},
    grid=major, grid style={popline!50},
    domain=10:100, samples=200,
    clip=false,
    title style={font=\small, align=center},
    title={Évolution de $N_2$ en fonction de $Z_2$ pour $Z_1=25$ dents, $N_1=1500\ \si{tr/min}$}
]
% Courbe hyperbole N2 = 37500 / Z2
\addplot[popcurve, thick, popblue] {37500/x};
% Points remarquables
\addplot[only marks, mark=*, mark size=2.5pt, poporange] coordinates {(25,1500)};
\addplot[only marks, mark=*, mark size=2.5pt, popgreen] coordinates {(75,500)};
\addplot[only marks, mark=*, mark size=2.5pt, poppurple] coordinates {(10,3750)};
% Lignes pointillées pour le point (75, 500)
\draw[popgreen, dashed, thin] (75,0) node[below, font=\tiny] {$Z_2=75$} -- (75,500) -- (0,500) node[left, font=\tiny] {$N_2=500$};
% Zones
\node[popgreenL, fill=popgreenL!50, rounded corners, inner sep=2pt, font=\tiny, anchor=south west] at (55, 1000) {Zone Réducteur ($r>1$, $N_2<N_1$)};
\node[poporangeL, fill=poporangeL!50, rounded corners, inner sep=2pt, font=\tiny, anchor=north east] at (20, 1800) {Zone Multiplicateur ($r<1$, $N_2>N_1$)};
% Ligne r=1 (Z2=Z1=25)
\draw[popdark, dashed, thin] (25,0) -- (25,2000) node[above, font=\tiny, pos=0.9] {$Z_2=Z_1$ ($r=1$)};
\end{axis}
\end{tikzpicture}
\legende{Courbe $N_2 = \frac{Z_1 N_1}{Z_2}$ (hyperbole). Pour $Z_1=25, N_1=1500$ : si on augmente $Z_2$ (roue menée plus grosse), $N_2$ chute rapidement. Le point $(75, 500)$ correspond à l'exemple de la pompe.}
\end{popfigure}

\courssub{Application concrète : le vélo, une boîte de vitesses à portée de main}

Le vélo est l'exemple parfait pour comprendre l'utilité de changer le rapport $r$. Sur un VTT ou un vélo de course camerounais (type \og Phoenix \fg{} ou \og Peugeot \fg{} qu'on voit partout), il y a :
\begin{itemize}
    \item À l'avant (pédalier) : 2 ou 3 plateaux (menants). Ex : \textbf{Petit plateau $Z_1 = 28$ dents}, \textbf{Grand plateau $Z_1 = 48$ dents}.
    \item À l'arrière (roue libre / cassette) : 6 à 11 pignons (menés). Ex : \textbf{Grand pignon $Z_2 = 34$ dents}, \textbf{Petit pignon $Z_2 = 11$ dents}.
\end{itemize}

\begin{center}
\begin{tabularx}{\linewidth}{|l|X|X|X|}\hline
\rowcolor{popblueL}
\textbf{Situation} & \textbf{Choix (Menante $\rightarrow$ Menée)} & \textbf{Rapport $r = Z_2/Z_1$} & \textbf{Résultat} \\ \hline
\textbf{Démarrage / Côte raide (Nkolbisson)} & Petit plateau (28) $\rightarrow$ Grand pignon (34) & $r = 34/28 \approx 1,21 > 1$ & \textbf{Réducteur fort.} Vitesse faible, force max. On monte sans s'épuiser. \\ \hline
\textbf{Plat / Vent favorable} & Grand plateau (48) $\rightarrow$ Pignon moyen (19) & $r = 19/48 \approx 0,40 < 1$ & \textbf{Multiplicateur.} Grande vitesse, effort modéré. \\ \hline
\textbf{Descente / Sprint} & Grand plateau (48) $\rightarrow$ Petit pignon (11) & $r = 11/48 \approx 0,23 \ll 1$ & \textbf{Multiplicateur fort.} Vitesse max, très peu de force. Attention aux genoux ! \\ \hline
\end{tabularx}
\end{center}

\begin{savaistu}[La boîte de vitesses automobile : une suite d'engrenages]
Dans une voiture (taxi jaune, 4x4 de la SONARA, bus de l'agence Voyageur), la \textbf{boîte de vitesses} n'est qu'un ensemble d'engrenages montés sur deux arbres (primaire et secondaire).
\begin{itemize}
    \item \textbf{1ère vitesse :} Petit pignon menant / Grande roue menée $\rightarrow r \approx 3,5$ à $4$. Couple maximal pour démarrer chargé.
    \item \textbf{4ème ou 5ème vitesse (prise directe ou overdrive) :} $r \approx 1$ ou $r < 1$. Vitesse de croisière, régime moteur bas (économie carburant, moins de bruit).
    \item \textbf{Marche arrière :} Un pignon intermédiaire (renvoi) est inséré $\rightarrow$ inverse le sens de rotation ($r$ négatif par convention).
\end{itemize}
Le levier de vitesse sélectionne \emph{quel couple de roues} est verrouillé sur l'arbre de sortie. C'est la technologie pure au service du confort et de l'efficacité !
\end{savaistu}

\courssub{Exercice résolu : dimensionner un engrenage pour un projet}

\begin{exoresolu}[Choisir les dents pour un ventilateur]
\textbf{Énoncé :} Un élève de 3ème veut construire un petit ventilateur pour sa chambre. Il a un petit moteur DC (récupéré sur un vieux jouet) qui tourne à \textbf{$N_1 = 6\,000\ \si{tr/min}$}. Il veut que l'hélice tourne à \textbf{$N_2 = 1\,500\ \si{tr/min}$} pour ne pas faire trop de bruit et avoir un bon flux d'air. Il dispose de roues dentées en plastique avec des nombres de dents standards : 10, 12, 15, 20, 24, 30, 36, 40, 48, 60 dents.
\textbf{Question :} Quel couple de roues $(Z_1, Z_2)$ choisir parmi le stock ? La roue menante est sur le moteur.

\tcblower
\textbf{Solution détaillée :}
\begin{enumerate}
    \item \textbf{Données :} $N_1 = 6000\ \si{tr/min}$ (menante), $N_2 = 1500\ \si{tr/min}$ (menée). Inconnues : $Z_1, Z_2$ (entiers dans la liste).
    \item \textbf{Calcul du rapport requis :} $r = \frac{N_1}{N_2} = \frac{6000}{1500} = 4$.
    \item \textbf{Interprétation :} $r = 4 > 1$, c'est un \textbf{réducteur}. Il faut $Z_2 > Z_1$. Précisément $Z_2 = 4 \times Z_1$.
    \item \textbf{Recherche dans la liste :} On teste les valeurs de $Z_1$ (menante, donc plutôt petite) :
        \begin{itemize}
            \item $Z_1 = 10 \Rightarrow Z_2 = 40$ (40 est dans la liste !) $\rightarrow$ \textbf{Candidat 1}.
            \item $Z_1 = 12 \Rightarrow Z_2 = 48$ (48 est dans la liste !) $\rightarrow$ \textbf{Candidat 2}.
            \item $Z_1 = 15 \Rightarrow Z_2 = 60$ (60 est dans la liste !) $\rightarrow$ \textbf{Candidat 3}.
            \item $Z_1 = 20 \Rightarrow Z_2 = 80$ (80 pas dans la liste).
        \end{itemize}
    \item \textbf{Choix technologique :} Les trois couples fonctionnent mathématiquement. On choisira souvent le couple le plus compact ou le plus robuste. $(Z_1=15, Z_2=60)$ donne des dents plus grandes (module plus gros si même diamètre primitif) donc plus solide pour un ventilateur qui vibre. $(Z_1=10, Z_2=40)$ est plus compact.
    \item \textbf{Vérification :} Pour $(15, 60)$ : $Z_1 N_1 = 15 \times 6000 = 90\,000$. $Z_2 N_2 = 60 \times 1500 = 90\,000$. \textbf{OK}.
\end{enumerate}
\textbf{Réponse :} Trois solutions possibles : $(10; 40)$, $(12; 48)$ ou $(15; 60)$. L'élève choisira $(15; 60)$ pour la robustesse.
\end{exoresolu}

\courssub{Sécurité et entretien : prolonger la vie des dents}

Un engrenage bien calculé ne sert à rien s'il casse au bout de deux semaines. Au Cameroun, avec la poussière rouge (latérite), l'humidité et les surcharges fréquentes (moto-taxi avec 3 passagers + bagages), l'entretien est vital.

\begin{aretenir}[Règles d'or pour des engrenages durables]
\begin{enumerate}
    \item \textbf{Lubrification :} Graisse (graisse au lithium, graisse verte) pour les vitesses lentes / charges lourdes (réducteur de moto, treuil). Huile (SAE 80W-90, ATF) pour les vitesses rapides (boîte de vitesses voiture, différentiel). \textbf{Vérifier le niveau tous les 5\,000 km ou 6 mois.}
    \item \textbf{Propreté :} La poussière est abrasive. Les carters doivent être étanches (joints spi en bon état). Nettoyer les évents (respirateurs) bouchés par la boue.
    \item \textbf{Alignement :} Arbres parallèles (pour engrenages cylindriques). Un mauvais alignement use les dents en biseau (contact sur le bord) $\rightarrow$ casse rapide.
    \item \textbf{Jeu de fonctionnement :} Un léger jeu entre dents est nécessaire (dilatation thermique, lubrifiant). Jeu trop grand = choc à chaque inversion (claquement). Jeu nul = grippage.
    \item \textbf{Ne jamais forcer à froid :} La graisse est épaisse. Laisser tourner au ralenti 2-3 min avant de charger (moto, groupe électrogène).
\end{enumerate}
\end{aretenir}

\begin{attention}[Signes d'usure ou de danger — Arrêtez la machine !]
\begin{itemize}
    \item \textbf{Bruit anormal :} \og Clac-clac \fg{} régulier (dent cassée), sifflement aigu (manque d'huile / mauvais alignement), grondement sourd (roulement mort).
    \item \textbf{Vibrations :} Senties au guidon, au levier de vitesse ou au plancher.
    \item \textbf{Fuite d'huile :} Au niveau des joints de carter ou des arbres de sortie. Perte de lubrifiant = mort programmée.
    \item \textbf{Échauffement :} Carter trop chaud pour y tenir la main (\textgreater 60-70°C). Vérifier niveau et qualité d'huile.
\end{itemize}
\end{attention}

\courssub{Bilan de la leçon : ce qu'il faut retenir pour le BEPC}

\begin{aretenir}[Synthèse : Les engrenages — Notions clés]
\begin{itemize}
    \item \textbf{Vocabulaire :} Menante ($Z_1, N_1$), Menée ($Z_2, N_2$), Prise directe, Pignon, Roue, Renvoi.
    \item \textbf{Sens de rotation :} Engrenage externe $\rightarrow$ sens opposés. Engrenage interne / Renvoi $\rightarrow$ même sens. Chaine / Courroie non croisée $\rightarrow$ même sens.
    \item \textbf{Relation fondamentale :} $\boxed{Z_1 N_1 = Z_2 N_2}$ (Conservation du nombre de dents par unité de temps).
    \item \textbf{Rapport de transmission :} $\boxed{r = \frac{Z_2}{Z_1} = \frac{N_1}{N_2}}$ (Sans unité).
    \item \textbf{Types :} $r > 1$ : Réducteur ($N_2 < N_1$, force $\uparrow$). $r < 1$ : Multiplicateur ($N_2 > N_1$, force $\downarrow$). $r = 1$ : Direct.
    \item \textbf{Méthode calcul :} Identifier menante/menée $\rightarrow$ Écrire $Z_1 N_1 = Z_2 N_2$ $\rightarrow$ Isoler l'inconnue $\rightarrow$ Vérifier cohérence ($r$ vs $N_2/N_1$).
    \item \textbf{Sécurité :} Lubrification adaptée, étanchéité, alignement, surveillance bruit/température.
\end{itemize}
\end{aretenir}

\begin{exercice}[Entraînement BEPC — Engrenage de moto-taxi]
Le moteur d'une moto-taxi (type \og Bajaj \fg{} ou \og TVS \fg{}) tourne à \SI{5000}{tr/min} au régime de croisière. Le pignon de sortie de boîte (menant) a \textbf{15 dents}. La couronne de roue arrière (menée) a \textbf{45 dents}.
\begin{enumerate}
    \item Calculer le rapport de transmission $r$ de la transmission finale. Quel type d'engrenage est-ce ?
    \item Calculer la vitesse de rotation $N_2$ de la roue arrière (en \si{tr/min}).
    \item Si le diamètre de la roue arrière est de \SI{60}{cm}, calculer la vitesse de la moto en \si{km/h} (arrondi au km/h). \textit{Indice : $v = \pi \times D \times N_2$ (en m/min), puis conversion.}
    \item Le conducteur veut plus de reprises en côte. Doit-il monter une couronne de \textbf{48 dents} ou de \textbf{42 dents} ? Justifier par le calcul du nouveau $r$ et de $N_2$.
\end{enumerate}
\end{exercice}

\begin{corrige}[Exercice Entraînement BEPC]
\begin{enumerate}
    \item $r = Z_2/Z_1 = 45/15 = 3$. $r > 1 \Rightarrow$ \textbf{Réducteur}.
    \item $N_2 = N_1 / r = 5000 / 3 \approx \textbf{1667 \si{tr/min}}$. (Vérif : $15 \times 5000 = 75000 = 45 \times 1666,7$).
    \item Périmètre roue $P = \pi \times 0,60 \approx 1,885$ m. $v = 1,885 \times 1667 \approx 3142$ m/min. $v = 3142 \times 60 / 1000 \approx \textbf{189 km/h}$. (Note : vitesse théorique max, en réalité pertes, aérodynamisme, régime max plus bas).
    \item Pour plus de force (reprises), il faut \textbf{augmenter $r$} (plus réducteur). Couronne 48 dents : $r = 48/15 = 3,2$ ($N_2 = 1563\ \si{tr/min}$). Couronne 42 dents : $r = 42/15 = 2,8$ ($N_2 = 1786\ \si{tr/min}$). Il faut choisir \textbf{48 dents} (rapport plus grand, roue tourne moins vite, plus de couple).
\end{enumerate}
\end{corrige}

\coursec{Sécurité et entretien des machines à engrenages}

Après avoir compris comment calculer le \cle{rapport de transmission} et prévoir le \cle{sens de rotation} des arbres, nous devons aborder un aspect tout aussi vital : la \cle{sécurité} et l'\cle{entretien}. Un engrenage bien conçu mais mal entretenu ou utilisé sans précaution devient une source de danger mortel. Dans nos ateliers, nos moulins de quartier ou sur les chantiers, le respect de quelques règles simples sauve des vies et préserve le matériel.

\courssub{Les risques liés aux engrenages : le point d'engrènement, zone critique}

\paragraph{Observation.}
Regarde une bétonnière en marche sur un chantier, ou le moulin à maïs du quartier quand il broie les grains. Tu vois les dents des roues s'emboîter les unes dans les autres. C'est là que la force se transmet. Mais c'est aussi là que se concentre le danger.

\begin{experience}[Visualiser la zone de happement]
\textbf{Matériel :} Deux roues dentées en plastique (type jeu de construction) ou deux engrenages métalliques légers montés sur un support, un morceau de tissu (vieux chiffon), des gants épais, lunettes de protection.
\textbf{Protocole :}
\begin{enumerate}
    \item Faire tourner lentement une roue motrice à la main (ou à très basse vitesse avec un moteur réduit) pour entraîner la roue menée.
    \item Approcher \emph{très prudemment} le morceau de tissu du \textbf{point d'engrènement} (là où les dents s'engagent) \textbf{sans y mettre les doigts}.
    \item Observer ce qui se passe au tissu.
\end{enumerate}
\textbf{Observations :} Le tissu est instantanément \textbf{happé}, tiré vers l'intérieur avec une force considérable, et se déchire ou s'enroule autour des arbres.
\textbf{Interprétation :} Au point d'engrènement, les surfaces des dents glissent les unes sur les autres avec une pression énorme. Tout corps mou (doigt, cheveu, vêtement large, chiffon) qui entre dans cet espace est aspiré par l'effet de coin et l'adhérence. La vitesse de happement est telle qu'il est \textbf{impossible de retirer sa main à temps}.
\textbf{Conclusion :} Le point d'engrènement est une \cle{zone de happement} (ou zone de pincement) \textbf{interdite d'accès} quand la machine tourne.
\legende{Expérience n°1 : Démonstration du happement au point d'engrènement. Ne jamais reproduire avec les doigts.}
\end{experience}

\begin{popfigure}
\begin{tikzpicture}[scale=1.2, >=Stealth]
    % Engrenage 1 (moteur)
    \foreach \i in {0,1,...,11} {
        \draw[popblue, thick, fill=popblueL] ({90-\i*30}:1.5) -- ({90-\i*30}:2.2) -- ({105-\i*30}:2.2) -- ({105-\i*30}:1.5) -- cycle;
    }
    \draw[popblue, thick, fill=white] (0,0) circle (1.2);
    \node[popblue, font=\bfseries] at (0,0) {Moteur};

    % Engrenage 2 (mené)
    \begin{scope}[xshift=3.7cm]
        \foreach \i in {0,1,...,11} {
            \draw[poporange, thick, fill=poporangeL] ({90-\i*30}:1.5) -- ({90-\i*30}:2.2) -- ({105-\i*30}:2.2) -- ({105-\i*30}:1.5) -- cycle;
        }
        \draw[poporange, thick, fill=white] (0,0) circle (1.2);
        \node[poporange, font=\bfseries] at (0,0) {Mené};
    \end{scope}

    % Flèches de rotation
    \draw[popdark, thick, ->] (1.8, 1.5) arc (90:-30:1.5);
    \node[popdark, above] at (1.8, 3) {Sens horaire};
    \draw[popdark, thick, ->] (3.7+1.8, 1.5) arc (90:210:1.5);
    \node[popdark, above] at (5.5, 3) {Sens anti-horaire};

    % Zone de danger
    \draw[poppink, very thick, dashed, decorate, decoration={snake, amplitude=1pt, segment length=4pt}] (1.5,0) -- (2.2,0);
    \draw[poppink, very thick, dashed, decorate, decoration={snake, amplitude=1pt, segment length=4pt}] (3.7+1.5,0) -- (3.7+2.2,0);
    \node[poppink, align=center, font=\bfseries\small] at (1.85, -0.8) {ZONE DE\\HAPPEMENT};
    \node[poppink, align=center, font=\bfseries\small] at (3.7+1.85, -0.8) {ZONE DE\\HAPPEMENT};

    % Carter de protection (représentation partielle)
    \draw[popgreen, very thick] (0,2.5) -- (3.7,2.5) -- (3.7,-2.5) -- (0,-2.5) -- cycle;
    \draw[popgreen, very thick] (0,2.5) arc (90:270:1.2);
    \draw[popgreen, very thick] (3.7,2.5) arc (90:-90:1.2);
    \node[popgreen, rotate=90, font=\small] at (-1.2,0) {CARTER DE PROTECTION};

    % Flèche danger
    \draw[poppink, very thick, ->] (1.85, -1.8) -- (1.85, -2.3);
    \node[poppink, below, font=\small] at (1.85, -2.5) {DANGER EXTRÊME};
\end{tikzpicture}
\legende{Schéma du point d'engrènement : les dents s'engagent (zone pointillée rose). Le carter vert empêche l'accès aux doigts. Les sens de rotation sont opposés.}
\end{popfigure}

\begin{attention}[Le point d'engrènement : la zone la plus dangereuse]
Dans toute machine à engrenages (moulin à farine, scie circulaire, bétonnière, boîte de vitesse de moto, groupe électrogène), \textbf{le point où les dents s'engrènent est la zone la plus dangereuse}.
\begin{itemize}
    \item \textbf{Force de happement :} Elle peut broyer les os d'un doigt en une fraction de seconde.
    \item \textbf{Vitesse :} Même à bas régime, le couple est maximal au démarrage. L'inertie des roues rend l'arrêt brutal impossible.
    \item \textbf{Corps étrangers :} Cheveux longs, manches larges, bijoux (chaînes, bracelets, bagues), chiffons de nettoyage sont aspirés instantanément.
\end{itemize}
\textbf{Règle d'or :} \emph{Jamais, au grand jamais, on n'approche les mains, un outil ou un chiffon d'un engrenage en mouvement.}
\end{attention}

\courssub{Les quatre consignes de sécurité vitales}

Ces consignes s'appliquent à \textbf{toute} machine comportant des organes en rotation (engrenages, courroies, poulies, arbres, lames).

\begin{propriete}[Consigne 1 : Distance de sécurité]
\textbf{Ne jamais approcher les mains d'un engrenage en mouvement.}
Garde une distance d'au moins \SI{50}{\cm} (longueur du bras) de la zone d'engrènement tant que la machine tourne. Utilise un outil long (bâton, tige) si tu dois écarter un déchet coincé \emph{à l'extérieur} de la zone de danger, mais \textbf{jamais} dans l'engrènement lui-même.
\end{propriete}

\begin{methode}[Procédure d'intervention en sécurité : « Arrêter – Débrancher – Vérifier – Intervenir »]
Avant \textbf{toute} intervention (réglage, graissage, changement de pièce, nettoyage, débouchage), suis \textbf{scrupuleusement} ces 4 étapes :
\begin{enumerate}
    \item \textbf{ARRÊTER} la machine par son bouton d'arrêt normal (ou couper le contact moteur).
    \item \textbf{DÉBRANCHER} l'alimentation : tirer la fiche secteur, couper le disjoncteur général, retirer la clé de contact (moto, groupe), fermer le robinet de gaz (groupe électrogène). \emph{Objectif : empêcher tout redémarrage accidentel (court-circuit, bouton appuyé par erreur, redémarrage auto après coupure de courant).}
    \item \textbf{VÉRIFIER L'IMMOBILISATION} : attendre l'arrêt complet des roues (inertie !), essayer de tourner l'arbre à la main (si accessible et sans danger) pour confirmer qu'il est libre et immobile.
    \item \textbf{INTERVENIR} : seulement alors, faire le graissage, le réglage ou la réparation.
\end{enumerate}
\textbf{Astuce pro :} Accroche ton cadenas personnel sur le disjoncteur ou la prise (consignation) si tu dois travailler longtemps : toi seul as la clé.
\end{methode}

\begin{propriete}[Consigne 3 : Les carters de protection]
\textbf{Utiliser les carters de protection prévus sur les machines.}
\begin{itemize}
    \item Sur une \textbf{bétonnière} : le carter qui couvre la couronne et le pignon.
    \item Sur un \textbf{moulin à maïs/manioc} : le carter qui entoure les meules ou les rouleaux.
    \item Sur une \textbf{scie circulaire} : le capot mobile qui recouvre la lame.
    \item Sur une \textbf{moto} : le carter de chaîne (ou le carter de boîte de vitesse).
\end{itemize}
\textbf{Interdiction formelle :} \emph{Ne jamais démonter un carter de protection pour « voir » ou « aller plus vite »}. Ne jamais faire fonctionner une machine dont le carter est absent, cassé ou mal fermé. C'est une infraction grave au code du travail et au bon sens.
\end{propriete}

\begin{propriete}[Consigne 4 : Tenue adaptée]
\textbf{Porter des vêtements ajustés, attacher les cheveux longs, ne pas porter de bijoux pendants.}
\begin{itemize}
    \item \textbf{Vêtements :} Pas de manches larges, pas de bas de pantalon évasé, pas de cordon de capuche qui pend. Rentrer la chemise dans le pantalon.
    \item \textbf{Cheveux :} Attachés solidement (queue de cheval, chignon) et rentrés sous une casquette ou un filet.
    \item \textbf{Bijoux :} \textbf{Aucun} anneau, bague, bracelet, chaîne, montre à bracelet métallique. Une bague coincée dans un engrenage arrache le doigt (dégantage) avant que la machine ne s'arrête.
    \item \textbf{EPI :} Lunettes de protection (projections de graisse, copeaux), chaussures de sécurité (écrasement), gants \textbf{uniquement pour manipuler des pièces à l'arrêt} (jamais sur machine en marche : le gant se coince et tire la main).
\end{itemize}
\end{propriete}

\begin{exemplebox}[Consignes affichées sur le moulin à maïs du quartier « Maman Pauline »]
À l'entrée du local, une pancarte peinte en rouge sur fond blanc rappelle :
\begin{center}
\begin{tabular}{|c|}\hline
\textbf{MOULIN À MAÏS – CONSIGNES DE SÉCURITÉ} \\ \hline
1. \textbf{INTERDIT} d'approcher les mains de la trémie en marche. \\ \hline
2. \textbf{ARRÊTER LE MOTEUR} et \textbf{DÉBRANCHER} avant tout graissage. \\ \hline
3. \textbf{CHEVEUX ATTACHÉS}, \textbf{MANCHES RETROUSSÉES}, \textbf{PAS DE BIJOUX}. \\ \hline
4. \textbf{NE JAMAIS} enlever le carter de protection des rouleaux. \\ \hline
5. En cas de bourrage : \textbf{COUPER LE CONTACT} $\rightarrow$ attendre l'arrêt $\rightarrow$ déboucher avec un \textbf{bâton}. \\ \hline
\textbf{LA VIE N'A PAS DE PRIX. LE MAÏS EN A UN.} \\ \hline
\end{tabular}
\end{center}
Cet affichage, simple et visible, a permis d'éviter plusieurs accidents graves depuis son installation.
\end{exemplebox}

\courssub{Entretien des engrenages : garantir la durée de vie et la sécurité}

Un engrenage mal entretenu s'use prématurément, fait du bruit, vibre, chauffe, et finit par casser (dents arrachées, arbre tordu), projetant des éclats métalliques dangereux.

\begin{definition}[Entretien préventif des engrenages]
L'entretien préventif consiste en trois actions régulières, planifiées selon la notice du constructeur ou l'intensité d'utilisation (quotidien, hebdomadaire, mensuel) :
\begin{enumerate}
    \item \textbf{Graissage / Lubrification :} Appliquer le \textbf{bon lubrifiant} (graisse épaisse pour engrenages ouverts à basse vitesse, huile de boîte pour engrenages en bain d'huile) aux \textbf{bons endroits} (dents, paliers, roulements) à la \textbf{bonne fréquence}. \emph{Exemple :} Graisser la couronne de la bétonnière chaque matin ; vidanger l'huile de la boîte de vitesse de la moto tous les \SI{10000}{\km}.
    \item \textbf{Vérification de l'usure des dents :} Contrôler visuellement (et au pied à coulisse si précis) l'épaisseur des dents, l'absence de \cle{piqûres} (petits trous dus à la fatigue), d'\cle{écaillage} (bords cassés), d'\cle{usure en escalier} (profil déformé). Une dent usée ne transmet plus correctement l'effort et bruite.
    \item \textbf{Remplacement des pièces endommagées :} Une dent cassée, un arbre voilé, un roulement qui gronde doivent être remplacés \textbf{immédiatement}. Ne jamais « souder » une dent cassée : la structure métallique est modifiée, la rupture sera brutale.
\end{enumerate}
\end{definition}

\begin{experience}[Contrôle visuel et auditif d'un engrenage (TP atelier)]
\textbf{Matériel :} Engrenage démonté (pignon + couronne) ou machine accessible (moto sur béquille, moulin arrêté et débranché), lampe torche, chiffon, graisse adaptée, clé à pipe.
\textbf{Protocole :}
\begin{enumerate}
    \item \textbf{Nettoyer} les dents avec le chiffon (enlever l'ancienne graisse noircie, la poussière, le sable).
    \item \textbf{Inspecter} chaque dent : chercher des fissures, des manques de matière, une usure asymétrique.
    \item \textbf{Tester le jeu latéral} (mouvement de va-et-vient angulaire) : il doit être faible et régulier.
    \item \textbf{Faire tourner} l'arbre lentement à la main : sentir les points durs, écouter les cliquetis ou grattements.
    \item \textbf{Regraisser} : déposer un filet de graisse neuve sur le flanc de chaque dent (côté poussée).
    \item \textbf{Remonter} le carter, revisser les boulons au couple.
\end{enumerate}
\textbf{Interprétation :} Un engrenage propre, bien graissé, sans jeu anormal et silencieux à la main tournera longtemps sans chauffer. Le sable et la vieille graisse sont abrasifs : ils usent les dents 10 fois plus vite.
\legende{TP Entretien : Nettoyage, inspection, graissage. Toujours machine à l'arrêt et débranchée.}
\end{experience}

\begin{attention}[Pièges fréquents lors de l'entretien]
\begin{itemize}
    \item \textbf{Graisser « par-dessus » l'ancienne graisse sale :} Le mélange sable/vieille graisse/nouvelle graisse devient une pâte à roder qui use les dents en quelques heures. \textbf{Nettoyer avant de graisser.}
    \item \textbf{Mettre trop de graisse :} L'excès chauffe par cisaillement, fait fuir les joints, projette de la graisse partout. \textbf{Quantité modérée, régulière.}
    \item \textbf{Utiliser de l'huile moteur au lieu de graisse d'engrenage :} L'huile coule, ne tient pas sur les dents verticales, ne protège pas contre les chocs. \textbf{Respecter la préconisation constructeur (grade EP, viscosité).}
    \item \textbf{Oublier de refermer le carter :} Après l'entretien, \textbf{revisser tous les boulons}. Un carter mal fermé vibre, se déforme, laisse entrer la poussière et l'eau.
\end{itemize}
\end{attention}

\courssub{Sensibilisation : la réalité des accidents dans nos ateliers}

\begin{savaistu}[Accidents de machines agricoles et artisanales au Cameroun]
La plupart des accidents graves (amputations de doigts, de mains, scalping, fractures) survenus dans les \textbf{moulins villageois} (maïs, manioc, arachide), les \textbf{scieries artisanales}, les \textbf{ateliers de soudure/mécanique} et sur les \textbf{chantiers} (bétonnières) ont une cause commune : \textbf{le non-respect de la procédure « Arrêter – Débrancher – Vérifier » avant intervention}.
\begin{itemize}
    \item \textbf{Scénario type 1 :} Le moulin bourre. L'opérateur essaie de pousser le maïs avec la main ou un bâton \emph{sans couper le moteur}. La main est happée par les rouleaux.
    \item \textbf{Scénario type 2 :} On veut graisser la chaîne de la moto ou le pignon de la bétonnière « vite fait » pendant qu'elle tourne au ralenti. Le chiffon ou le gant est happé.
    \item \textbf{Scénario type 3 :} Après une coupure de courant ENEO, le groupe électrogène redémarre seul (démarrage auto) pendant que le technicien a la main sur le carter ouvert pour vérifier l'huile.
\end{itemize}
\textbf{Statistique :} Selon les rapports de la CNPS (Caisse Nationale de Prévoyance Sociale) et des inspections du travail, \textbf{plus de 70\%} des accidents du travail impliquant des machines tournantes sont dus à une \textbf{absence de consignation (coupure effective de l'énergie)} avant intervention.
\textbf{Message :} La paresse, la précipitation ou l'oubli d'une seule consigne changent une vie à jamais. \textbf{Toi, futur technicien, ingénieur ou artisan, tu as le devoir de faire respecter ces règles autour de toi.}
\end{savaistu}

\courssub{Exercice résolu : Analyse de situation à risque}

\begin{exoresolu}[Situation au moulin du quartier]
{Énoncé :}
Monsieur K. gère un moulin à maïs à moteur thermique. Un client apporte un sac de maïs humide. Le moulin bourbe (bourrage). Monsieur K., pressé par la file d'attente, introduit une main gantée dans la trémie pour débloquer les grains, pendant que le moteur tourne au ralenti. Son gant est happé, son doigt est broyé.
\begin{enumerate}
    \item Identifier \textbf{trois} fautes de sécurité commises par Monsieur K.
    \item Expliquer pourquoi le gant a \textbf{aggravé} la blessure au lieu de la protéger.
    \item Rédiger la procédure correcte qu'il aurait dû suivre (étapes numérotées).
\end{enumerate}
\tcblower
{Solution détaillée :}
\begin{enumerate}
    \item \textbf{Faute 1 :} Intervention sur machine en marche (non-respect de l'arrêt complet).
    \item \textbf{Faute 2 :} Introduction de la main (même gantée) dans la zone d'alimentation/trémie qui communique avec la zone d'engrènement des rouleaux.
    \item \textbf{Faute 3 :} Port de gants \textbf{sur machine en mouvement}. Le gant est un facteur aggravant majeur.
    \item \textbf{Pourquoi le gant aggrave :} Sur machine arrêtée, le gant protège des coupures/saletés. Sur machine en marche, le tissu du gant s'accroche aux aspérités des rouleaux ou des dents \textbf{bien plus fort que la peau nue}. Au lieu de glisser ou de se déchirer net, le gant s'enroule autour de l'arbre et \textbf{tire le doigt/la main à l'intérieur} avec une force décuplée. Le gant devient un piège qui arrache les tendons et broie l'os. \textbf{Règle : JAMAIS de gants près d'organes en rotation.}
    \item \textbf{Procédure correcte :}
    \begin{enumerate}
        \item Couper le contact du moteur thermique (clé sur « OFF » ou bouton d'arrêt d'urgence).
        \item Attendre l'arrêt \textbf{complet} et \textbf{silencieux} des rouleaux (inertie du volant).
        \item Vérifier l'immobilisation en essayant de tourner la poulie à la main (si accessible sans danger).
        \item Déboucher la trémie avec un \textbf{bâton en bois} ou une tige métallique longue (jamais la main).
        \item Refermer les trappes, redémarrer le moteur, reprendre le travail.
    \end{enumerate}
\end{enumerate}
\end{exoresolu}

\courssub{Bilan de la leçon : Les engrenages, puissance et responsabilité}

\begin{aretenir}[Essentiel sur les engrenages – Sécurité \& Entretien]
\begin{itemize}
    \item \textbf{Point d'engrènement = Zone de happement mortelle.} Interdiction absolue d'y approcher mains, outils, chiffons, vêtements, cheveux, bijoux.
    \item \textbf{4 Consignes vitales :} 1. Distance de sécurité. 2. \textbf{Arrêter → Débrancher → Vérifier → Intervenir} (Consignation). 3. Carters de protection \textbf{toujours} en place et fermés. 4. Tenue ajustée, cheveux attachés, \textbf{pas de bijoux, pas de gants en rotation}.
    \item \textbf{Entretien = Nettoyer + Inspecter + Graisser (juste dose, bon produit) + Refermer carter.}
    \item \textbf{Usure anormale (bruit, jeu, chauffe) = Arrêt immédiat + Réparation/Remplacement.}
    \item \textbf{Responsabilité :} Appliquer et faire appliquer ces règles protège ton intégrité physique et celle de tes collègues. C'est la marque du vrai professionnel.
\end{itemize}
\end{aretenir}

\begin{exercice}[Mise en situation : La bétonnière du chantier]
Sur un chantier de construction à Yaoundé, une bétonnière électrique (monophasé \SI{230}{\V}) fonctionne. L'opérateur veut graisser la couronne dentée et le pignon.
\begin{enumerate}
    \item Dessine le schéma simplifié (couronne + pignon + carter) et entoure en rouge la zone de happement.
    \item Liste dans l'ordre les 4 actions exactes qu'il doit faire \textbf{avant} d'ouvrir le carter pour graisser.
    \item Il a une bombe de graisse « multi-usage » et un pot de graisse « engrenages EP 2 ». Laquelle choisir ? Justifie.
    \item Pourquoi est-il interdit de graisser « par-dessus » l'ancienne graisse noircie par la poussière de ciment ?
\end{enumerate}
\end{exercice}

\begin{corrige}[Exercice : La bétonnière du chantier]
\begin{enumerate}
    \item \textit{Schéma :} Deux cercles en contact (gros = couronne, petit = pignon). Zone de contact hachurée en rouge : « ZONE DE HAPPEMENT ». Rectangle autour = Carter.
    \item \textbf{1. ARRÊTER} (bouton stop). \textbf{2. DÉBRANCHER} la fiche \SI{230}{\V} (ou couper disjoncteur chantier). \textbf{3. VÉRIFIER} l'arrêt total (attendre 30 s, vérifier visuellement). \textbf{4. OUVRIR} le carter et graisser.
    \item \textbf{Graisse « engrenages EP 2 »}. La graisse multi-usage n'a pas les additifs « Extrême Pression » (EP) nécessaires pour supporter les chocs et pressions énormes entre les dents d'une bétonnière chargée. Elle serait éjectée ou écrasée en quelques minutes.
    \item L'ancienne graisse est contaminée par le \textbf{sable} et la \textbf{poussière de ciment} (abrasifs). Graisser par-dessus crée une « pâte à roder » qui use les dents 10 à 20 fois plus vite qu'un métal nu. Il faut \textbf{nettoyer au chiffon/solvant} avant de regraisser.
\end{enumerate}
\end{corrige}

\begin{savaistu}[Le « cadenas de consignation » : l'outil du pro]
Dans l'industrie et les grands ateliers (SONARA, ENEO, usines, grands chantiers), on utilise un \textbf{cadenas de consignation} (souvent rouge, avec une étiquette « NE PAS ACTIONNER – MAINTENANCE EN COURS – NOM : ... »).
\begin{enumerate}
    \item L'opérateur arrête la machine.
    \item Il met \textbf{son cadenas personnel} sur le disjoncteur ou la prise verrouillée.
    \item Il garde \textbf{la seule clé} dans sa poche.
    \item Personne d'autre ne peut redémarrer la machine tant que son cadenas est là.
\end{enumerate}
C'est la \textbf{seule} méthode qui garantit à 100\% qu'aucun collègue, aucune automatique, aucune surtension ne redémarrera la machine pendant que tes mains sont dedans. Adopte ce réflexe dès maintenant, même pour ton groupe électrogène ou ta moto : \textbf{clé de contact dans TA poche} = consignation simple et efficace.
\end{savaistu}

\newpage

\coursec{Activité d'intégration}

\coursec{Les engrenages}

\courssub{Objectifs de la séance}
À l'issue de cette séance, tu seras capable de :
\begin{itemize}
    \item Identifier et nommer les organes d'un engrenage (roue menante, roue menée, pignon, dent, cercle primitif).
    \item Déterminer le sens de rotation des roues engrenées (engrenage externe, interne, vis sans fin).
    \item Calculer le \cle{rapport de transmission} $R$ et l'utiliser pour trouver une vitesse de rotation inconnue.
    \item Caractériser une transmission en \cle{réducteur} ou \cle{multiplicateur} et en déduire l'évolution du couple.
    \item Citer les principaux types d'engrenages, leurs avantages, inconvénients et domaines d'application.
    \item Appliquer les consignes de sécurité et les règles d'entretien (graissage, alignement) sur une machine à engrenages.
\end{itemize}

\courssub{Découverte expérimentale : Manipulation d'engrenages}
\begin{experience}[Observation de kits d'engrenages (ou vélos, perceuses, moulins)]
\textbf{Matériel :} Kit de roues dentées (différents diamètres, nombres de dents), arbres, support ; ou observation d'une transmission de vélo / perceuse / moulin à maïs.
\textbf{Protocole :}
\begin{enumerate}
    \item Assembler deux roues de tailles différentes (ex: 12 dents et 36 dents) sur deux axes parallèles. Faire tourner la petite roue (menante) d'un tour complet. Compter le nombre de tours de la grande roue (menée). Noter le sens de rotation des deux roues.
    \item Répéter en faisant tourner la grande roue (menante) : compter les tours de la petite roue (menée). Noter les sens.
    \item Essayer d'engrener trois roues de suite. Observer le sens de la troisième par rapport à la première.
    \item (Si possible) Bloquer légèrement la roue menée avec la main : ressentir l'effort (couple) sur la menante. Comparer selon le rapport de dents.
\end{enumerate}
\textbf{Observations attendues :}
\begin{itemize}
    \item Deux roues extérieures tournent en \textbf{sens inverse}.
    \item La petite roue fait plus de tours que la grande (vitesse plus grande).
    \item Quand on mène la grande roue, la petite tourne plus vite (multiplicateur) mais l'effort à fournir sur la grande est plus grand (couple résistant plus fort).
    \item Avec trois roues : la 1ère et la 3ème tournent dans le \textbf{même sens}.
\end{itemize}
\textbf{Interprétation :} Le nombre de dents qui s'engrènent est le même pour les deux roues. La vitesse de rotation est inversement proportionnelle au nombre de dents. Le sens s'inverse à chaque contact externe.
\end{experience}

\courssub{Définitions et vocabulaire}
\begin{definition}[Organes d'un engrenage]
Un \textbf{engrenage} est un mécanisme de transmission de mouvement et de force par contact de dents entre deux roues.
\begin{itemize}
    \item \textbf{Roue menante (ou pignon) :} Reçoit l'énergie du moteur (arbre d'entrée). C'est le \textbf{moteur} du système.
    \item \textbf{Roue menée :} Transmet l'énergie à l'organe récepteur (meule, roue de véhicule, arbre de sortie). C'est l'\textbf{utilisateur}.
    \item \textbf{Dent :} Partie saillante assurant le contact. Profil le plus souvent \textbf{évolutif} (courbe assurant un mouvement régulier).
    \item \textbf{Cercle primitif :} Cercle imaginaire où s'effectue théoriquement le roulement sans glissement des deux roues. Son rayon $R$ est proportionnel au nombre de dents $Z$ : $R = \frac{m \cdot Z}{2}$ ($m$ = module, en mm).
    \item \textbf{Module $m$ :} Caractéristique de la taille des dents. \textbf{Deux roues engrenées ont obligatoirement le même module.}
\end{itemize}
\end{definition}

\begin{propriete}[Condition d'engrenage]
Pour que deux roues puissent engrener correctement, elles doivent avoir \textbf{le même module $m$} (et le même angle de pression, généralement 20°).
\[
m = \frac{2R}{Z} = \frac{\text{Pas}}{\pi} \quad (\text{en mm})
\]
\end{propriete}

\courssub{Types d'engrenages et sens de rotation}
\begin{popfigure}
\begin{tikzpicture}[scale=0.7, every node/.style={font=\small}]
% Engrenage droit externe
\begin{scope}[xshift=0cm]
\draw[fill=popblue!10, draw=popblue, thick] (0,0) circle (1.5);
\foreach \a in {0,30,...,330} \draw[popblue, thick] ({1.2*cos(\a)},{1.2*sin(\a)}) -- ({1.5*cos(\a)},{1.5*sin(\a)});
\draw[fill=popblue!10, draw=popblue, thick] (4,0) circle (1);
\foreach \a in {0,30,...,330} \draw[popblue, thick] ({4+0.7*cos(\a)},{0.7*sin(\a)}) -- ({4+1*cos(\a)},{1*sin(\a)});
\node[align=center] at (2,-2.2) {\textbf{Engrenage droit externe}\\(Axes parallèles, sens inversés)};
\draw[->, thick, poporange] (0,1.8) arc (90:-270:0.4);
\draw[->, thick, popgreen] (4,1.3) arc (90:450:0.4);
\end{scope}

% Engrenage droit interne
\begin{scope}[xshift=7cm]
\draw[fill=popgreen!10, draw=popgreen, thick] (0,0) circle (2);
\foreach \a in {0,20,...,340} \draw[popgreen, thick] ({1.7*cos(\a)},{1.7*sin(\a)}) -- ({2*cos(\a)},{2*sin(\a)});
\draw[fill=poporange!10, draw=poporange, thick] (0,0) circle (0.8);
\foreach \a in {0,20,...,340} \draw[poporange, thick] ({0.8*cos(\a)},{0.8*sin(\a)}) -- ({0.5*cos(\a)},{0.5*sin(\a)});
\node[align=center] at (0,-2.7) {\textbf{Engrenage droit interne}\\(Axes parallèles, \textbf{même sens})};
\draw[->, thick, poporange] (0,1) arc (90:-270:0.4);
\draw[->, thick, popgreen] (0,1.8) arc (90:-270:0.4);
\end{scope}

% Engrenage conique
\begin{scope}[xshift=0cm, yshift=-4.5cm]
\draw[fill=poppurple!10, draw=poppurple, thick] (0,0) ellipse (1.5 and 0.5);
\draw[poppurple, thick] (-1.5,0) -- (-0.5,1.5) -- (0.5,1.5) -- (1.5,0);
\draw[poppurple, thick] (-1.5,0) -- (-0.5,-1.5) -- (0.5,-1.5) -- (1.5,0);
\draw[fill=poppink!10, draw=poppink, thick] (2.5,1.5) ellipse (0.5 and 1.5);
\draw[poppink, thick] (2.5,0) -- (3.5,0.5) -- (3.5,-0.5) -- (2.5,0);
\draw[poppink, thick] (2.5,3) -- (3.5,2.5) -- (3.5,3.5) -- (2.5,3);
\node[align=center] at (1.2,-2.5) {\textbf{Engrenage conique}\\(Axes concourants, souvent 90°)};
\draw[->, thick, poppurple] (0,0.8) arc (90:-90:0.3);
\draw[->, thick, poppink] (2.5,2.3) arc (0:180:0.3);
\end{scope}

% Vis sans fin
\begin{scope}[xshift=7cm, yshift=-4.5cm]
\draw[fill=popgold!10, draw=popgold, thick] (0,0) ellipse (1.5 and 0.4);
\foreach \x in {-1.2,-0.6,0,0.6,1.2} \draw[popgold, thick] (\x,-0.4) -- (\x+0.1,0) -- (\x,-0.4) -- (\x-0.1,0) -- cycle;
\draw[fill=popblue!10, draw=popblue, thick] (3,0) circle (1.2);
\foreach \a in {0,15,...,345} \draw[popblue, thick] ({3+0.9*cos(\a)},{0.9*sin(\a)}) -- ({3+1.2*cos(\a)},{1.2*sin(\a)});
\node[align=center] at (1.5,-2.5) {\textbf{Vis sans fin + Roue}\\(Axes perpendiculaires, non réversible)};
\draw[->, thick, popgold] (0,0.6) -- (-1.5,0.6);
\draw[->, thick, popblue] (3,1.4) arc (90:450:0.4);
\end{scope}
\end{tikzpicture}
\legende{Les quatre grands types d'engrenages rencontrés en technologie.}
\end{popfigure}

\begin{propriete}[Sens de rotation selon le type]
\begin{itemize}
    \item \textbf{Engrenage externe (droits, hélicoïdaux) :} Axes parallèles $\rightarrow$ Sens \textbf{inversés}.
    \item \textbf{Engrenage interne :} Axes parallèles $\rightarrow$ \textbf{Même sens}.
    \item \textbf{Engrenage conique :} Axes concourants (souvent 90°) $\rightarrow$ Sens selon géométrie (s'inversent si vue de face).
    \item \textbf{Vis sans fin + roue :} Axes perpendiculaires non sécants. La vis (menante) entraîne la roue. \textbf{Non réversible} (la roue ne peut pas entraîner la vis si angle d'hélice petit). Utilisé pour le verrouillage (treuil, porte de garage).
\end{itemize}
\end{propriete}

\courssub{Avantages et inconvénients des engrenages}
\begin{aretenir}[Bilan comparatif]
\begin{center}
\begin{tabularx}{\linewidth}{|l|X|X|}\hline
\rowcolor{popblueL} \textbf{Critère} & \textbf{Avantages} & \textbf{Inconvénients} \\ \hline
\textbf{Transmission} & Rapport de transmission \textbf{exact, constant} (pas de glissement comme courroie). & Bruyant à grande vitesse (claquement dents). \\ \hline
\textbf{Puissance / Couple} & Supporte \textbf{très forts couples} et puissances élevées (moteurs camions, moulins, industrie). & Exige \textbf{lubrification} rigoureuse (graisse/huile) et alignement précis. \\ \hline
\textbf{Encombrement} & \textbf{Compact} pour gros rapports (surtout vis sans fin, planétaires). & Coût de fabrication élevé (usinage dents de précision). \\ \hline
\textbf{Sens} & Permet inversion sens (externe) ou angle 90° (conique, vis). & Usure par frottement (poussière abrasive = danger). \\ \hline
\end{tabularx}
\end{center}
\textbf{Applications courantes au Cameroun :} Boîte de vitesse moto/voiture (droits/hélicoïdaux), différentiel (coniques), moulin à maïs (droits en fonte), treuil de forage (vis sans fin), réducteurs industriels (planétaires).
\end{aretenir}

\courssub{Rapport de transmission}
\begin{propriete}[Définition et formules]
Pour un engrenage (externe ou interne), le \textbf{rapport de transmission $R$} est défini par :
\[
R = \frac{Z_{\text{menée}}}{Z_{\text{menante}}} = \frac{N_{\text{menante}}}{N_{\text{menée}}} = \frac{\omega_{\text{menante}}}{\omega_{\text{menée}}}
\]
\begin{itemize}
    \item $Z$ : nombre de dents (entier, sans unité).
    \item $N$ : vitesse de rotation en \textbf{tr/min} (tours par minute) ou $\omega$ en \textbf{rad/s}.
    \item $R$ est un \textbf{nombre sans unité}.
\end{itemize}
\emph{Remarque :} Pour une chaîne de plusieurs engrenages, le rapport global est le produit des rapports de chaque engrenage ($R_{\text{tot}} = R_1 \times R_2 \times \dots$). Une roue intermédiaire (porte-satellite) ne change pas la valeur de $R$ (simplification des $Z$), elle ne change que le sens.
\end{propriete}

\begin{methode}[Calculer une vitesse de sortie]
\begin{enumerate}
    \item Identifier la roue menante ($Z_1$, $N_1$) et la roue menée ($Z_2$).
    \item Calculer $R = \frac{Z_2}{Z_1}$.
    \item Calculer $N_2 = \frac{N_1}{R}$.
    \item Conclure : Réducteur ($R>1$, $N_2<N_1$) ou Multiplicateur ($R<1$, $N_2>N_1$).
\end{enumerate}
\end{methode}

\begin{exemplebox}[Exemple rapide]
Un pignon de 15 dents entraîne une roue de 60 dents. Moteur à 1800 tr/min.
$R = 60/15 = 4$. $N_{\text{sortie}} = 1800/4 = 450$ tr/min. C'est un réducteur ($R=4>1$). Le couple est multiplié par 4.
\end{exemplebox}

\courssub{Sécurité et entretien des machines à engrenages}
\begin{attention}[Risques majeurs et prévention]
\begin{itemize}
    \item \textbf{Risque d'entraînement :} Vêtements amples, cheveux longs, gants, bijoux aspirés dans l'engrenage $\rightarrow$ \textbf{Arrachage, broyage membres}. \textbf{Parade :} Carter \textbf{fermé et verrouillé} en fonctionnement. Arrêt total (cadenassage) avant maintenance. Pas de gants près parties tournantes.
    \item \textbf{Risque de projection :} Dents cassées (fatigue, choc), éclats de pièce usinée, lubrifiant projeté. \textbf{Parade :} Lunettes de protection obligatoires. Carter étanche.
    \item \textbf{Risque thermique / incendie :} Échauffement par manque de graisse, frottement excessif. \textbf{Parade :} \textbf{Graissage régulier} (graisse épaisse type NLGI 2 ou huile EP selon vitesse/charge). Vérifier niveau huile bain d'huile. Contrôler température (main / sonde).
    \item \textbf{Alignement :} Arbres non parallèles $\rightarrow$ usure en biseau, bruit, casse prématurée. \textbf{Parade :} Contrôle alignement (règle, comparateur) au montage et après vibrations.
\end{itemize}
\end{attention}

\courssub{Activité d'intégration : Le moulin à maïs de Nkolbisson}
\begin{popfigure}
\begin{tikzpicture}[scale=0.9, every node/.style={font=\small}]
    % Moteur
    \draw[fill=popblue!20, draw=popblue, thick, rounded corners=5pt] (0,0) rectangle (3,2);
    \node[align=center] at (1.5,1) {\textbf{Moteur électrique}\\(Monophasé 220 V)\\(1440 tr/min)};
    \draw[->, thick, poporange] (3,1) -- (4,1) node[midway, above] {Arbre moteur};
    
    % Pignon (menant)
    \draw[fill=poporange!30, draw=poporange, thick] (5,1) circle (1.2cm);
    \foreach \angle in {0,15,...,345} {
        \draw[poporange, thick] ({5+0.9*cos(\angle)},{1+0.9*sin(\angle)}) -- ({5+1.2*cos(\angle)},{1+1.2*sin(\angle)});
    }
    \node at (5,1) {$Z_1 = 24$};
    \node[below] at (5,-0.5) {\textbf{Pignon (Roue menante)}};
    
    % Engrenement
    \draw[<->, thick, popgreen, dashed] (5,1) -- (9,1);
    
    % Roue menée (Meule)
    \draw[fill=popgreen!20, draw=popgreen, thick] (10,1) circle (2.5cm);
    \foreach \angle in {0,11.25,...,348.75} {
        \draw[popgreen, thick] ({10+2.2*cos(\angle)},{1+2.2*sin(\angle)}) -- ({10+2.5*cos(\angle)},{1+2.5*sin(\angle)});
    }
    \node at (10,1) {$Z_2 = 96$};
    \node[below] at (10,-1.5) {\textbf{Grosse roue (Roue menée)}};
    
    % Meule
    \draw[fill=poppurple!20, draw=poppurple, thick, rounded corners=5pt] (10,-3.5) rectangle (14,-2.5);
    \node[align=center] at (12,-3) {\textbf{Meule (Pierre à moudre)}\\(Solidaire de la roue menée)};
    \draw[thick, poppurple] (10,1) -- (10,-2.5) node[midway, right] {Arbre de la meule};
    
    % Sens de rotation Moteur
    \draw[->, thick, popblue] (2.5,2.5) arc (90:-270:0.5);
    \node[popblue] at (2.5,3.2) {Sens horaire};
    
    % Flèche sens pignon
    \draw[->, thick, poporange] (6.2,1.2) arc (45:-225:0.7);
    \node[poporange] at (6.2,2.2) {?};
    
    % Flèche sens roue
    \draw[->, thick, popgreen] (11.5,2.5) arc (45:405:1.0);
    \node[popgreen] at (11.5,3.8) {?};
\end{tikzpicture}
\legende{Schéma simplifié de la transmission du moulin à maïs. Le moteur entraîne le pignon (24 dents) qui engrène la grosse roue (96 dents) solidaire de la meule.}
\end{popfigure}

Dans le village de \textbf{Nkolbisson}, à la périphérie de Yaoundé, le moulin à maïs du groupement d'intérêt économique (GIE) « \emph{Ngono Nkwa} » fonctionne grâce à un moteur électrique monophasé \SI{220}{V} récupéré sur un vieux groupe électrogène. Ce moteur tourne à une vitesse de rotation \textbf{$N_{\text{moteur}} = 1440$~tr/min}.

L'arbre du moteur est solidaire d'un \textbf{pignon} de \textbf{$Z_1 = 24$~dents}. Ce pignon engrène directement une \textbf{grosse roue} de \textbf{$Z_2 = 96$~dents}, fixée sur le même arbre que la meule en pierre qui écrase le maïs pour produire la farine de \emph{ndolé} et de \emph{kwem}.

Le président du GIE, M. \emph{Atangana}, souhaite comprendre le fonctionnement de sa machine pour mieux l'entretenir et former les jeunes apprentis. Il a remarqué que la meule tourne « moins vite » que le moteur, mais qu'elle a « plus de force » pour broyer les grains durs.

\courssub{Consignes}
À partir du schéma et des données ci-dessus, réponds aux questions suivantes en justifiant chaque étape de ton raisonnement.

\begin{enumerate}
    \item \textbf{Identification des organes :} Nomme la \cle{roue menante} et la \cle{roue menée} dans ce système. Justifie ton choix en te basant sur la définition du cours.
    \item \textbf{Sens de rotation :} Le moteur tourne dans le sens horaire (comme les aiguilles d'une montre). Détermine le sens de rotation de la meule. Représente-le par une flèche sur le schéma de ton papier.
    \item \textbf{Rapport de transmission et vitesse de la meule :}
    \begin{enumerate}
        \item Calcule le \cle{rapport de transmission} $R$ de cet engrenage.
        \item Calcule la vitesse de rotation $N_{\text{meule}}$ de la meule en tr/min.
    \end{enumerate}
    \item \textbf{Nature de la transmission :} S'agit-il d'un \cle{réducteur} ou d'un \cle{multiplicateur} ? Justifie ta réponse par le calcul et par l'analyse des vitesses.
    \item \textbf{Intérêt technologique pour le meunier :} M. Atangana dit : « La meule tourne moins vite, mais elle a plus de force. » Explique scientifiquement pourquoi cette affirmation est correcte en utilisant la notion de \cle{couple} (moment de force) et la conservation de l'énergie (puissance). Pourquoi est-ce indispensable pour moudre du maïs sec ?
    \item \textbf{Consignes de sécurité :} Propose \textbf{deux consignes de sécurité} impératives à afficher sur le carter de protection du moulin pour protéger les opérateurs (risques mécaniques).
    \item \textbf{Prolongement (Modification du système) :} Pour augmenter la production, on remplace la grosse roue de 96 dents par une roue de \textbf{48 dents} (le pignon reste à 24 dents).
    \begin{enumerate}
        \item Calcule le nouveau rapport de transmission $R'$ et la nouvelle vitesse de la meule $N'_{\text{meule}}$.
        \item Que se passe-t-il pour la force de broyage (couple) ? Est-ce une bonne idée pour moudre du maïs très sec ? Justifie.
    \end{enumerate}
\end{enumerate}

\courssub{Rappels de cours utiles pour l'activité}
\begin{definition}[Roue menante et roue menée]
La \textbf{roue menante} (ou pignon) est celle qui \emph{reçoit} l'énergie du moteur (arbre d'entrée). La \textbf{roue menée} est celle qui \emph{transmet} l'énergie à l'organe récepteur (meule, roue de voiture, etc.).
\end{definition}

\begin{propriete}[Sens de rotation en engrenage externe]
Deux roues à denture extérieure engrenées tournent \textbf{toujours en sens inverse} l'une de l'autre.
\end{propriete}

\begin{propriete}[Rapport de transmission $R$]
Pour un engrenage, le rapport de transmission est le rapport du nombre de dents de la roue menée sur le nombre de dents de la roue menante. Il relie les vitesses de rotation :
\[
R = \frac{Z_{\text{menée}}}{Z_{\text{menante}}} = \frac{N_{\text{menante}}}{N_{\text{menée}}}
\]
\emph{Attention :} $R$ est un nombre \textbf{sans unité}. $N$ s'exprime en tr/min (ou en rad/s).
\end{propriete}

\begin{propriete}[Réducteur et Multiplicateur]
\begin{itemize}
    \item \textbf{Réducteur :} $R > 1$ (ou $N_{\text{sortie}} < N_{\text{entrée}}$). La vitesse diminue, le couple \textbf{augmente}.
    \item \textbf{Multiplicateur :} $R < 1$ (ou $N_{\text{sortie}} > N_{\text{entrée}}$). La vitesse augmente, le couple \textbf{diminue}.
\end{itemize}
\end{propriete}

\begin{propriete}[Puissance et Couple (Niveau 3ème - Approche qualitative/quantitative)]
En négligeant les pertes par frottement (rendement $\approx 1$), la \textbf{puissance} est conservée : $P_{\text{entrée}} = P_{\text{sortie}}$.
Comme $P = C \cdot \omega$ (Couple $\times$ Vitesse angulaire), si la vitesse $\omega$ diminue (réducteur), le couple $C$ \textbf{augmente} proportionnellement.
\[
C_{\text{sortie}} \approx R \times C_{\text{entrée}} \quad \text{(pour un réducteur)}
\]
\end{propriete}

\begin{attention}[Pièges fréquents]
\begin{itemize}
    \item Ne pas confondre $Z_{\text{menante}}$ et $Z_{\text{menée}}$ dans la formule de $R$. La roue menée est au numérateur.
    \item Ne pas oublier que le sens de rotation s'inverse à \emph{chaque} engrenage externe.
    \item « Plus de force » au niveau de la meule signifie « plus grand couple », pas « plus grande puissance ». La puissance moteur ne change pas (sauf pertes).
\end{itemize}
\end{attention}

\courssub{Corrigé commenté de l'activité}
\begin{exoresolu}[Analyse complète de la transmission du moulin de Nkolbisson]
\textbf{Énoncé complet :} Voir consignes ci-dessus.

\tcblower

\textbf{Solution détaillée et commentée :}

\textbf{1. Identification des organes}
\begin{itemize}
    \item \textbf{Roue menante :} C'est le \textbf{pignon de 24 dents}. Il est solidaire de l'arbre du \textbf{moteur} (source d'énergie). C'est lui qui « tire » l'autre roue.
    \item \textbf{Roue menée :} C'est la \textbf{grosse roue de 96 dents}. Elle est solidaire de l'arbre de la \textbf{meule} (organe récepteur). Elle « subit » l'entraînement du pignon.
\end{itemize}
\emph{Justification :} Par définition, la menante est du côté de la source (moteur), la menée du côté de l'utilisateur (meule).

\textbf{2. Sens de rotation}
\begin{itemize}
    \item Le moteur tourne dans le \textbf{sens horaire}.
    \item Le pignon (menant) est solidaire du moteur $\rightarrow$ Il tourne aussi dans le \textbf{sens horaire}.
    \item La grosse roue (menée) engrène \emph{extérieurement} avec le pignon.
    \item \textbf{Propriété :} Deux roues extérieures engrenées tournent en sens inverse.
    \item \textbf{Conclusion :} La grosse roue (et donc la meule) tourne dans le \textbf{sens anti-horaire} (trigonométrique).
\end{itemize}
\begin{center}
\textbf{\textcolor{popgreen}{$\leftarrow$ Sens anti-horaire (Meule)}} \quad $\leftarrow$ Engrenement \quad \textbf{\textcolor{poporange}{$\rightarrow$ Sens horaire (Moteur/Pignon)}}
\end{center}

\textbf{3. Rapport de transmission et vitesse de la meule}
\begin{itemize}
    \item \textbf{Données :} $Z_{\text{menante}} = Z_1 = 24$~dents ; $Z_{\text{menée}} = Z_2 = 96$~dents ; $N_{\text{menante}} = N_{\text{moteur}} = 1440$~tr/min.
    \item \textbf{Calcul du rapport $R$ :}
    \[
    R = \frac{Z_{\text{menée}}}{Z_{\text{menante}}} = \frac{96}{24} = \mathbf{4}
    \]
    $R$ est un nombre pur (sans unité). Cela signifie que pour \textbf{4 tours du moteur}, la meule fait \textbf{1 tour}.
    \item \textbf{Calcul de $N_{\text{meule}}$ :}
    On utilise la relation $R = \frac{N_{\text{menante}}}{N_{\text{menée}}}$.
    \[
    N_{\text{menée}} = N_{\text{meule}} = \frac{N_{\text{menante}}}{R} = \frac{1440}{4} = \mathbf{360~tr/min}
    \]
    \emph{Vérification :} $1440 / 360 = 4 = R$. C'est cohérent.
\end{itemize}

\textbf{4. Nature de la transmission}
\begin{itemize}
    \item $R = 4 > 1$ (ou $N_{\text{meule}} = 360 < N_{\text{moteur}} = 1440$).
    \item C'est un \textbf{RÉDUCTEUR de vitesse}.
    \item \textbf{Justification :} La vitesse de rotation en sortie (meule) est 4 fois plus petite qu'en entrée (moteur).
\end{itemize}

\textbf{5. Intérêt technologique (Couple et Broyage)}
\begin{itemize}
    \item \textbf{Analyse énergétique :} La puissance fournie par le moteur ($P_{\text{moteur}}$) est transmise à la meule ($P_{\text{meule}}$). En supposant un bon rendement (peu de frottements dans les paliers, graissage correct), $P_{\text{meule}} \approx P_{\text{moteur}}$.
    \item \textbf{Relation Puissance - Couple - Vitesse :} $P = C \times \omega$ (avec $\omega$ proportionnel à $N$).
    \item Comme $N$ a été \textbf{divisé par 4} ($R=4$), pour conserver $P$ constante, le couple $C$ a été \textbf{multiplié par 4}.
    \[
    C_{\text{meule}} \approx R \times C_{\text{moteur}} = 4 \times C_{\text{moteur}}
    \]
    \item \textbf{Interprétation pour M. Atangana :} « Plus de force » signifie scientifiquement \textbf{un couple moteur 4 fois plus grand} sur l'arbre de la meule.
    \item \textbf{Pourquoi c'est indispensable :} Moudre du maïs sec (grains durs) demande un \textbf{effort de cisaillement et de compression énorme}. Un couple élevé permet d'écraser les grains sans que le moteur ne cale (ne s'arrête par manque de couple). Si on branchait la meule directement sur le moteur (sans réducteur), la vitesse serait trop grande (1440 tr/min = danger, éclatement pierre) et le couple trop faible : le moteur bloquerait au premier grain dur.
\end{itemize}

\textbf{6. Consignes de sécurité (Affichage obligatoire sur le carter)}
\begin{enumerate}
    \item \textbf{INTERDIT D'OUVRIR LE CARTER OU D'INTRODUIRE LES MAINS/OUTILS PENDANT LE FONCTIONNEMENT.} Risque d'arrachage, de broyage des doigts ou de projection de débris (dents cassées, éclats de maïs). \emph{Arrêt total + temps d'arrêt complet avant toute intervention.}
    \item \textbf{PORT DES EPI OBLIGATOIRE : Lunettes de protection + Chaussures fermées + Pas de vêtements amples / cheveux attachés.} Risque d'entraînement (vêtements, cheveux) dans l'engrenage découvert ou projection de poussière de maïs dans les yeux.
\end{enumerate}
\emph{Autres consignes valides :} Vérifier le graissage hebdomadaire ; Ne pas dépasser le débit de maïs préconisé (surcharge moteur) ; Interrupteur d'urgence accessible.

\textbf{7. Prolongement : Remplacement par une roue de 48 dents}
\begin{itemize}
    \item \textbf{Nouvelles données :} $Z'_{\text{menée}} = 48$~dents ; $Z_{\text{menante}} = 24$~dents (inchangé).
    \item \textbf{Nouveau rapport $R'$ :}
    \[
    R' = \frac{Z'_{\text{menée}}}{Z_{\text{menante}}} = \frac{48}{24} = \mathbf{2}
    \]
    \item \textbf{Nouvelle vitesse $N'_{\text{meule}}$ :}
    \[
    N'_{\text{meule}} = \frac{N_{\text{moteur}}}{R'} = \frac{1440}{2} = \mathbf{720~tr/min}
    \]
    \item \textbf{Conséquence sur le couple (Force de broyage) :}
    Le nouveau rapport $R' = 2$ est \textbf{deux fois plus petit} que l'ancien ($R=4$).
    Le couple sur la meule sera : $C'_{\text{meule}} \approx 2 \times C_{\text{moteur}}$.
    \textbf{Comparaison :} $C'_{\text{meule}} = \frac{1}{2} C_{\text{meule (ancien)}}$.
    \textbf{Le couple est divisé par 2.} La meule a \textbf{deux fois moins de force} pour broyer.
    \item \textbf{Avis :} C'est une \textbf{MAUVAISE IDÉE} pour du maïs sec.
    \begin{itemize}
        \item La vitesse double (720 tr/min au lieu de 360 tr/min) $\rightarrow$ Risque d'échauffement de la pierre, de projection de farine, d'usure prématurée.
        \item Le couple est divisé par 2 $\rightarrow$ Le moteur risque de \textbf{caler} (s'arrêter brutalement) sur les grains durs, ce qui peut l'endommager (surchauffe du bobinage) et user l'embrayage ou les courroies s'il y en a.
    \end{itemize}
    \emph{Conclusion :} On garde la roue de 96 dents pour le maïs sec. Une roue de 48 dents pourrait convenir pour un produit plus tendre (ex: arachides déjà décortiquées) mais pas pour la farine de maïs villageoise.
\end{itemize}
\end{exoresolu}

\courssub{Bilan de la séance}
Cette séance t'a permis de \textbf{valider tes compétences} sur le thème des engrenages dans un contexte professionnel réaliste (maintenance rurale / agro-alimentaire) :

\begin{itemize}
    \item \textbf{Vocabulaire maîtrisé :} Tu sais identifier sans ambiguïté la roue menante (côté énergie) et la roue menée (côté utilisateur), même si la menée est plus grosse.
    \item \textbf{Kinématique :} Tu appliques la règle du \textbf{sens inverse} pour un engrenage externe et tu calcules correctement le \textbf{rapport de transmission} $R = Z_{\text{menée}}/Z_{\text{menante}}$ et la vitesse de sortie $N_{\text{sortie}} = N_{\text{entrée}}/R$.
    \item \textbf{Analyse énergétique (Niveau 3ème) :} Tu fais le lien crucial entre \textbf{vitesse} et \textbf{couple} via la conservation de la puissance ($P = C \omega$). Tu comprends qu'un \textbf{réducteur} ($R>1$) \textbf{augmente le couple} (la « force » de travail) au détriment de la vitesse. C'est le principe de base de la boîte de vitesse d'une moto-taxi (rapports bas pour démarrer en côte) ou du moulin.
    \item \textbf{Dimensionnement / Choix technique :} Le prolongement t'a montré qu'on ne choisit pas un engrenage au hasard : le rapport $R$ détermine le couple disponible. Pour broyer du dur (maïs, mil, manioc sec), il faut un \textbf{grand $R$} (grosse différence de dents). Pour aller vite (roues de vélo, multiplicateur de vélo), il faut un \textbf{petit $R$}.
    \item \textbf{Sécurité :} Tu as identifié les risques majeurs (entraînement, projection, coupure) et formulé des consignes claires, impératives et affichables, conformes à la réglementation du travail (Code du travail camerounais, sécurité des machines).
\end{itemize}

\begin{aretenir}[Ce qu'il faut retenir pour le BEPC]
\begin{itemize}
    \item \textbf{Engrenage externe :} Sens de rotation \textbf{inversés}. $R = \frac{Z_{\text{menée}}}{Z_{\text{menante}}} = \frac{N_{\text{menante}}}{N_{\text{menée}}}$.
    \item \textbf{Réducteur ($R>1$) :} Vitesse $\downarrow$, Couple $\uparrow$ (Force $\uparrow$). Utilisation : Moulin, treuil, boîte de vitesse (rapports bas), perceuse.
    \item \textbf{Multiplicateur ($R<1$) :} Vitesse $\uparrow$, Couple $\downarrow$. Utilisation : Vélo (grand plateau / petit pignon), dynamo, ventilateur.
    \item \textbf{Sécurité :} Carter fermé, arrêt avant maintenance, EPI (lunettes, chaussures, pas de vêtements amples), graissage régulier.
\end{itemize}
\end{aretenir}

\begin{savaistu}[Le moulin à maïs au Cameroun : un patrimoine technique]
Le moulin à meule de pierre (souvent en granit ou en pierre volcanique des monts Mandara ou de l'Adamaoua) est l'outil central de la transformation du maïs en \emph{farine de maïs} (pour le \emph{couscous}, le \emph{kwem}, le \emph{ndolé}). 
\begin{itemize}
    \item \textbf{Traditionnel :} Actionné à la main (meule dormante + meule courante) ou par force animale (âne, bœuf). Débit faible, pénibilité forte.
    \item \textbf{Motorisé (années 70-80 à aujourd'hui) :} Moteur thermique (essence/diesel) ou électrique (ENEO ou groupe électrogène). La transmission par \textbf{engrenages droits en fonte} (souvent récupérés sur des camions ou fabriqués localement) est la solution la plus robuste pour supporter les chocs des grains durs et la poussière.
    \item \textbf{Entretien critique :} Le \textbf{graissage} (graisse épaisse type « graisse à roulement » ou huile de vidange épaissie) est vital. Sans graisse, les dents s'usent en quelques semaines (jeu excessif, bruit, casse). L'alignement des arbres (parallélisme) évite l'usure en biseau.
    \item \textbf{Innovation locale :} Certains artisans (ex: à Douala Bassa ou Yaoundé Mvog-Ada) conçoivent des \textbf{multiplicateurs} pour adapter des moteurs rapides (3000 tr/min) sur des meules lentes, ou des \textbf{réducteurs à deux étages} (pignon $\rightarrow$ roue intermédiaire $\rightarrow$ grosse roue) pour obtenir des rapports $R > 10$ avec des moteurs de faible couple.
\end{itemize}
Comprendre les engrenages, c'est comprendre comment le Cameroun transforme ses récoltes !
\end{savaistu}

\newpage

\coursec{Méthodes et exercices résolus}

\coursec{Les engrenages}

\courssub{Découverte expérimentale : Manipulation d'un kit d'engrenages}
\begin{experience}[Observation du mouvement et de la transmission par engrenages]
    \textbf{Matériel :} Kit d'engrenages pédagogiques (pignons, couronnes, crémaillère, vis sans fin, arbres, support), crayon, règle.
    \textbf{Protocole :}
    \begin{enumerate}
        \item \textbf{Montage 1 (Engrenage droit extérieur) :} Emboîter un pignon $Z_1=12$ dents et une couronne $Z_2=24$ dents sur deux axes parallèles. Tourner le pignon (menant) d'un tour complet dans le sens horaire. Compter le nombre de tours de la couronne (menée) et noter son sens de rotation.
        \item \textbf{Montage 2 (Engrenage intérieur) :} Remplacer la couronne par une couronne intérieure (dents à l'intérieur) de même module. Recommencer l'observation.
        \item \textbf{Montage 3 (Vis sans fin) :} Monter une vis sans fin (filet unique) et sa roue hélicoïdale (axes croisés à 90°). Tourner la vis d'un tour. Observer le déplacement de la roue (nombre de pas/dents franchis) et le sens. Essayer de tourner la roue pour faire tourner la vis (test d'irréversibilité).
        \item \textbf{Montage 4 (Crémaillère) :} Fixer une crémaillère. Faire rouler un pignon dessus. Observer la transformation rotation $\leftrightarrow$ translation.
    \end{enumerate}
    \textbf{Observations attendues :}
    \begin{itemize}
        \item Montage 1 : Couronne tourne 1/2 tour, sens \textbf{inverse}. Contact par roulement sans glissement (théorique).
        \item Montage 2 : Couronne intérieure tourne 1/2 tour, \textbf{même sens}.
        \item Montage 3 : Roue avance d'1 dent (1/ $Z_{roue}$ tour). Sens selon hélice. \textbf{Impossible} de faire tourner la vis en actionnant la roue (irréversibilité statique).
        \item Montage 4 : Rotation du pignon $\rightarrow$ Translation rectiligne de la crémaillère.
    \end{itemize}
    \textbf{Interprétation :} Les engrenages transmettent un mouvement de rotation entre axes (parallèles, concourants, croisés) ou transforment rotation/translation. Le \textbf{rapport de transmission} dépend du nombre de dents. Le \textbf{type de denture} (droite, hélicoïdale, vis) détermine le sens, le bruit, la charge et la réversibilité.
    \textbf{Conclusion :} On définit le \cle{module} $m$, le \cle{pas} $p$, le \cle{cercle primitif} et le \cle{rapport de transmission} $R$ pour quantifier ces transmissions.
\end{experience}

\courssub{Définitions et vocabulaire}
\begin{definition}[Engrenage]
    Un \cle{engrenage} est un organe mécanique denté qui transmet un mouvement de rotation (et une puissance) entre deux arbres, par engrènement successif des dents, \textbf{sans glissement} (théorique) au niveau des cercles primitifs.
\end{definition}

\begin{definition}[Module, Pas, Cercles de référence]
    \begin{itemize}
        \item \cle{Module} $m$ (en mm) : Quotient du pas par $\pi$. C'est la cote fondamentale normalisée (ISO 54). Série courante : 1 ; 1,25 ; 1,5 ; 2 ; 2,5 ; 3 ; 4 ; 5 ; 6 ; 8 ; 10 mm.
        \item \cle{Pas} $p$ (en mm) : Longueur d'arc sur le cercle primitif entre deux profils de dents correspondants. $p = \pi \times m$.
        \item \cle{Cercle primitif} (diamètre $D$) : Cercle théorique de contact pur (roulement sans glissement). $D = m \times Z$.
        \item \cle{Cercle de tête} (diamètre $D_a$) : Sommet des dents. $D_a = D + 2h_a \approx D + 2m$ (denture normale).
        \item \cle{Cercle de pied} (diamètre $D_f$) : Fond des creux. $D_f = D - 2h_f \approx D - 2,5m$.
        \item \cle{Jeu de flanc} $j$ : Écart angulaire/linéaire entre les flancs d'une dent menante et la dent menée. Indispensable pour éviter le grippage et permettre la lubrification.
    \end{itemize}
\end{definition}

\begin{propriete}[Rapport de transmission $R$ (Convention Menant / Menée)]
    Pour un engrenage idéal (rendement $\eta = 1$), le \cle{rapport de transmission} $R$ (sans unité) est défini par les égalités suivantes :
    \[
    R = \frac{\omega_{\text{menant}}}{\omega_{\text{menée}}} = \frac{N_{\text{menant}}}{N_{\text{menée}}} = \frac{Z_{\text{menée}}}{Z_{\text{menant}}} = \frac{D_{\text{menée}}}{D_{\text{menant}}} = \frac{C_{\text{menée}}}{C_{\text{menant}}}
    \]
    \begin{itemize}
        \item Indice \textbf{1 (Menant / Entrée / Moteur)} : Celui qui \textbf{entraîne}.
        \item Indice \textbf{2 (Menée / Sortie / Utilisateur)} : Celui qui \textbf{reçoit} le mouvement.
    \end{itemize}
    \textbf{Conséquences :}
    \begin{itemize}
        \item $R > 1$ : \textbf{Réducteur} (Vitesse $\downarrow$, Couple $\uparrow$). Ex: Réducteur de moto, boîte de vitesse 1ère.
        \item $R < 1$ : \textbf{Multiplicateur} (Overdrive) (Vitesse $\uparrow$, Couple $\downarrow$). Ex: 4ème/5ème vitesse, alternateur.
        \item $R = 1$ : \textbf{Renvoi d'angle} (Vitesse =, Couple =). Ex: Différentiel, boîte de direction (parfois).
    \end{itemize}
    \textbf{Réel (avec rendement $\eta < 1$) :} $C_2 = \eta \cdot R \cdot C_1$ \quad et \quad $P_2 = \eta \cdot P_1$.
\end{propriete}

\courssub{Types d'engrenages et sens de rotation}
\begin{aretenir}[Classification selon la géométrie des axes]
    \begin{center}
        \begin{tabularx}{\linewidth}{|l|X|X|X|}\hline
            \rowcolor{popblueL} \textbf{Type} & \textbf{Axes} & \textbf{Denture} & \textbf{Sens rotation (Ext/Int)} \\ \hline
            \textbf{Engrenage droit} & Parallèles & Dents $\parallel$ axe & Ext: Inverse \quad Int: Même \\ \hline
            \textbf{Engrenage hélicoïdal} & Parallèles & Dents inclinées (hélice) & Ext: Inverse \quad Int: Même \\ \hline
            \textbf{Engrenage conique} & Concourants (souvent 90°) & Dents sur cône & Ext: Inverse \\ \hline
            \textbf{Vis sans fin} & Croisés (souvent 90°) & Vis (filet) + Roue hélicoïdale & Règle main droite (vis) \\ \hline
            \textbf{Crémaillère} & Axe $\perp$ Plan translation & Droite (sur plan) & Rotation $\leftrightarrow$ Translation \\ \hline
        \end{tabularx}
    \end{center}
    \textbf{Règle du sens (Engrenage extérieur) :} \textbf{Inversion} à chaque maille. \textbf{Roue folle (idler) :} Inverse le sens, \textbf{ne change pas le rapport $R$} (ses dents s'annulent : $Z_{folle}/Z_{menant} \times Z_{menée}/Z_{folle} = Z_{menée}/Z_{menant}$).
\end{aretenir}

\begin{savaistu}[L'engrenage au Cameroun : Du moto-taxi à la SONARA]
    \begin{itemize}
        \item \textbf{Moto-taxi (Bensikin) :} La transmission finale (sortie boîte $\rightarrow$ roue arrière) utilise un \textbf{pignon menant} (petit, acier traité, $Z \approx 14-16$) et une \textbf{couronne menée} (grande, $Z \approx 40-45$), souvent à \textbf{denture hélicoïdale} pour le silence et la résistance aux chocs des nids-de-poule. L'entraxe est réglable (tendeur de chaîne) mais le module doit être identique ($m \approx 4-5$ mm).
        \item \textbf{SONARA (Raffinerie de Limbé) :} Les pompes centrifuges et compresseurs sont entraînés par des \textbf{réducteurs à engrenages hélicoïdaux} (parallèles) ou \textbf{coniques} (angle droit) de forte puissance (MW). L'alignement laser des axes et la qualité de l'huile (viscosité ISO VG 220/320, additifs EP) sont critiques pour la durée de vie (20 000 h visées).
        \item \textbf{Histoire :} La machine d'Anticythère (Grèce, ~100 av. J.-C.) utilisait déjà des engrenages à dents triangulaires pour calculer les positions astronomiques. Léonard de Vinci a dessiné les premiers engrenages hélicoïdaux et coniques.
    \end{itemize}
\end{savaistu}

\courssub{Avantages et inconvénients des engrenages}
\begin{aretenir}[Choix technologique : Tableau de synthèse]
    \begin{center}
        \begin{tabularx}{\linewidth}{|l|X|X|}\hline
            \rowcolor{popblueL} \textbf{Type} & \textbf{Avantages clés} & \textbf{Inconvénients clés} \\ \hline
            \textbf{Droit} & Simple, peu cher, \textbf{pas de poussée axiale}, réversible, rendement élevé (0,98-0,99) & Bruit à haute vitesse, chocs à l'engrènement (contact ligne brutal) \\ \hline
            \textbf{Hélicoïdal} & \textbf{Silencieux} (engrènement progressif), charge élevée (contact surface), réversible & \textbf{Poussée axiale} forte (butées à rouleaux obligatoires), plus cher à usiner \\ \hline
            \textbf{Conique} & Transmission d'angle (90° standard), compact, réversible & Complexe à usiner/rectifier, réglage précis du cône (contact) requis \\ \hline
            \textbf{Vis sans fin} & \textbf{Très grand rapport} en 1 étage (1 tour vis = 1 dent roue), \textbf{Irréversible} (sécurité/freinage), axes croisés compacts & Rendement faible (0,4-0,7, glissement), \textbf{échauffement}, jeu axial, usure rapide (bronze/acier) \\ \hline
            \textbf{Crémaillère} & Transformation \textbf{Rotation $\leftrightarrow$ Translation}, course illimitée (théorique) & Course limitée par longueur, usure/précision guidage, bruit \\ \hline
        \end{tabularx}
    \end{center}
\end{aretenir}

\courssub{Rapport de transmission : calculs et applications}
\begin{methode}[Calculer le rapport de transmission (Vitesse, Couple, Dents)]
    \textbf{Objectif :} Utiliser les formules fondamentales pour dimensionner ou vérifier une transmission.
    \begin{enumerate}
        \item \textbf{Identifier les grandeurs connues} : $N_1, N_2$ (tr/min ou rad/s), $Z_1, Z_2$, $C_1, C_2$ (N$\cdot$m), $D_1, D_2$ (m ou mm).
        \item \textbf{Choisir la formule} (Menant = 1, Menée = 2) :
              \[
              R = \frac{N_1}{N_2} = \frac{Z_2}{Z_1} = \frac{D_2}{D_1} = \frac{C_2}{C_1} \quad (\eta=1)
              \]
        \item \textbf{Calculer l'inconnue} (isoler la variable).
        \item \textbf{Vérifier l'ordre de grandeur} : Réducteur ($N_2 < N_1$, $C_2 > C_1$) ou Multiplicateur ($N_2 > N_1$, $C_2 < C_1$).
        \item \textbf{Si rendement $\eta$ donné} : $C_2 = \eta \cdot R \cdot C_1$ ; $P_2 = \eta \cdot P_1$.
    \end{enumerate}
    \textbf{Attention :} $R$ sans unité. Au BEPC, on attend l'égalité $\frac{N_1}{N_2} = \frac{Z_2}{Z_1}$.
\end{methode}

\begin{methode}[Calculer le rapport global d'un train d'engrenages à plusieurs étages]
    \textbf{Objectif :} Déterminer la vitesse et le couple de sortie d'une boîte de vitesse à plusieurs paires en série.
    \begin{enumerate}
        \item \textbf{Décomposer} la chaîne cinématique en étages (paires menant/mené successives).
        \item \textbf{Calculer le rapport partiel $R_i$} de chaque étage : $R_i = \frac{Z_{\text{menée},i}}{Z_{\text{menant},i}}$.
        \item \textbf{Calculer le rapport global $R_{\text{tot}}$} : Produit des rapports partiels.
              \[
              R_{\text{tot}} = R_1 \times R_2 \times \dots \times R_k = \frac{N_{\text{entrée}}}{N_{\text{sortie}}}
              \]
              \emph{Rappel :} Roues solidaires d'un même arbre $\rightarrow$ même $N$. Roue folle (libre sur arbre) $\rightarrow$ ne change pas $R$ (inverse seulement le sens).
        \item \textbf{Calculer la sortie} : $N_{\text{sortie}} = \frac{N_{\text{entrée}}}{R_{\text{tot}}}$, $C_{\text{sortie}} = \eta_{\text{tot}} \cdot R_{\text{tot}} \cdot C_{\text{entrée}}$.
    \end{enumerate}
\end{methode}

\courssub{Sécurité et entretien des machines à engrenages}
\begin{methode}[Analyser la sécurité et l'entretien d'une transmission par engrenages]
    \textbf{Objectif :} Identifier les risques, les protections obligatoires et les opérations de maintenance préventive.
    \begin{enumerate}
        \item \textbf{Risques majeurs (Identification) :}
              \begin{itemize}
                  \item \textbf{Prise / Cisaillement} : Zone d'engrènement (doigts, vêtements, cheveux). \textbf{Zone interdite d'accès en mouvement.}
                  \item \textbf{Projection} : Cassure de dent, éclat de lubrifiant, corps étranger.
                  \item \textbf{Bruit} : Niveau sonore > 85 dB(A) $\rightarrow$ Risque surdité (port EPI bouchons/casque).
                  \item \textbf{Chaleur} : Carter chaud (brûlure), surchauffe huile (risque incendie si fuite).
              \end{itemize}
        \item \textbf{Protections collectives (Obligatoires Code du Travail / Normes CM) :}
              \begin{itemize}
                  \item \textbf{Carter fermé} (étanche à l'huile, empêche tout accès aux dents).
                  \item \textbf{Verrouillage / Consignation (LOTO)} : Impossible d'ouvrir le carter si machine en marche (interrupteur de sécurité sur capot).
                  \item \textbf{Pictogrammes} : "Danger pièces mobiles", "Ne pas ouvrir en marche", "Port EPI obligatoire".
              \end{itemize}
        \item \textbf{Entretien préventif (Gamme type) :}
              \begin{itemize}
                  \item \textbf{Niveau / Qualité huile} : Vérification journalière (jauge / hublot). Vidange périodique (heures de fonctionnement, ex: 2000 h ou 6 mois).
                  \item \textbf{État des dents} : Contrôle visuel (usure, piqûres/pitting, cassure) à chaque vidange.
                  \item \textbf{Alignement / Jeu} : Contrôle du jeu de flanc (calepinage) et alignement axes (vibrations, bruit anormal).
                  \item \textbf{Respiration / Filtre} : Vérifier le reniflard (évite surpression / aspiration poussière/eau).
              \end{itemize}
        \item \textbf{Consignation (Lockout/Tagout) - Procédure stricte avant intervention :}
              Arrêt machine $\rightarrow$ Coupure source énergie (électrique, hydraulique) $\rightarrow$ \textbf{Verrouillage} (cadenas personnel) $\rightarrow$ \textbf{Étiquetage} (nom, date, motif) $\rightarrow$ Vérification absence tension/rotation (essai démarrage) $\rightarrow$ Intervention.
    \end{enumerate}
\end{methode}

\courssub{Méthodes et exercices résolus}

\begin{methode}[Identifier les éléments constitutifs d'un engrenage]
    \textbf{Objectif :} Reconnaître et nommer les pièces principales sur un schéma ou une pièce réelle.
    \begin{enumerate}
        \item \textbf{Observer} la pièce : repérer les \cle{dents} (parties saillantes), les \cle{creux} (espaces entre les dents), le \cle{jeu} (espace entre dent et creux de l'engrenage mené).
        \item \textbf{Repérer les cercles de référence} (souvent en traits pointillés sur les schémas) :
              \begin{itemize}
                  \item Cercle \cle{primitif} (cercle de référence théorique où s'engrènent les dents, rayon $R$). Trait \textbf{mixte (pointillé-tiret)}.
                  \item Cercle de \cle{tête} (sommet des dents, rayon $R_a = R + h_a$). Trait \textbf{discontinu (pointillés)}.
                  \item Cercle de \cle{pied} (fond des creux, rayon $R_f = R - h_f$). Trait \textbf{pointillé serré}.
              \end{itemize}
        \item \textbf{Noter les cotes caractéristiques} : le \cle{module} $m$ (en mm), le \cle{pas} $p = \pi \cdot m$, le nombre de dents $Z$.
        \item \textbf{Calculer le diamètre primitif} : $D = m \times Z$ (relation fondamentale).
    \end{enumerate}
    \textbf{Réflexe :} Sur un dessin technique, le cercle primitif est tracé en trait mixte (pointillé-tiret). C'est lui qui sert aux calculs de transmission.
\end{methode}

\begin{exoresolu}[Lecture de schéma et calcul de dimensions d'un pignon de moto-taxi]
    Sur le schéma du pignon d'attaque d'une moto-taxi (type Yamaha 125), on relève : module $m = \SI{2,5}{mm}$, nombre de dents $Z_1 = 14$. Le pignon entraîne une couronne de $Z_2 = 42$ dents de même module.
    \begin{enumerate}
        \item Nomme les cercles (1), (2), (3) sur la figure ci-contre.
        \item Calcule le diamètre primitif $D_1$ du pignon et $D_2$ de la couronne.
        \item Calcule le pas $p$ et la distance d'axes $a$ (approchée) pour un engrenage droit à denture droite.
    \end{enumerate}
    \begin{popfigure}
        \begin{tikzpicture}[scale=0.8, >=Stealth]
            % Pignon : R1 = 1.75cm (35mm), Ra1 = 2.0cm, Rf1 = 1.5cm
            % Couronne : R2 = 5.25cm (105mm), Ra2 = 5.5cm, Rf2 = 5.0cm
            % Entraxe a = 7.0cm (70mm)
            \draw[popblue, thick] (0,0) circle (1.75); % Primitif
            \draw[popblue, dashed] (0,0) circle (2.0); % Tête
            \draw[popblue, densely dotted] (0,0) circle (1.5); % Pied
            \node[popblue] at (0, -2.5) {Pignon $Z_1=14$};
            % Couronne
            \draw[poporange, thick] (7,0) circle (5.25); % Primitif (Centre corrigé à 7)
            \draw[poporange, dashed] (7,0) circle (5.5); % Tête
            \draw[poporange, densely dotted] (7,0) circle (5.0); % Pied
            \node[poporange] at (7, -6.0) {Couronne $Z_2=42$};
            % Flèches cotations
            \draw[<->, popmuted] (0,1.75) -- (0,2.0) node[midway, right] {$h_a$};
            \draw[<->, popmuted] (0,1.5) -- (0,1.75) node[midway, left] {$h_f$};
            \draw[<->, popdark] (-1.75,0) -- (1.75,0) node[midway, above] {$D_1$};
            \draw[<->, popdark] (0, -1.0) -- (7, -1.0) node[midway, above] {$a$}; % Entraxe centre à centre
            % Dents symboliques
            \foreach \angle in {0, 25.714, ..., 360} {
                \draw[popblue, thick] (\angle:1.75) -- (\angle:2.0);
                \draw[popblue, thick] (\angle:1.5) -- (\angle:1.75);
            }
            \foreach \angle in {0, 8.571, ..., 360} {
                \draw[poporange, thick] (7,0) ++(\angle:5.25) -- ++(\angle:0.25);
                \draw[poporange, thick] (7,0) ++(\angle:5.0) -- ++(\angle:0.25);
            }
        \end{tikzpicture}
        \legende{Schéma simplifié d'un engrenage droit : pignon d'attaque et couronne menée (vue de face partielle). Cercles primitifs en trait mixte.}
    \end{popfigure}
    \tcblower
    \textbf{Solution détaillée}
    \begin{enumerate}
        \item \textbf{Nommage des cercles (norme ISO / dessin technique) :}
              \begin{itemize}
                  \item Cercle (1) : Trait mixte (pointillé-tiret) $\rightarrow$ \textbf{Cercle primitif} (diamètre $D$).
                  \item Cercle (2) : Trait discontinu (pointillés) $\rightarrow$ \textbf{Cercle de tête} (diamètre $D_a$).
                  \item Cercle (3) : Trait pointillé serré $\rightarrow$ \textbf{Cercle de pied} (diamètre $D_f$).
              \end{itemize}
        \item \textbf{Calcul des diamètres primitifs} (Formule : $D = m \times Z$)
              \begin{align*}
                  D_1 &= m \times Z_1 = 2,5 \times 14 = \SI{35,0}{mm} \\
                  D_2 &= m \times Z_2 = 2,5 \times 42 = \SI{105,0}{mm}
              \end{align*}
        \item \textbf{Calcul du pas $p$ et de la distance d'axes $a$}
              \begin{itemize}
                  \item Pas : $p = \pi \times m = \pi \times 2,5 \approx \SI{7,85}{mm}$.
                  \item Distance d'axes (engrenage extérieur) : $a = \frac{D_1 + D_2}{2} = \frac{35 + 105}{2} = \SI{70,0}{mm}$.
              \end{itemize}
    \end{enumerate}
    \textbf{Conclusion :} Le pignon a un diamètre primitif de \SI{35}{mm}, la couronne de \SI{105}{mm}. L'entraxe à réaliser sur le carter moteur est de \SI{70}{mm}. Le pas de \SI{7,85}{mm} correspond à l'arc de cercle primitif entre deux dents successives.
\end{exoresolu}

\begin{methode}[Déterminer le sens de rotation dans un train d'engrenages]
    \textbf{Objectif :} Prévoir le sens de rotation de la sortie connaissant l'entrée et le nombre d'engrenages.
    \begin{enumerate}
        \item \textbf{Identifier le type d'engrenage} : \textbf{Extérieur} (dents à l'extérieur, le plus courant) ou \textbf{Intérieur} (dents à l'intérieur, couronne creuse).
        \item \textbf{Appliquer la règle du sens} :
              \begin{itemize}
                  \item \textbf{Engrenage extérieur :} Les deux roues tournent en \textbf{sens inverse} (si l'une tourne horaire, l'autre anti-horaire).
                  \item \textbf{Engrenage intérieur :} Les deux roues tournent dans le \textbf{même sens}.
                  \item \textbf{Vis sans fin / Roue hélicoïdale croisée :} Règle de la main droite (vis) ou observation de l'hélice.
              \end{itemize}
        \item \textbf{Propager le sens} roue par roue depuis l'entrée (moteur) vers la sortie (roue, arbre de sortie).
        \item \textbf{Compter le nombre d'inversions} : Pair $\rightarrow$ même sens que l'entrée ; Impair $\rightarrow$ sens inverse.
    \end{enumerate}
    \textbf{Astuce :} Dessine une petite flèche sur la première roue et retourne-la à chaque engrenage extérieur.
\end{methode}

\begin{exoresolu}[Sens de rotation dans la boîte de vitesse d'un groupe électrogène]
    Le moteur thermique d'un groupe électrogène (arbre d'entrée) tourne dans le sens \textbf{horaire} (vue de face). Il entraîne un arbre intermédiaire via un engrenage extérieur $E_1$/$E_2$. Cet arbre intermédiaire porte un deuxième pignon $E_3$ qui engrène extérieurement avec la roue de sortie $E_4$ (alternateur). Un troisième engrenage extérieur $E_5$ (roue folle / \emph{idler}) est inséré entre $E_3$ et $E_4$ pour rapprocher les axes sans changer le rapport.
    \begin{enumerate}
        \item Donne le sens de rotation de l'alternateur ($E_4$) par rapport au moteur.
        \item À quoi sert la roue folle $E_5$ sur le plan du sens de rotation ? Sur le plan du rapport de transmission ?
    \end{enumerate}
    \tcblower
    \textbf{Solution détaillée}
    \begin{enumerate}
        \item \textbf{Propagation des sens (Engrenages extérieurs = inversion à chaque maille) :}
              \begin{itemize}
                  \item Moteur (Entrée) : \textbf{Horaire}.
                  \item $E_1$ (solidaire moteur) : Horaire.
                  \item $E_2$ (engrène sur $E_1$, extérieur) : \textbf{Anti-horaire} (1ère inversion).
                  \item $E_3$ (solidaire $E_2$ sur même arbre) : \textbf{Anti-horaire}.
                  \item $E_5$ (roue folle, engrène sur $E_3$, extérieur) : \textbf{Horaire} (2ème inversion).
                  \item $E_4$ (Sortie alternateur, engrène sur $E_5$, extérieur) : \textbf{Anti-horaire} (3ème inversion).
              \end{itemize}
              \textbf{Résultat :} L'alternateur tourne en \textbf{sens anti-horaire} (sens inverse du moteur car 3 inversions = nombre impair).
        \item \textbf{Rôle de la roue folle $E_5$ :}
              \begin{itemize}
                  \item \textbf{Sens de rotation :} Elle ajoute une inversion. Sans elle ($E_3$ engrenant directement $E_4$), il y aurait 2 inversions $\rightarrow$ même sens que le moteur. Avec elle, 3 inversions $\rightarrow$ sens inverse. \textbf{Elle permet de choisir le sens de sortie.}
                  \item \textbf{Rapport de transmission :} Son nombre de dents $Z_5$ ne figure \textbf{pas} dans le calcul du rapport global ($R = \frac{Z_2}{Z_1} \times \frac{Z_4}{Z_3}$). Elle ne change \textbf{pas} le rapport de vitesse ni de couple. Elle ne sert qu'à \textbf{combler l'entraxe} (distance d'axes) ou inverser le sens.
              \end{itemize}
    \end{enumerate}
    \textbf{Conclusion :} L'alternateur tourne en sens inverse du moteur. La roue folle est un artifice technologique pour adapter la géométrie (entraxe) ou la cinématique (sens) sans modifier le rapport de réduction.
\end{exoresolu}

\begin{exoresolu}[Dimensionnement du réducteur d'une pompe à eau manuelle]
    Une pompe à eau de forage (type India Mark II) est actionnée par un volant manuel. L'opérateur tourne le volant à $N_1 = \SI{60}{tr/min}$. Le mécanisme interne (bielle-manivelle) nécessite une vitesse de rotation lente sur l'arbre de pompage $N_2 = \SI{15}{tr/min}$ pour un débit correct. Le couple résistant sur l'arbre de pompage est estimé à $C_2 = \SI{120}{N.m}$. Le rendement global de la transmission par engrenages droits est $\eta = 0,85$ (85\%).
    \begin{enumerate}
        \item Détermine le rapport de transmission $R$ nécessaire.
        \item Si le pignon d'attaque (menant) a $Z_1 = 12$ dents, combien de dents $Z_2$ doit avoir la roue menée ?
        \item Calcule le couple $C_1$ que l'opérateur doit appliquer sur le volant (arbre entrant).
    \end{enumerate}
    \tcblower
    \textbf{Solution détaillée}
    \begin{enumerate}
        \item \textbf{Rapport de transmission $R$ (définition vitesse) :}
              \[
              R = \frac{N_1}{N_2} = \frac{60}{15} = 4
              \]
              C'est un \textbf{réducteur} de rapport 4 ($N_1 = 4 N_2$). La vitesse est divisée par 4.
        \item \textbf{Nombre de dents $Z_2$ (relation géométrique) :}
              \[
              R = \frac{Z_2}{Z_1} \implies Z_2 = R \times Z_1 = 4 \times 12 = 48 \text{ dents}
              \]
              \emph{Vérification :} $Z_2 > Z_1$, cohérent pour un réducteur (roue menée plus grosse).
        \item \textbf{Couple d'entrée $C_1$ (relation énergétique avec rendement) :}
              \begin{itemize}
                  \item Puissance sortie : $P_2 = C_2 \cdot \omega_2$. Puissance entrée : $P_1 = C_1 \cdot \omega_1$.
                  \item Rendement : $\eta = \frac{P_2}{P_1} = \frac{C_2 \omega_2}{C_1 \omega_1} = \frac{C_2}{C_1} \cdot \frac{1}{R}$.
                  \item On en déduit : $C_1 = \frac{C_2}{\eta \cdot R}$.
              \end{itemize}
              \[
              C_1 = \frac{120}{0,85 \times 4} = \frac{120}{3,4} \approx \SI{35,3}{N.m}
              \]
              \emph{Sans frottements ($\eta=1$), $C_1 = 30$ N$\cdot$m. Les frottements demandent un effort supplémentaire.}
    \end{enumerate}
    \textbf{Conclusion :} Il faut une couronne de 48 dents. L'opérateur doit fournir un couple d'environ \SI{35,3}{N.m} sur le volant. Ce couple est accessible pour un adulte (force d'environ 35 N sur un volant de 1 m de rayon, ou 70 N sur 0,5 m).
\end{exoresolu}

\begin{exoresolu}[Choix technologique : Direction de voiture vs Réducteur de moto vs Treuil de grue]
    \begin{enumerate}
        \item Sur une voiture particulière, la direction transforme la rotation du volant (axe horizontal) en rotation des fusées d'essieu (axes verticaux). Le système doit être \textbf{réversible} (roues $\rightarrow$ volant pour ressentir la route). Quel type d'engrenage est utilisé dans la \textbf{boîte de direction à vis et écrou} (ancienne technologie) ou \textbf{pignon-crémaillère} (moderne) ? Justifie.
        \item Sur une moto, le réducteur primaire (sortie boîte de vitesse $\rightarrow$ couronne arrière) transmet une forte puissance à haute vitesse entre axes parallèles. Le silence de fonctionnement est recherché. Quel type de denture choisir ? Justifie.
        \item Sur un treuil de levage de charge lourde (grue de chantier), on doit empêcher la charge de redescendre si le moteur coupe (sécurité). Les axes sont perpendiculaires et non concourants. Quel type choisir ? Justifie.
    \end{enumerate}
    \tcblower
    \textbf{Solution détaillée}
    \begin{enumerate}
        \item \textbf{Direction voiture :}
              \begin{itemize}
                  \item \textbf{Pignon-Crémaillère (moderne) :} Axe volant (rotation) $\perp$ Axe crémaillère (translation). \textbf{Avantage :} Réversible (info route $\rightarrow$ volant), précis, compact, peu d'entretien.
                  \item \textbf{Vis à billes / Vis sans fin à rouleaux (ancien "vis et écrou") :} Axes croisés. \textbf{Avantage :} Réversible si angle d'hélice élevé (vis à billes), démultiplication intégrée.
              \end{itemize}
              \textbf{Choix retenu aujourd'hui : Pignon-Crémaillère.} Justification : Réversibilité obligatoire, encombrement minimal, coût maîtrisé.
        \item \textbf{Réducteur moto (sortie boîte / couronne) :}
              \begin{itemize}
                  \item Axes parallèles. Puissance élevée, vitesse élevée.
                  \item \textbf{Engrenage à denture hélicoïdale.} Justification : Engrènement progressif $\rightarrow$ \textbf{Silencieux} (confort, réglementation bruit), surface de contact plus grande $\rightarrow$ \textbf{Résistance} accrue aux charges de choc. Inconvénient (poussée axiale) géré par des butées à rouleaux dans le carter.
              \end{itemize}
        \item \textbf{Treuil de grue (sécurité anti-retour) :}
              \begin{itemize}
                  \item Axes croisés (90°), non concourants. Exigence : \textbf{Irréversibilité} (charge ne redescend pas seule).
                  \item \textbf{Vis sans fin (simple filet, acier / bronze).} Justification : Seule technologie standard offrant l'\textbf{irréversibilité statique} naturelle (si angle d'hélice $<$ angle de frottement). Rapport de réduction très grand en un seul étage (compact). Inconvénient (rendement $\approx 0,4-0,6$, échauffement) accepté car fonctionnement intermittent.
              \end{itemize}
    \end{enumerate}
    \textbf{Conclusion :} Le choix dicte la géométrie (axes), la cinématique (réversible/irréversible), l'acoustique et le coût. Il n'y a pas de "meilleur" engrenage universel, mais le \textbf{mieux adapté} au cahier des charges.
\end{exoresolu}

\begin{methode}[Calculer le module et vérifier la résistance (approche initiation - Formule de Lewis simplifiée)]
    \textbf{Objectif :} Déterminer le module normalisé à partir de l'effort tangent ou vérifier la tenue en flexion.
    \begin{enumerate}
        \item \textbf{Calculer le couple $C$} sur l'arbre concerné (entrée ou sortie).
        \item \textbf{Calculer l'effort tangent $F_t$} (force au cercle primitif) :
              \[
              F_t = \frac{2C}{D} = \frac{2C}{m Z}
              \]
              \textbf{Unités obligatoires :} $C$ en \textbf{N$\cdot$m}, $D$ en \textbf{m} $\rightarrow$ $F_t$ en \textbf{N}. (Ou $C$ en N$\cdot$mm, $D$ en mm).
        \item \textbf{Choisir le module $m$} (série normalisée : 1, 1,25, 1,5, 2, 2,5, 3, 4, 5, 6, 8, 10... mm) pour que la contrainte de flexion $\sigma_f$ reste admissible.
        \item \textbf{Formule de Lewis simplifiée (initiation) :}
              \[
              \sigma_f = \frac{F_t}{b \cdot m \cdot Y}
              \]
              \begin{itemize}
                  \item $b$ = largeur de dent (souvent $b \approx 8m$ à $12m$).
                  \item $Y$ = coefficient de forme de la dent (dépend de $Z$ et angle de pression, $\approx 0,3$ à $0,4$ pour $Z>20$).
              \end{itemize}
        \item \textbf{Vérifier :} $\sigma_f \le \sigma_{\text{adm}}$ (Acier trempé $\approx 600-1000$ MPa, Fonte $\approx 150$ MPa, Plastique $\approx 50$ MPa).
    \end{enumerate}
    \textbf{Au BEPC :} On se limite souvent au calcul de $m$ par la relation $D = mZ$ ou au choix dans un catalogue. La vérification de contrainte est en option "Technologie" approfondie.
\end{methode}

\begin{exoresolu}[Choix du module pour un engrenage de mélangeuse à béton]
    Le tambour d'une bétonnière (mélangeuse) de chantier est entraîné par un pignon $Z_1 = 18$ dents sur l'arbre du moteur électrique (vitesse $N_1 = \SI{1500}{tr/min}$, puissance $P = \SI{2,2}{kW}$). La couronne $Z_2 = 108$ dents est en fonte. L'entraxe imposé par le châssis est $a = \SI{315}{mm}$.
    \begin{enumerate}
        \item Vérifie que l'entraxe correspond à un module normalisé. Détermine $m$.
        \item Calcule le couple $C_1$ sur l'arbre moteur.
        \item Calcule l'effort tangent $F_t$ sur les dents du pignon.
        \item Sachant que la largeur de dent $b = 10m$ et le coefficient de forme $Y = 0,32$, calcule la contrainte de flexion $\sigma_f$ dans les dents du pignon (acier, $\sigma_{\text{adm}} = \SI{700}{MPa}$). Le choix est-il validé ?
    \end{enumerate}
    \tcblower
    \textbf{Solution détaillée}
    \begin{enumerate}
        \item \textbf{Module imposé par l'entraxe :}
              \[
              a = \frac{D_1 + D_2}{2} = \frac{m(Z_1 + Z_2)}{2} \implies m = \frac{2a}{Z_1 + Z_2} = \frac{2 \times 315}{18 + 108} = \frac{630}{126} = \SI{5,0}{mm}
              \]
              $m = 5$ mm est une valeur \textbf{normalisée} (série standard). \textbf{Validé.}
        \item \textbf{Couple moteur $C_1$ :}
              \[
              P = C_1 \cdot \omega_1 = C_1 \cdot \frac{2\pi N_1}{60} \implies C_1 = \frac{60 P}{2\pi N_1} = \frac{60 \times 2200}{2\pi \times 1500} \approx \SI{14,0}{N.m}
              \]
        \item \textbf{Effort tangent $F_t$ (sur pignon) :}
              Diamètre primitif pignon : $D_1 = m Z_1 = 5 \times 18 = \SI{90}{mm} = \SI{0,090}{m}$.
              \[
              F_t = \frac{2 C_1}{D_1} = \frac{2 \times 14,0}{0,090} \approx \SI{311}{N}
              \]
        \item \textbf{Contrainte de flexion $\sigma_f$ (Lewis) :}
              \begin{itemize}
                  \item $b = 10m = \SI{50}{mm}$.
                  \item $m = \SI{5}{mm}$.
                  \item $Y = 0,32$.
              \end{itemize}
              \[
              \sigma_f = \frac{F_t}{b \cdot m \cdot Y} = \frac{311}{50 \times 5 \times 0,32} = \frac{311}{80} \approx \SI{3,89}{MPa}
              \]
              \textbf{Attention unités :} $F_t$ en N, $b$ et $m$ en mm $\rightarrow$ $\sigma_f$ en N/mm$^2$ = MPa.
              Comparaison : $\sigma_f = \SI{3,9}{MPa} \ll \sigma_{\text{adm}} = \SI{700}{MPa}$.
              \textbf{Validé largement.} La denture est surdimensionnée pour la flexion (souvent le cas pour les gros modules à basse vitesse, l'usure de surface / pitting devient le critère limitant).
    \end{enumerate}
    \textbf{Conclusion :} Le module normalisé imposé par l'entraxe est $m = 5$ mm. La contrainte de flexion est très faible (\SI{3,9}{MPa}), la résistance mécanique est largement assurée.
\end{exoresolu}

\begin{exoresolu}[Diagnostic de panne sur un réducteur de convoyeur à bande (SONARA / Cimenterie)]
    Un réducteur à engrenages hélicoïdaux (axes parallèles) d'un convoyeur à bande transporteur de clinker (cimenterie) présente un bruit anormal (ronronnement grave croissant) et une élévation de température du carter (\SI{95}{\celsius} mesurée au laser, normale \SI{70}{\celsius}). L'opérateur a remis de l'huile neuve il y a 2 mois.
    \begin{enumerate}
        \item Propose 3 causes probables classées par ordre de probabilité.
        \item Pour chaque cause, indique l'action corrective et la mesure de sécurité à prendre \textbf{avant} d'ouvrir le carter.
        \item Pourquoi l'huile neuve n'a-t-elle pas résolu le problème ?
    \end{enumerate}
    \tcblower
    \textbf{Solution détaillée}
    \begin{enumerate}
        \item \textbf{Causes probables (Diagnostic) :}
              \begin{enumerate}
                  \item \textbf{Désalignement des axes / Mauvaise mise en contact (Cause n°1).} Les engrenages hélicoïdaux sont très sensibles à l'alignement. Un fléchissement du châssis ou un montage imparfait crée un contact en bord de dent $\rightarrow$ surcharge locale, bruit, chaleur.
                  \item \textbf{Usure / Pitting (Picotage) avancé sur les flancs (Cause n°2).} Charge cyclique lourde (clinker) $\rightarrow$ fatigue de contact $\rightarrow$ perte de profil $\rightarrow$ bruit "ronronnement" et rendement chute $\rightarrow$ chaleur.
                  \item \textbf{Défaillance roulements (paliers) (Cause n°3).} Bruit grave peut venir des roulements (billes/rouleaux) qui supportent les poussées axiales des hélicoïdaux. Si cage cassée ou manque graisse $\rightarrow$ chaleur transmise au carter.
              \end{enumerate}
        \item \textbf{Actions correctives \& Sécurité (Consignation LOTO) :}
              \begin{center}
                  \begin{tabularx}{\linewidth}{|l|X|X|}\hline
                      \rowcolor{popblueL} \textbf{Cause} & \textbf{Action Corrective} & \textbf{Sécurité AVANT ouverture (Consignation)} \\ \hline
                      Désalignement & Réalignement précis (pied à coulisse / laser), calage pieds moteur/réducteur. & \textbf{1. Arrêt d'urgence.} \textbf{2. Coupure disjoncteur + Cadenas (LOTO).} \textbf{3. Attente refroidissement (< 50°C).} \textbf{4. Vérif. absence rotation (clé sur arbre).} \\ \hline
                      Usure dents & Remplacement jeu pignon/couronne (par paire !). Contrôle dureté. & Même procédure + \textbf{Étiquetage "Ne pas démarrer"}. \\ \hline
                      Roulements & Remplacement roulements (par jeu), vérif. butées axiales. & Même procédure + Vidange huile chaude (risque brûlure). \\ \hline
                  \end{tabularx}
              \end{center}
        \item \textbf{Pourquoi l'huile neuve a échoué :} L'huile assure la lubrification (film) et le refroidissement. Elle \textbf{ne peut pas corriger} un défaut géométrique (alignement, usure de forme, roulement HS). Le bruit et la chaleur sont des \textbf{symptômes mécaniques}, pas un manque de lubrifiant (sauf si huile inadaptée : viscosité trop basse pour la charge, mais ici le problème est structurel).
    \end{enumerate}
    \textbf{Conclusion :} La maintenance curative impose le diagnostic de cause racine (alignement/usure). La sécurité (Consignation LOTO) est \textbf{non négociable} avant tout démontage de carter. Le remplacement se fait toujours \textbf{par jeu complet} (pignon + couronne) pour respecter le rodage.
\end{exoresolu}

\begin{exoresolu}[Analyse de la boîte de vitesse 4 vitesses d'une moto-taxi (type CG 125)]
    Schéma simplifié de la boîte de vitesse (arbre primaire $\rightarrow$ arbre secondaire). Arbre primaire (entrée embrayage) : $Z_{p1}=13, Z_{p2}=18, Z_{p3}=22, Z_{p4}=26$. Arbre secondaire (sortie) : $Z_{s1}=35, Z_{s2}=28, Z_{s3}=24, Z_{s4}=21$. Les paires en prise sont : 1ère : $p1/s1$, 2ème : $p2/s2$, 3ème : $p3/s3$, 4ème : $p4/s4$. Le moteur tourne à $N_m = \SI{8000}{tr/min}$ au régime max. Rendement boîte $\eta = 0,95$.
    \begin{enumerate}
        \item Calcule le rapport de boîte $R_b$ pour chaque vitesse.
        \item Calcule la vitesse de rotation de l'arbre de sortie (secondaire) $N_s$ pour chaque vitesse.
        \item Si le couple moteur max est $C_m = \SI{12}{N.m}$, calcule le couple à la sortie de boîte $C_s$ pour la 1ère et la 4ème.
        \item Commente l'évolution du couple et de la vitesse. À quoi correspond la 4ème ?
    \end{enumerate}
    \tcblower
    \textbf{Solution détaillée}
    \begin{enumerate}
        \item \textbf{Rapports de boîte $R_b = \frac{Z_{\text{secondaire}}}{Z_{\text{primaire}}}$ (Menant = Primaire, Menée = Secondaire) :}
              \begin{align*}
                  \text{1ère : } R_1 &= \frac{35}{13} \approx 2,69 \\
                  \text{2ème : } R_2 &= \frac{28}{18} \approx 1,56 \\
                  \text{3ème : } R_3 &= \frac{24}{22} \approx 1,09 \\
                  \text{4ème : } R_4 &= \frac{21}{26} \approx 0,808
              \end{align*}
        \item \textbf{Vitesse sortie $N_s = \frac{N_m}{R_b}$ (Calcul exact avec fractions) :}
              \begin{align*}
                  N_{s1} &= \frac{8000}{35/13} = 8000 \times \frac{13}{35} \approx \SI{2971}{tr/min} \\
                  N_{s2} &= \frac{8000}{28/18} = 8000 \times \frac{18}{28} \approx \SI{5143}{tr/min} \\
                  N_{s3} &= \frac{8000}{24/22} = 8000 \times \frac{22}{24} \approx \SI{7333}{tr/min} \\
                  N_{s4} &= \frac{8000}{21/26} = 8000 \times \frac{26}{21} \approx \SI{9905}{tr/min}
              \end{align*}
        \item \textbf{Couple sortie $C_s = \eta \cdot R_b \cdot C_m$ ($\eta=0,95$) :}
              \begin{align*}
                  C_{s1} &= 0,95 \times \frac{35}{13} \times 12 \approx \SI{30,7}{N.m} \quad \text{(Fort couple pour démarrer/grimper)} \\
                  C_{s4} &= 0,95 \times \frac{21}{26} \times 12 \approx \SI{9,2}{N.m} \quad \text{(Couple réduit, vitesse max)}
              \end{align*}
        \item \textbf{Commentaire :}
              \begin{itemize}
                  \item \textbf{1ère vitesse :} $R > 1$ (Réducteur). Vitesse divisée $\approx 2,7$, Couple multiplié $\approx 2,6$. \textbf{Démarrage, côte, charge.}
                  \item \textbf{4ème vitesse :} $R < 1$ (Multiplicateur / \textbf{Overdrive}). Vitesse multipliée $\approx 1,24$, Couple divisé $\approx 0,77$. \textbf{Vitesse de pointe, route plate, économie carburant.}
                  \item Le passage des vitesses adapte le couple/vitesse au besoin (résistance au roulement, pente).
              \end{itemize}
    \end{enumerate}
    \textbf{Conclusion :} La boîte de vitesse offre une plage de rapport de $0,81$ à $2,69$. La 4ème est un overdrive (multiplicateur) permettant au moteur de tourner moins vite que la roue pour une vitesse de croisière économique.
\end{exoresolu}

\begin{attention}[Les 5 erreurs qui coûtent des points au BEPC]
    \begin{enumerate}
        \item \textbf{Confondre "Menant" et "Menée" dans les formules.} \emph{Règle :} $R = \frac{Z_{\text{menée}}}{Z_{\text{menant}}} = \frac{N_{\text{menant}}}{N_{\text{menée}}}$. Le \textbf{menant} est celui qui \textbf{entraîne} (moteur), la \textbf{menée} celui qui \textbf{reçoit} (utilisateur). Inverser les indices donne l'inverse du rapport.
        \item \textbf{Oublier le rendement $\eta$ pour le calcul du couple réel.} \emph{Piège :} Écrire $C_2 = R \cdot C_1$ au lieu de $C_2 = \eta \cdot R \cdot C_1$. Au BEPC, si $\eta$ n'est pas donné, on suppose $\eta=1$ (idéal), mais si $\eta$ est donné (ex: 0,85), son omission fait perdre la moitié des points du calcul.
        \item \textbf{Erreur de sens de rotation avec roue folle.} \emph{Classique :} Dire "la roue folle ne change rien". \textbf{Faux pour le sens !} Elle inverse le sens (engrenage extérieur). Elle ne change pas le \textbf{rapport} (valeur numérique). Distinguer \cle{Rapport} (valeur) et \cle{Sens} (cinématique).
        \item \textbf{Unités incohérentes dans la formule de Lewis ou Puissance.} \emph{Exemple :} $C$ en N$\cdot$m, $D$ en mm $\rightarrow$ $F_t$ faux. \textbf{Imposer :} $D$ en \textbf{mètres} pour $F_t$ en Newtons. Ou convertir $C$ en N$\cdot$mm. De même $\omega$ en rad/s (pas tr/min) pour $P = C\omega$ en Watts.
        \item \textbf{Confondre "Module" et "Pas".} \emph{Erreur :} $p = m$ (faux). \textbf{Relation :} $p = \pi \cdot m$. Le module $m$ est en mm, le pas $p$ aussi, mais $p > m$. Utiliser $m$ pour les calculs de diamètre ($D=mZ$), $p$ pour la conception de la denture (épaisseur dent = $p/2$).
    \end{enumerate}
    \textbf{Conseil final :} Relis l'énoncé pour identifier : 1. Qui entraîne qui ? (Menant/Menée). 2. Le rendement est-il donné ? 3. Les axes sont-ils parallèles, concourants, croisés ? 4. Les unités sont-elles homogènes (SI) avant de calculer ?
\end{attention}

\newpage

\coursec{Exercices}

\courssub{Je vérifie mes connaissances}

\begin{exercice}[QCM : Vocabulaire et sens de rotation]
\textbf{Complète chaque phrase en choisissant la bonne réponse.}
\begin{enumerate}
    \item La roue qui reçoit le mouvement du moteur s'appelle la roue \ldots
    \begin{itemize}
        \item A. menée
        \item B. menante
        \item C. motrice
        \item D. réceptrice
    \end{itemize}
    \item Dans un engrenage à deux roues dentées en prise directe, les deux roues tournent \ldots
    \begin{itemize}
        \item A. dans le même sens
        \item B. en sens inverse
        \item C. à la même vitesse angulaire
        \item D. sans frottement
    \end{itemize}
    \item Le rapport de transmission $r$ d'un engrenage est défini par \ldots
    \begin{itemize}
        \item A. $r = \frac{Z_{\text{menante}}}{Z_{\text{menée}}}$
        \item B. $r = \frac{N_{\text{menée}}}{N_{\text{menante}}}$
        \item C. $r = \frac{Z_{\text{menée}}}{Z_{\text{menante}}}$
        \item D. $r = N_{\text{menante}} \times Z_{\text{menée}}$
    \end{itemize}
    \item Si $r > 1$, l'engrenage est un \ldots
    \begin{itemize}
        \item A. multiplicateur de vitesse
        \item B. réducteur de vitesse
        \item C. inverseur de sens seulement
        \item D. système non fiable
    \end{itemize}
\end{enumerate}
\end{exercice}

\begin{exercice}[Vrai ou Faux ? Justifie]
\textbf{Pour chaque affirmation, coche Vrai ou Faux et justifie ta réponse par une phrase ou une formule.}
\begin{enumerate}
    \item Dans un train d'engrenages à trois roues, la roue intermédiaire change le sens de rotation de la roue menée par rapport à la roue menante.
    \item Le module d'un engrenage est le rapport du diamètre primitif sur le nombre de dents ; il doit être identique pour deux roues en prise.
    \item Un engrenage réducteur ($r > 1$) augmente le couple transmis à la roue menée.
    \item Les engrenages à denture hélicoïdale sont plus bruyants que les engrenages à denture droite à haute vitesse.
\end{enumerate}
\end{exercice}

\begin{exercice}[Définitions et notations]
\textbf{Donne la définition précise des termes suivants et indique l'unité le cas échéant :}
\begin{enumerate}
    \item Roue menante / Roue menée.
    \item Nombre de dents (noté $Z$).
    \item Vitesse de rotation (notée $N$).
    \item Rapport de transmission (noté $r$).
    \item Module (noté $m$).
    \item Diamètre primitif (noté $d$).
\end{enumerate}
\end{exercice}

\begin{exercice}[Calculs directs de rapport et de vitesse]
\textbf{Calcule la grandeur demandée. Donne la formule, le remplacement numérique et le résultat avec son unité.}
\begin{enumerate}
    \item Une roue menante de $Z_1 = 20$ dents entraîne une roue menée de $Z_2 = 60$ dents. Calculer le rapport de transmission $r$. Est-ce un réducteur ou un multiplicateur ?
    \item Un engrenage a un rapport de transmission $r = 4$. La roue menante tourne à $N_1 = 1\,200\ \text{tr/min}$. Calculer la vitesse de rotation $N_2$ de la roue menée.
    \item Une roue menée tourne à $N_2 = 300\ \text{tr/min}$ entraînée par une roue menante à $N_1 = 1\,500\ \text{tr/min}$. Déterminer le rapport de transmission $r$. Si $Z_1 = 18$ dents, combien de dents a la roue menée ($Z_2$) ?
    \item Un pignon de $Z_1 = 12$ dents (menant) engrenage avec une roue de $Z_2 = 48$ dents (menée). La vitesse du moteur est de $N_1 = 3\,000\ \text{tr/min}$. Calculer $r$ et $N_2$.
\end{enumerate}
\end{exercice}

\courssub{Je m'entraîne}

\begin{exercice}[Tableau de correspondance : engrenages simples]
\textbf{Complète le tableau ci-dessous. Chaque ligne correspond à un engrenage indépendant.}
\begin{center}
\begin{tabular}{|c|c|c|c|c|c|}\hline
\textbf{Config.} & $Z_1$ (dents) & $Z_2$ (dents) & $N_1$ (tr/min) & $N_2$ (tr/min) & $r$ \\ \hline
1 & 25 & 75 & 900 & \ldots & \ldots \\ \hline
2 & 15 & \ldots & 1\,800 & 450 & \ldots \\ \hline
3 & \ldots & 80 & 2\,400 & 600 & \ldots \\ \hline
4 & 30 & 90 & \ldots & 200 & \ldots \\ \hline
\end{tabular}
\end{center}
\emph{Indique pour chaque ligne s'il s'agit d'un réducteur ou d'un multiplicateur.}
\end{exercice}

\begin{exercice}[Engrenages et vélo : le dérailleur]
\textbf{Le pédalier (plateau) d'un VTT est la roue menante. Le pignon arrière (roue libre) est la roue menée.}
\begin{enumerate}
    \item Un cycliste utilise un plateau de $Z_1 = 42$ dents et un pignon de $Z_2 = 14$ dents. Calculer le rapport de transmission $r$. Combien de tours la roue arrière effectue-t-elle pour un tour de pédalier ?
    \item En côte, il change de pignon pour $Z_2 = 28$ dents (plateau inchangé). Calculer le nouveau rapport $r'$ et le nombre de tours de roue arrière par tour de pédalier.
    \item \textbf{Compare} les deux situations : dans quel cas l'effort sur les pédales est-il le plus faible ? Dans quel cas la vitesse du vélo est-elle la plus grande (pour un même rythme de pédalage) ? \textbf{Explique} en utilisant les notions de couple et de vitesse.
\end{enumerate}
\end{exercice}

\begin{exercice}[Train d'engrenages à trois roues]
\textbf{On considère un train d'engrenages composé de trois roues A, B, C en prise successives.}
\begin{itemize}
    \item Roue A (menante) : $Z_A = 20$ dents.
    \item Roue B (intermédiaire) : $Z_B = 40$ dents.
    \item Roue C (menée) : $Z_C = 60$ dents.
\end{itemize}
La roue A tourne à $N_A = 600\ \text{tr/min}$ dans le sens horaire.
\begin{enumerate}
    \item Dessine un schéma simplifié (cercles en prise) en indiquant le sens de rotation de chaque roue (flèches).
    \item Calculer la vitesse de rotation $N_B$ de la roue B.
    \item Calculer la vitesse de rotation $N_C$ de la roue C.
    \item Calculer le rapport de transmission global $r_{\text{global}} = \frac{N_A}{N_C}$. Vérifier que $r_{\text{global}} = \frac{Z_C}{Z_A}$ (la roue intermédiaire ne change pas le rapport global).
    \item Quel est le sens de rotation de la roue C par rapport à la roue A ?
\end{enumerate}
\end{exercice}

\begin{exercice}[Avantages, inconvénients et choix technologique]
\begin{enumerate}
    \item Cite \textbf{trois avantages} majeurs des engrenages par rapport aux transmissions par courroie ou par chaîne.
    \item Cite \textbf{trois inconvénients} des engrenages.
    \item \textbf{Justifie} pourquoi la boîte de vitesses d'une moto-taxi (type Yamaha 125 ou Bajaj Boxer) utilise des engrenages et non des courroies. Utilise les mots : \emph{couple}, \emph{précision}, \emph{lubrification}, \emph{compact}.
\end{enumerate}
\end{exercice}

\courssub{Je me prépare au Bac}

\begin{exercice}[Sujet type BEPC : Boîte de vitesses de moto-taxi]
\textbf{Contexte :} Au Cameroun, les moto-taxis (\emph{"benskin"}) utilisent des moteurs 4 temps monocylindres. La boîte de vitesses contient plusieurs paires d'engrenages en prise constante. On étudie le passage de la \textbf{1ère vitesse} (démarrage en côte) et de la \textbf{4ème vitesse} (vitesse de pointe sur route plate).

\textbf{Données :}
\begin{itemize}
    \item Arbre primaire (menant, sorti de l'embrayage) : pignon de $Z_{P1} = 14$ dents (1ère) et $Z_{P4} = 28$ dents (4ème).
    \item Arbre secondaire (mené, vers la chaîne finale) : couronne de $Z_{C1} = 42$ dents (1ère) et $Z_{C4} = 24$ dents (4ème).
    \item Régime moteur au démarrage : $N_{\text{moteur}} = 3\,000\ \text{tr/min}$.
    \item Rapport de la transmission finale (chaîne) : $r_{\text{finale}} = 3,0$ (pignon sortie boîte / couronne roue arrière).
\end{itemize}

\textbf{Questions :}
\begin{enumerate}
    \item Calculer le rapport de transmission de la boîte $r_{\text{boîte}}$ en 1ère vitesse et en 4ème vitesse.
    \item Calculer la vitesse de rotation de l'arbre secondaire $N_{\text{secondaire}}$ en 1ère et en 4ème vitesse.
    \item Calculer la vitesse de rotation de la roue arrière $N_{\text{roue}}$ dans les deux cas.
    \item \textbf{Analyse :} Comparer le couple disponible à la roue arrière en 1ère et en 4ème vitesse (rappelle la relation entre couple et rapport de transmission pour une puissance constante). Pourquoi la 1ère vitesse permet-elle de démarrer en côte avec un passager ?
    \item Sur le schéma ci-dessous (à reproduire), représente les deux paires d'engrenages (1ère et 4ème) en prise sur les arbres. Indique les sens de rotation des arbres primaire et secondaire (flèches).
\end{enumerate}
\end{exercice}

\begin{exercice}[Sujet type BEPC : Pompe à eau et groupe électrogène]
\textbf{Contexte :} Dans un quartier de Yaoundé, un groupe électrogène (moteur thermique) entraîne une pompe à eau centrifuge via un engrenage réducteur. Le moteur tourne à vitesse constante $N_M = 1\,500\ \text{tr/min}$. La pompe doit tourner à $N_P = 375\ \text{tr/min}$ pour un débit optimal.

\textbf{Données :}
\begin{itemize}
    \item Pignon sur le moteur (menant) : $Z_M = 20$ dents.
    \item Module de l'engrenage : $m = 4\ \text{mm}$.
    \item Largeur des dents : $b = 30\ \text{mm}$.
    \item Puissance transmise : $P = 5\ \text{kW}$.
\end{itemize}

\textbf{Questions :}
\begin{enumerate}
    \item Calculer le rapport de transmission $r$ nécessaire.
    \item Déterminer le nombre de dents $Z_P$ de la grande roue (menée) sur la pompe.
    \item Calculer les diamètres primitifs $d_M$ et $d_P$ des deux roues (rappel : $d = m \times Z$).
    \item Calculer l'entraxe $E$ (distance entre les axes des deux arbres).
    \item La puissance absorbée par la pompe est $P_{\text{utile}} = 4,5\ \text{kW}$. Calculer le rendement $\eta$ de cet engrenage. Où vont les pertes d'énergie ? Cite deux causes.
    \item \textbf{Sécurité :} Le carter de l'engrenage est ouvert pour inspection pendant que le groupe tourne. Cite \textbf{trois dangers} immédiats pour l'opérateur et la \textbf{consigne de sécurité} absolue à respecter avant toute intervention.
\end{enumerate}
\end{exercice}

\begin{exercice}[Sujet type BEPC : Système de direction à crémaillère (Technologie)]
\textbf{Contexte :} La direction d'une voiture (type Toyota Corolla très courant au Cameroun) transforme le mouvement de rotation du volant en mouvement de translation de la crémaillère (barre dentée) via un pignon.

\textbf{Données :}
\begin{itemize}
    \item Pignon de direction : $Z_{\text{pignon}} = 6$ dents (équivalent, car c'est un pignon sur crémaillère).
    \item Pas de la crémaillère : $p = 5\ \text{mm}$ (distance entre deux dents successives).
    \item On considère que le pignon fait un tour complet ($360^\circ$).
\end{itemize}

\textbf{Questions :}
\begin{enumerate}
    \item Quelle est la distance $d$ parcourue par la crémaillère (translation) pour un tour complet du pignon ? (Indice : $d = Z_{\text{pignon}} \times p$).
    \item Si le conducteur tourne le volant de $2$ tours complets de butée à butée (volant à gauche puis volant à droite), quelle est la course totale de la crémaillère ?
    \item Dans la réalité, le pignon a souvent plus de dents (ex: $Z = 12$ dents) et le pas est plus petit. \textbf{Explique} l'influence du nombre de dents du pignon et du pas sur :
    \begin{itemize}
        \item La démultiplication (effort au volant vs braquage des roues).
        \item La précision de la direction.
    \end{itemize}
    \item \textbf{Entretien :} Le soufflet de protection de la crémaillère est percé (poussière, eau, graisse qui sort). Quels sont les \textbf{deux risques majeurs} pour le mécanisme d'engrenage pignon/crémaillère ? Quelle est l'opération de maintenance préventive à faire tous les $50\,000\ \text{km}$ environ ?
\end{enumerate}
\end{exercice}

\newpage

\coursec{Corrigés des exercices}

\courssub{Je vérifie mes connaissances}

\begin{corrige}[Exercice 1 — QCM : Vocabulaire et sens de rotation]
\begin{enumerate}
    \item \textbf{Réponse B (menante).} La roue qui reçoit le mouvement du moteur est la roue \cle{menante} (on dit aussi roue \emph{motrice} ; le terme du cours est « menante »). La roue menée (réceptrice) est celle qui est entraînée.
    \item \textbf{Réponse B (en sens inverse).} Deux roues dentées en prise directe tournent toujours en sens inverse l'une de l'autre. Leurs vitesses de rotation ne sont égales que si $Z_1 = Z_2$, et il existe toujours des frottements (d'où la nécessité de la lubrification).
    \item \textbf{Réponse C.} Le rapport de transmission est :
    \[ r = \frac{Z_{\text{menée}}}{Z_{\text{menante}}} = \frac{N_{\text{menante}}}{N_{\text{menée}}} \]
    Attention au piège : $r$ est le rapport des nombres de dents \emph{menée sur menante}, mais le rapport des vitesses \emph{menante sur menée} (les vitesses varient en sens inverse des nombres de dents).
    \item \textbf{Réponse B (réducteur de vitesse).} Si $r > 1$, alors $N_{\text{menée}} = \dfrac{N_{\text{menante}}}{r} < N_{\text{menante}}$ : la roue menée tourne moins vite, c'est un \cle{réducteur} (mais le couple augmente). Si $r < 1$, c'est un multiplicateur.
\end{enumerate}
\end{corrige}

\begin{corrige}[Exercice 2 — Vrai ou Faux ? Justifie]
\begin{enumerate}
    \item \textbf{Faux.} Avec trois roues A--B--C en prise : si A tourne dans le sens horaire, B tourne en sens antihoraire et C tourne à nouveau dans le sens horaire. La roue intermédiaire \emph{rétablit} le sens : la roue menée C tourne dans le \textbf{même sens} que la roue menante A. (C'est avec \emph{deux} roues seulement que les sens sont inverses.)
    \item \textbf{Vrai.} Par définition, le module est $m = \dfrac{d}{Z}$ (diamètre primitif divisé par le nombre de dents). Pour que deux roues engrènent correctement, leurs dents doivent avoir la même taille, donc le \textbf{même module} : c'est une condition d'engrènement indispensable.
    \item \textbf{Vrai.} À puissance transmise constante ($P = C \times \omega$), si la vitesse diminue d'un facteur $r$, le couple augmente du même facteur : $C_{\text{menée}} \approx r \times C_{\text{menante}}$ (au rendement près). C'est pour cela qu'on démarre en 1ère vitesse.
    \item \textbf{Faux.} C'est l'inverse : la denture \cle{hélicoïdale} assure un contact progressif entre les dents, donc un fonctionnement \textbf{plus silencieux} et plus doux que la denture droite, surtout à haute vitesse. Son inconvénient est de créer des efforts axiaux (poussée le long de l'axe).
\end{enumerate}
\end{corrige}

\begin{corrige}[Exercice 3 — Définitions et notations]
\begin{enumerate}
    \item \textbf{Roue menante} : roue qui reçoit le mouvement du moteur et le transmet ; c'est la roue d'entrée. \textbf{Roue menée} : roue entraînée par la roue menante ; c'est la roue de sortie.
    \item \textbf{Nombre de dents $Z$} : nombre total de dents taillées sur la périphérie de la roue dentée. Grandeur sans unité (un nombre entier).
    \item \textbf{Vitesse de rotation $N$} : nombre de tours effectués par la roue par unité de temps. Unité usuelle : le tour par minute (tr/min).
    \item \textbf{Rapport de transmission $r$} : quotient du nombre de dents de la roue menée par celui de la roue menante :
    \[ r = \frac{Z_{\text{menée}}}{Z_{\text{menante}}} = \frac{N_{\text{menante}}}{N_{\text{menée}}} \quad \text{(sans unité)} \]
    \item \textbf{Module $m$} : grandeur caractérisant la taille des dents, définie par $m = \dfrac{d}{Z}$. Unité : le millimètre (mm). Deux roues en prise doivent avoir le même module.
    \item \textbf{Diamètre primitif $d$} : diamètre du cercle imaginaire (cercle primitif) sur lequel se fait l'engrènement sans glissement des deux roues. Unité : le millimètre (mm). Relation : $d = m \times Z$.
\end{enumerate}
\end{corrige}

\begin{corrige}[Exercice 4 — Calculs directs de rapport et de vitesse]
\begin{enumerate}
    \item Formule : $r = \dfrac{Z_2}{Z_1}$. Application : $r = \dfrac{60}{20} = \textbf{3}$. Comme $r > 1$, c'est un \textbf{réducteur} de vitesse (la roue menée tourne 3 fois moins vite).
    \item Formule : $N_2 = \dfrac{N_1}{r}$. Application : $N_2 = \dfrac{1\,200}{4} = \textbf{300 tr/min}$.
    \item Formule : $r = \dfrac{N_1}{N_2} = \dfrac{1\,500}{300} = \textbf{5}$. Puis $Z_2 = r \times Z_1 = 5 \times 18 = \textbf{90 dents}$.
    \item Rapport : $r = \dfrac{Z_2}{Z_1} = \dfrac{48}{12} = \textbf{4}$. Vitesse menée : $N_2 = \dfrac{N_1}{r} = \dfrac{3\,000}{4} = \textbf{750 tr/min}$.
\end{enumerate}
\textbf{Point de vigilance :} toujours écrire la formule littérale avant le remplacement numérique, et ne pas oublier l'unité (tr/min) dans le résultat.
\end{corrige}

\courssub{Je m'entraîne}

\begin{corrige}[Exercice 5 — Tableau de correspondance : engrenages simples]
On utilise $r = \dfrac{Z_2}{Z_1} = \dfrac{N_1}{N_2}$, donc $N_2 = \dfrac{N_1}{r}$ et $Z_2 = r \times Z_1$.
\begin{center}
\begin{tabular}{|c|c|c|c|c|c|c|}\hline
\textbf{Config.} & $Z_1$ & $Z_2$ & $N_1$ (tr/min) & $N_2$ (tr/min) & $r$ & Type \\ \hline
1 & 25 & 75 & 900 & \textbf{300} & \textbf{3} & Réducteur \\ \hline
2 & 15 & \textbf{60} & 1\,800 & 450 & \textbf{4} & Réducteur \\ \hline
3 & \textbf{20} & 80 & 2\,400 & 600 & \textbf{4} & Réducteur \\ \hline
4 & 30 & 90 & \textbf{600} & 200 & \textbf{3} & Réducteur \\ \hline
\end{tabular}
\end{center}
Détails :
\begin{enumerate}
    \item $r = \dfrac{75}{25} = 3$ ; $N_2 = \dfrac{900}{3} = 300$ tr/min.
    \item $r = \dfrac{1\,800}{450} = 4$ ; $Z_2 = 4 \times 15 = 60$ dents.
    \item $r = \dfrac{2\,400}{600} = 4$ ; $Z_1 = \dfrac{Z_2}{r} = \dfrac{80}{4} = 20$ dents.
    \item $r = \dfrac{90}{30} = 3$ ; $N_1 = r \times N_2 = 3 \times 200 = 600$ tr/min.
\end{enumerate}
Dans les quatre cas $r > 1$ : ce sont tous des \textbf{réducteurs}.
\end{corrige}

\begin{corrige}[Exercice 6 — Engrenages et vélo : le dérailleur]
\begin{enumerate}
    \item $r = \dfrac{Z_2}{Z_1} = \dfrac{14}{42} = \dfrac{1}{3} \approx 0,33$. Comme $r < 1$, c'est un \textbf{multiplicateur} : pour 1 tour de pédalier, la roue arrière fait $\dfrac{1}{r} = \textbf{3 tours}$.
    \item $r' = \dfrac{28}{42} = \dfrac{2}{3} \approx 0,67$. Pour 1 tour de pédalier, la roue arrière fait $\dfrac{1}{r'} = \textbf{1,5 tour}$.
    \item \textbf{Comparaison :}
    \begin{itemize}
        \item \textbf{Effort le plus faible} : situation 2 (pignon de 28 dents). Le rapport $r'$ est plus grand, donc le \cle{couple} transmis à la roue est plus grand pour le même effort musculaire : on « mouline » facilement, idéal en côte.
        \item \textbf{Vitesse la plus grande} : situation 1 (pignon de 14 dents). À rythme de pédalage égal, la roue fait 3 tours par tour de pédalier contre 1,5 : le vélo va deux fois plus vite, mais l'effort sur les pédales est plus grand.
    \end{itemize}
    \textbf{Bilan :} grand rapport $\Rightarrow$ grand couple, faible vitesse (côte) ; petit rapport $\Rightarrow$ faible couple, grande vitesse (plat ou descente). C'est exactement le principe de la boîte de vitesses.
\end{enumerate}
\end{corrige}

\begin{corrige}[Exercice 7 — Train d'engrenages à trois roues]
\begin{enumerate}
    \item Schéma simplifié :
    \begin{popfigure}
    \begin{center}
    \begin{tikzpicture}[scale=0.9]
        \draw[thick,popblue,fill=popblueL] (0,0) circle (1);
        \draw[thick,poporange,fill=poporangeL] (2.5,0) circle (1.5);
        \draw[thick,popgreen,fill=popgreenL] (6,0) circle (2);
        \fill (0,0) circle (2pt); \fill (2.5,0) circle (2pt); \fill (6,0) circle (2pt);
        \node at (0,0) [above=14pt] {\textbf{A} ($Z_A=20$)};
        \node at (2.5,0) [above=20pt] {\textbf{B} ($Z_B=40$)};
        \node at (6,0) [above=26pt] {\textbf{C} ($Z_C=60$)};
        \draw[-{Stealth},thick,popdark] (0.7,0.7) arc (45:135:1);
        \draw[-{Stealth},thick,popdark] (3.56,-1.06) arc (-45:-135:1.5);
        \draw[-{Stealth},thick,popdark] (4.59,1.41) arc (135:45:2);
    \end{tikzpicture}
    \end{center}
    \legende{Train de trois roues en prise : A (sens horaire), B (sens antihoraire), C (sens horaire). Les rayons sont proportionnels aux nombres de dents.}
    \end{popfigure}
    \item Contact A--B : $N_B = N_A \times \dfrac{Z_A}{Z_B} = 600 \times \dfrac{20}{40} = \textbf{300 tr/min}$ (sens antihoraire).
    \item Contact B--C : $N_C = N_B \times \dfrac{Z_B}{Z_C} = 300 \times \dfrac{40}{60} = \textbf{200 tr/min}$ (sens horaire).
    \item Rapport global : $r_{\text{global}} = \dfrac{N_A}{N_C} = \dfrac{600}{200} = \textbf{3}$. Vérification : $\dfrac{Z_C}{Z_A} = \dfrac{60}{20} = 3$. \textbf{Conforme} : la roue intermédiaire B ne modifie pas le rapport global (son nombre de dents se simplifie : $r = \dfrac{Z_B}{Z_A} \times \dfrac{Z_C}{Z_B} = \dfrac{Z_C}{Z_A}$).
    \item La roue C tourne dans le \textbf{même sens} que la roue A (sens horaire). La roue intermédiaire sert donc à \emph{conserver le sens de rotation} et à régler l'entraxe, pas à changer le rapport.
\end{enumerate}
\end{corrige}

\begin{corrige}[Exercice 8 — Avantages, inconvénients et choix technologique]
\begin{enumerate}
    \item \textbf{Trois avantages :}
    \begin{itemize}
        \item Transmission \textbf{sans glissement} : le rapport de vitesse est exact et constant (pas de patinage comme avec une courroie).
        \item Capacité à transmettre des \textbf{couples élevés} et des puissances importantes.
        \item Grande \textbf{durabilité} et bon rendement (95 à 98 \%), encombrement réduit (système compact).
    \end{itemize}
    \item \textbf{Trois inconvénients :}
    \begin{itemize}
        \item \textbf{Bruit} de fonctionnement (surtout en denture droite à haute vitesse).
        \item Nécessité d'une \textbf{lubrification} permanente et d'un carter étanche (coût d'entretien).
        \item \textbf{Coût de fabrication} élevé (taillage précis des dents) et impossibilité d'écarter les axes (les roues doivent se toucher).
    \end{itemize}
    \item \textbf{Justification (moto-taxi) :} la boîte de vitesses d'une moto doit transmettre un \emph{couple} important dans un espace très réduit : les engrenages sont \emph{compacts} et offrent une \emph{précision} parfaite du rapport (aucun glissement, passage de vitesse net). Logés dans un carter, ils bénéficient d'une \emph{lubrification} permanente à l'huile, ce qui garantit leur durée de vie malgré les efforts. Une courroie patinerait sous le couple et occuperait trop de place.
\end{enumerate}
\end{corrige}

\courssub{Je me prépare au Bac}

\begin{corrige}[Exercice 9 — Sujet type BEPC : Boîte de vitesses de moto-taxi]
\textbf{Barème indicatif (10 pts) :} Q1 : 2 pts ; Q2 : 2 pts ; Q3 : 2 pts ; Q4 : 2,5 pts ; Q5 : 1,5 pt.

\begin{enumerate}
    \item Rapport de boîte : $r_{\text{boîte}} = \dfrac{Z_{\text{couronne}}}{Z_{\text{pignon}}}$.
    \begin{itemize}
        \item 1ère : $r_1 = \dfrac{42}{14} = \textbf{3}$ (réducteur).
        \item 4ème : $r_4 = \dfrac{24}{28} \approx \textbf{0,86}$ (multiplicateur).
    \end{itemize}
    \item Vitesse de l'arbre secondaire : $N_{\text{secondaire}} = \dfrac{N_{\text{moteur}}}{r_{\text{boîte}}}$.
    \begin{itemize}
        \item 1ère : $N_{S1} = \dfrac{3\,000}{3} = \textbf{1\,000 tr/min}$.
        \item 4ème : $N_{S4} = \dfrac{3\,000}{0,857} \approx \textbf{3\,500 tr/min}$.
    \end{itemize}
    \item Vitesse de la roue arrière : $N_{\text{roue}} = \dfrac{N_{\text{secondaire}}}{r_{\text{finale}}}$ avec $r_{\text{finale}} = 3$.
    \begin{itemize}
        \item 1ère : $N_{\text{roue}} = \dfrac{1\,000}{3} \approx \textbf{333 tr/min}$.
        \item 4ème : $N_{\text{roue}} = \dfrac{3\,500}{3} \approx \textbf{1\,167 tr/min}$.
    \end{itemize}
    \item \textbf{Analyse :} à puissance moteur constante, le couple à la roue est proportionnel au rapport de transmission total : $C_{\text{roue}} \approx r_{\text{total}} \times C_{\text{moteur}}$. En 1ère, $r_{\text{total}} = 3 \times 3 = 9$ : le couple est multiplié par 9, la roue tourne lentement mais avec une grande force — c'est ce qui permet de \textbf{démarrer en côte avec un passager}. En 4ème, $r_{\text{total}} \approx 0,86 \times 3 \approx 2,6$ : faible couple mais grande vitesse, adaptée à la route plate.
    \item Schéma : deux paires de roues en prise entre l'arbre primaire (haut) et l'arbre secondaire (bas) ; pour la 1ère, petit pignon (14 dents) sur grande couronne (42 dents) ; pour la 4ème, grand pignon (28 dents) sur petite couronne (24 dents). Les deux arbres tournent en \textbf{sens inverse} (flèches opposées) car chaque paire est en prise directe.
\end{enumerate}
\textbf{Points de vigilance :} citer la condition « puissance constante » pour la relation couple/vitesse ; ne pas inverser pignon et couronne dans le rapport ; arrondir proprement ($0,857$) et garder les unités.
\end{corrige}

\begin{corrige}[Exercice 10 — Sujet type BEPC : Pompe à eau et groupe électrogène]
\textbf{Barème indicatif (10 pts) :} Q1 : 1 pt ; Q2 : 1,5 pt ; Q3 : 2 pts ; Q4 : 1,5 pt ; Q5 : 2 pts ; Q6 : 2 pts.

\begin{enumerate}
    \item Rapport nécessaire : $r = \dfrac{N_M}{N_P} = \dfrac{1\,500}{375} = \textbf{4}$.
    \item Nombre de dents de la roue menée : $Z_P = r \times Z_M = 4 \times 20 = \textbf{80 dents}$.
    \item Diamètres primitifs ($d = m \times Z$, avec $m = 4$ mm) :
    \begin{itemize}
        \item $d_M = 4 \times 20 = \textbf{80 mm}$ ;
        \item $d_P = 4 \times 80 = \textbf{320 mm}$.
    \end{itemize}
    \item Entraxe : $E = \dfrac{d_M + d_P}{2} = \dfrac{80 + 320}{2} = \textbf{200 mm}$.
    \item Rendement : $\eta = \dfrac{P_{\text{utile}}}{P} = \dfrac{4,5}{5} = 0,9 = \textbf{90 \%}$. Les 0,5 kW perdus (10 \%) se transforment en \textbf{chaleur} à cause : (1) des \emph{frottements} entre les dents en contact ; (2) du \emph{barattage de l'huile} et des frottements dans les roulements (on accepte aussi : chocs entre dents, usure).
    \item \textbf{Sécurité — trois dangers :}
    \begin{itemize}
        \item \emph{Happement} des doigts, des vêtements ou des cheveux par les dents en rotation (écrasement, amputation) ;
        \item Projection d'\emph{huile chaude} ou de particules dans les yeux ;
        \item Contact avec des pièces en rotation ou brûlantes (brûlures, coupures).
    \end{itemize}
    \textbf{Consigne absolue :} \textbf{arrêter le groupe électrogène et attendre l'immobilisation complète} (et le refroidissement) avant toute ouverture du carter ou intervention ; ne jamais travailler sur une machine en mouvement.
\end{enumerate}
\textbf{Points de vigilance :} le rendement s'exprime en pourcentage ; l'entraxe est la \emph{demi-somme} des diamètres primitifs (pas la somme) ; pour la sécurité, une consigne précise est attendue, pas une phrase vague.
\end{corrige}

\begin{corrige}[Exercice 11 — Sujet type BEPC : Système de direction à crémaillère]
\textbf{Barème indicatif (10 pts) :} Q1 : 2 pts ; Q2 : 2 pts ; Q3 : 3 pts ; Q4 : 3 pts.

\begin{enumerate}
    \item Pour un tour complet, le pignon fait passer $Z_{\text{pignon}}$ dents, donc la crémaillère avance de :
    \[ d = Z_{\text{pignon}} \times p = 6 \times 5 = \textbf{30 mm}. \]
    \item Pour 2 tours complets : course totale $= 2 \times 30 = \textbf{60 mm}$ (soit 6 cm de déplacement de la barre de direction).
    \item \textbf{Influence des paramètres :}
    \begin{itemize}
        \item \textbf{Démultiplication (effort) :} plus le pignon a de dents (ou plus le pas est grand), plus la crémaillère avance pour un tour de volant : la direction est plus « directe » mais demande \emph{plus d'effort}. À l'inverse, un petit pignon et un petit pas donnent une grande démultiplication : il faut tourner davantage le volant, mais l'\emph{effort est faible} (le couple fourni par le conducteur est amplifié).
        \item \textbf{Précision :} un petit pas et un nombre de dents élevé donnent un déplacement fin et progressif de la crémaillère : le braquage des roues est plus \emph{précis} et plus doux, sans à-coups.
    \end{itemize}
    \item \textbf{Entretien — deux risques majeurs si le soufflet est percé :}
    \begin{itemize}
        \item Pénétration de \emph{poussière et d'eau} : abrasion et corrosion des dents du pignon et de la crémaillère $\Rightarrow$ usure rapide, jeu dans la direction, voire grippage (perte de contrôle du véhicule) ;
        \item \emph{Fuite de la graisse} de lubrification $\Rightarrow$ frottements à sec, échauffement et détérioration du mécanisme.
    \end{itemize}
    \textbf{Maintenance préventive (tous les 50\,000 km environ) :} contrôler l'état des soufflets de crémaillère, les remplacer s'ils sont fendus, et \emph{regraisser} le mécanisme pignon/crémaillère (contrôle du jeu de direction).
\end{enumerate}
\textbf{Points de vigilance :} bien distinguer translation (crémaillère, en mm) et rotation (pignon, en tours) ; relier explicitement démultiplication et effort ; en technologie, toujours conclure par l'aspect \emph{sécurité} (une direction défaillante met des vies en danger).
\end{corrige}

\newpage

\coursec{Fiche bilan}

\courssub{Fiche bilan : Les engrenages}

\begin{popfigure}
\begin{center}\tcbox[colback=popblue,colframe=popblue,coltext=white,arc=9pt,boxsep=4pt,left=6pt,right=6pt]{\bfseries\small Les engrenages}\end{center}
\vspace{-2pt}
\begin{tcbraster}[raster columns=3,raster equal height=rows,raster column skip=6pt,raster row skip=6pt,enhanced,blankest]
\begin{tcolorbox}[enhanced,colback=poporange,colframe=poporange,coltext=white,boxrule=0.9pt,arc=3pt,left=4pt,right=4pt,top=3pt,bottom=3pt,boxsep=1pt,halign=left]\bfseries\scriptsize Définitions \& Vocabulaire Roue dentée, module, primitif, pas\end{tcolorbox}
\begin{tcolorbox}[enhanced,colback=popgreen,colframe=popgreen,coltext=white,boxrule=0.9pt,arc=3pt,left=4pt,right=4pt,top=3pt,bottom=3pt,boxsep=1pt,halign=left]\bfseries\scriptsize Types \& Sens de rotation Droits, hélicoïdaux, coniques, vis sans fin Engrenage extérieur / intérieur\end{tcolorbox}
\begin{tcolorbox}[enhanced,colback=poppurple,colframe=poppurple,coltext=white,boxrule=0.9pt,arc=3pt,left=4pt,right=4pt,top=3pt,bottom=3pt,boxsep=1pt,halign=left]\bfseries\scriptsize Avantages \& Inconvénients + Précis, puissant, compact - Bruit, lubrification, coût\end{tcolorbox}
\begin{tcolorbox}[enhanced,colback=poppink,colframe=poppink,coltext=white,boxrule=0.9pt,arc=3pt,left=4pt,right=4pt,top=3pt,bottom=3pt,boxsep=1pt,halign=left]\bfseries\scriptsize Rapport de transmission $R = \frac{Z_2}{Z_1} = \frac{\omega_1}{\omega_2} = \frac{D_2}{D_1}$ Calcul vitesse / couple\end{tcolorbox}
\begin{tcolorbox}[enhanced,colback=popgold,colframe=popgold,coltext=white,boxrule=0.9pt,arc=3pt,left=4pt,right=4pt,top=3pt,bottom=3pt,boxsep=1pt,halign=left]\bfseries\scriptsize Sécurité \& Entretien Carter, graissage, alignement, arrêt d'urgence, EPI\end{tcolorbox}
\begin{tcolorbox}[enhanced,colback=popblue!70!popgreen,colframe=popblue!70!popgreen,coltext=white,boxrule=0.9pt,arc=3pt,left=4pt,right=4pt,top=3pt,bottom=3pt,boxsep=1pt,halign=left]\bfseries\scriptsize Applications au Cameroun Moto-taxi, groupe électrogène, moteur voiture, moulin à maïs\end{tcolorbox}
\end{tcbraster}

\legende{Carte mentale récapitulative de la leçon sur les engrenages.}
\end{popfigure}

\begin{bilan}

\textbf{1. Définitions clés}
\begin{itemize}
  \item \cle{Engrenage} : Ensemble de deux roues dentées (pignon et roue) qui s'engrènent pour transmettre un mouvement de rotation.
  \item \cle{Pignon} : La plus petite roue (moteur le plus souvent). \cle{Roue} : La plus grande roue (récepteur).
  \item \cle{Cercle primitif} : Cercle imaginaire où s'effectue le contact pur roulement sans glissement. Rayon $R$.
  \item \cle{Module} ($m$) : Taille des dents. $m = \frac{D}{Z}$ (en mm). Dents identiques $\Rightarrow$ même module.
  \item \cle{Pas} ($p$) : Distance entre deux dents successives sur le cercle primitif. $p = \pi \times m$.
  \item \cle{Déport} : Distance entre les axes des deux roues. $E = \frac{D_1 + D_2}{2} = \frac{m(Z_1+Z_2)}{2}$.
\end{itemize}

\textbf{2. Types d'engrenages et sens de rotation}
\begin{center}
\begin{tabularx}{\linewidth}{|l|X|X|}\hline
\rowcolor{popblueL}
\textbf{Type} & \textbf{Caractéristiques} & \textbf{Sens de rotation} \\ \hline
Engrenage droit (extérieur) & Axes parallèles, dents droites & \textbf{Inverses} (pignon $\rightarrow$ roue) \\ \hline
Engrenage intérieur & Axes parallèles, dents à l'intérieur & \textbf{Même sens} \\ \hline
Engrenage hélicoïdal & Axes parallèles, dents inclinées & Inverses (plus silencieux, charge axiale) \\ \hline
Engrenage conique & Axes concourants (souvent $90^\circ$) & Inverses (selon montage) \\ \hline
Vis sans fin / Roue & Axes croisés ($90^\circ$), grand rapport & Non réversible (vis $\rightarrow$ roue seulement) \\ \hline
\end{tabularx}
\end{center}

\textbf{3. Avantages et Inconvénients}
\begin{itemize}
  \item \textbf{Avantages :} Transmission synchrone (pas de glissement), rendement élevé ($> 95\%$), compacité, grande puissance transmissible, rapport de transmission précis.
  \item \textbf{Inconvénients :} Bruit et vibrations (surtout droits), nécessité de lubrification (graisse/huile), usure des dents, coût de fabrication élevé, axes fixes (distance imposée).
\end{itemize}

\textbf{4. Rapport de transmission (Formules à connaître par cœur)}
\begin{center}
\begin{tabularx}{\linewidth}{|c|X|c|}\hline
\rowcolor{poporangeL}
\cle{Grandeur} & \cle{Formule} & \cle{Unité} \\ \hline
Rapport de transmission & $R = \frac{Z_2}{Z_1} = \frac{D_2}{D_1} = \frac{\omega_1}{\omega_2} = \frac{N_1}{N_2}$ & Sans unité \\ \hline
Vitesse angulaire & $\omega_2 = \frac{\omega_1}{R}$ & rad/s ou tr/min \\ \hline
Couple (idéal) & $C_2 = R \times C_1$ & N$\cdot$m \\ \hline
Puissance (idéal) & $P_2 = P_1$ (conservation) & Watt \\ \hline
\end{tabularx}
\end{center}
\emph{Convention :} Indice 1 = Pignon (entrée), Indice 2 = Roue (sortie). $R > 1$ : Réducteur (vitesse $\downarrow$, couple $\uparrow$). $R < 1$ : Multiplicateur.

\textbf{5. Méthode express : Calculer les caractéristiques d'un engrenage}
\begin{enumerate}
  \item Identifier $Z_1$ (pignon) et $Z_2$ (roue) ou les diamètres $D_1, D_2$.
  \item Calculer le rapport $R = \frac{Z_2}{Z_1}$.
  \item Déduire la vitesse de sortie : $N_2 = \frac{N_1}{R}$ (ou $\omega_2 = \frac{\omega_1}{R}$).
  \item Déduire le couple de sortie : $C_2 = R \times C_1$ (rendement $\eta \approx 1$ si non précisé).
  \item Vérifier le déport $E = \frac{m(Z_1+Z_2)}{2}$ pour le montage.
\end{enumerate}

\textbf{6. Sécurité et Entretien}
\begin{itemize}
  \item \textbf{Carter de protection} obligatoire (évite pincement doigts, projection graisse).
  \item \textbf{Lubrification} régulière (graisse ou bain d'huile) selon préconisations constructeur.
  \item \textbf{Alignement} axes et parallélisme à contrôler (usure asymétrique sinon).
  \item \textbf{Arrêt d'urgence} accessible près de la machine.
  \item \textbf{EPI :} Lunettes, gants (hors zone engrenage), chaussures de sécurité, pas de vêtements amples.
\end{itemize}

\end{bilan}

\begin{aretenir}[Checklist avant le Bac]
\begin{itemize}
  \item $\square$ Je sais définir : engrenage, pignon, roue, cercle primitif, module, pas, déport.
  \item $\square$ Je sais nommer 4 types d'engrenages et donner le sens de rotation pour l'engrenage droit extérieur et intérieur.
  \item $\square$ Je connais la formule du rapport de transmission $R = \frac{Z_2}{Z_1} = \frac{\omega_1}{\omega_2}$.
  \item $\square$ Je sais calculer la vitesse de sortie $N_2$ connaissant $N_1$ et $R$.
  \item $\square$ Je sais calculer le couple de sortie $C_2$ connaissant $C_1$ et $R$.
  \item $\square$ Je sais expliquer pourquoi un engrenage est plus précis qu'une courroie (pas de glissement).
  \item $\square$ Je cite 2 avantages et 2 inconvénients des engrenages.
  \item $\square$ Je sais calculer le déport $E = \frac{m(Z_1+Z_2)}{2}$ pour un montage.
  \item $\square$ Je connais 3 règles de sécurité pour une machine à engrenages (carter, graissage, arrêt urgence, EPI).
  \item $\square$ Je donne 2 exemples d'applications au Cameroun (moto-taxi, groupe électrogène, moulin, voiture).
  \item $\square$ Je sais distinguer réducteur ($R>1$) et multiplicateur ($R<1$).
  \item $\square$ Je sais que la puissance est conservée (hors pertes) : $P_1 = P_2$.
\end{itemize}
\end{aretenir}

\findecours

\end{document}
